% Economic Life and Occupations — Anglais 1re Tech — Cours signé OnBuch+
% Programme officiel MINESEC. Compiler avec : tectonic N02-economic-life-and-occupations.tex
\newcommand{\DOCMATIERE}{Anglais}
\newcommand{\DOCNIVEAU}{1re Tech}
\newcommand{\DOCMODULE}{Anglais – classe de Première technique}
\newcommand{\DOCLECON}{N02}
\newcommand{\DOCTITRE}{Economic Life and Occupations}
% OnBuch+ — préambule « cours signés » ENRICHI (compilé avec tectonic / XeLaTeX)
% Système visuel « Pop » d'OnBuch+ : fond crème (jamais blanc), bordures encre
% épaisses, ombres « dures » décalées, titres Archivo Black en capitales.
% Version MATHÉMATIQUES : graphes (pgfplots/tikz), boîtes pédagogiques
% (définition, propriété, méthode, attention, le savais-tu, exercice résolu,
% exercice, TP, bilan), page de garde et signature OnBuch+ ancrée en bas.
%
% Macros attendues dans main.tex AVANT \input{preamble} :
%   \DOCMATIERE  \DOCNIVEAU  \DOCMODULE  \DOCLECON (numéro)  \DOCTITRE
\documentclass[11pt]{article}

\usepackage{fontspec}
\usepackage[french]{babel}
\usepackage[a4paper,margin=2cm,top=3.4cm,bottom=2.8cm,headheight=1.4cm,footskip=1.3cm]{geometry}
\usepackage{amsmath,amssymb,amsfonts}
\usepackage{mathtools} % \xmapsto, \coloneqq... souvent utilisés par les modèles
\usepackage{cancel} % simplifications barrées (bilans de forces)
% \abs{z} (module, valeur absolue) et \norm{u} (norme), utilisés par les modèles
\providecommand{\abs}{}\let\abs\relax\DeclarePairedDelimiter{\abs}{\lvert}{\rvert}
\providecommand{\norm}{}\let\norm\relax\DeclarePairedDelimiter{\norm}{\lVert}{\rVert}
\usepackage[table]{xcolor} % [table] : fournit aussi \rowcolors (alternance de lignes)
\usepackage{pagecolor}
\usepackage{tikz}
\usetikzlibrary{arrows.meta,positioning,calc,shapes.geometric,shapes.misc,shapes.arrows,decorations.pathmorphing,decorations.pathreplacing,decorations.markings,patterns,shadows,mindmap,trees,fit,backgrounds,matrix,intersections,angles,quotes}
\usepackage{pgfplots}
\usepackage[version=4]{mhchem} % équations nucléaires \ce{^{235}_{92}U}
\usepackage[european,siunitx]{circuitikz} % schémas électriques
\pgfplotsset{compat=1.18}
\usepgfplotslibrary{fillbetween,groupplots} % aires entre courbes (calcul intégral)
\usepackage{enumitem}
\usepackage{fancyhdr}
% [most] charge toutes les bibliothèques tcolorbox (listings, minted, raster,
% documentation...) et gonfle inutilement la mémoire TeX sur un document long
% (~300 boîtes) : on ne charge que ce qui est réellement utilisé.
\usepackage[skins,breakable]{tcolorbox}
\usepackage{graphicx}
\usepackage{colortbl}
\usepackage{booktabs}
\usepackage{tabularx}
\usepackage{multirow}
\usepackage{diagbox} % case d'en-tête diagonale des tableaux à double entrée
% Colonnes extensibles centrée / gauche / droite (C, L, R), souvent utilisées par les modèles
\newcolumntype{C}{>{\centering\arraybackslash}X}
\newcolumntype{L}{>{\raggedright\arraybackslash}X}
\newcolumntype{R}{>{\raggedleft\arraybackslash}X}
\usepackage{array}
\usepackage{multicol}
\usepackage{siunitx}
% parse-numbers=false : accepte aussi \SI{2,5 \times 10^{-3}}{...}
\sisetup{locale=FR,output-decimal-marker={,},group-minimum-digits=4,parse-numbers=false}
\DeclareSIUnit\ohm{\ensuremath{\symup{\Omega}}} % U+2126 absent de la police
\DeclareSIUnit\division{div} % réglages d'oscilloscope (ms/div, V/div)
\DeclareSIUnit\tr{tr} % tours (vitesse de rotation, ex. \SI{1500}{\tr\per\minute})
\DeclareSIUnit\molar{M} % concentration molaire (ex. \SI{3}{\molar}, \SI{1}{\micro\molar})
\DeclareSIUnit\an{an} % année (le modèle confond parfois avec \year, un compteur TeX) — ex. \SI{5}{\kilo\meter\per\an}
\DeclareSIUnit\FCFA{FCFA} % franc CFA (ex. \SI{500}{\FCFA\per\kilo\gram})
\DeclareSIUnit\degreeBrix{\textdegree Brix} % teneur en sucre d'un sirop/jus (réfractométrie)
\DeclareSIUnit\month{mois} % le modèle utilise parfois \month, un compteur TeX, comme unité de durée
\DeclareSIUnit\microliter{\micro\liter} % le modèle écrit parfois \microliter en un seul token
\DeclareSIUnit\million{millions} % dénombrement (ex. \SI{15}{\million\per\mL} spermatozoïdes)
\usepackage{qrcode}
% \degree : commande fréquemment utilisée par les modèles pour un angle ou une
% température en dehors de \SI{}{} (non fournie par les paquets chargés).
\providecommand{\degree}{^{\circ}}
% \qed (fin de démonstration), utilisé par les modèles sans amsthm
\providecommand{\qed}{\hfill$\square$}
% \barre : double barre des valeurs interdites dans un tableau de variations (tkz-tab)
\providecommand{\barre}{\ensuremath{\|}}
% environnement proof (amsthm non chargé) utilisé par les modèles
\ifdefined\proof\else\newenvironment{proof}[1][Démonstration]{\par\noindent\textit{#1.}\ }{\qed\par}\fi
\usepackage{lastpage}
\usepackage{hyperref}
\hypersetup{hidelinks}

% ============================================================
% Palette « Pop » OnBuch+ — mêmes hex que la classe P (plus_ui.dart)
% ============================================================
\definecolor{poporange}{HTML}{F59321}
\definecolor{popdark}{HTML}{E07A0C}
\definecolor{popcream}{HTML}{FFF6E8}
\definecolor{popink}{HTML}{1C1714}
\definecolor{poppurple}{HTML}{7A5AE0}
\definecolor{popgreen}{HTML}{2E9E6A}
\definecolor{poppink}{HTML}{E0457B}
\definecolor{popblue}{HTML}{2F7DE1}
\definecolor{popgold}{HTML}{C98A12}
\definecolor{popmuted}{HTML}{8A7F76}
\definecolor{popline}{HTML}{EADFD2}
\definecolor{popgray}{HTML}{8A7F76}
\colorlet{popgrey}{popgray}
% Alias tolérants (le modèle nomme parfois les couleurs différemment)
\colorlet{popwhite}{white}
\colorlet{popblack}{popink}
\colorlet{popred}{poppink}
\colorlet{popyellow}{popgold}
% Teintes claires (fonds de tableaux/encadrés) — le modèle nomme parfois
% les couleurs avec un suffixe L (ex. popblueL) sans les avoir définies.
\colorlet{poporangeL}{poporange!20}
\colorlet{popdarkL}{popdark!20}
\colorlet{poppurpleL}{poppurple!20}
\colorlet{popgreenL}{popgreen!20}
\colorlet{poppinkL}{poppink!20}
\colorlet{popblueL}{popblue!20}
\colorlet{popgoldL}{popgold!20}
\colorlet{popredL}{popred!20}
\colorlet{popgrayL}{popgray!20}
% Teintes claires pour fonds de tableaux / zones
\colorlet{popblueL}{popblue!10!white}
\colorlet{poporangeL}{poporange!14!white}
\colorlet{popgreenL}{popgreen!12!white}
\colorlet{poppurpleL}{poppurple!12!white}
\colorlet{poppinkL}{poppink!12!white}
\colorlet{popgoldL}{popgold!14!white}

\pagecolor{popcream}

% ============================================================
% Typographie
% ============================================================
\defaultfontfeatures{Ligatures=TeX}
\setmainfont{PlusJakartaSans-Medium.ttf}[Path=../../../fonts/,BoldFont=PlusJakartaSans-Bold.ttf]
% Maths en sans-serif (Fira Math) avec chiffres et lettres droites de Plus
% Jakarta, pour que les formules se fondent dans le texte.
\usepackage[math-style=ISO,bold-style=ISO]{unicode-math}
\setmathfont{FiraMath-Regular.otf}[Path=../../../fonts/]
\setmathfont{PlusJakartaSans-Medium.ttf}[Path=../../../fonts/,range={up/num,up/{Latin,latin},\mathup}]
\usepackage{icomma} % virgule décimale sans espace en mode math
% Caractères mathématiques Unicode absents des polices de texte (Plus Jakarta,
% Archivo Black) : rendus par leur équivalent mathématique, en texte comme en
% formule — sinon ils s'affichent en carré vide (ex. « Arithmétique dans ℤ »).
\usepackage{newunicodechar}
\newunicodechar{ℝ}{\ensuremath{\mathbb{R}}}
\newunicodechar{ℤ}{\ensuremath{\mathbb{Z}}}
\newunicodechar{ℂ}{\ensuremath{\mathbb{C}}}
\newunicodechar{ℕ}{\ensuremath{\mathbb{N}}}
\newunicodechar{ℚ}{\ensuremath{\mathbb{Q}}}
\newunicodechar{𝔻}{\ensuremath{\mathbb{D}}}
\newunicodechar{α}{\ensuremath{\alpha}}
\newunicodechar{β}{\ensuremath{\beta}}
\newunicodechar{γ}{\ensuremath{\gamma}}
\newunicodechar{δ}{\ensuremath{\delta}}
\newunicodechar{ε}{\ensuremath{\varepsilon}}
\newunicodechar{θ}{\ensuremath{\theta}}
\newunicodechar{λ}{\ensuremath{\lambda}}
\newunicodechar{μ}{\ensuremath{\mu}}
\newunicodechar{σ}{\ensuremath{\sigma}}
\newunicodechar{τ}{\ensuremath{\tau}}
\newunicodechar{φ}{\ensuremath{\varphi}}
\newunicodechar{ψ}{\ensuremath{\psi}}
\newunicodechar{ω}{\ensuremath{\omega}}
\newunicodechar{Σ}{\ensuremath{\Sigma}}
% majuscules produites par \MakeUppercase dans les titres (α-aminés -> Α)
\newunicodechar{Α}{\ensuremath{\alpha}}
\newunicodechar{Β}{\ensuremath{\beta}}
\newunicodechar{Γ}{\ensuremath{\Gamma}}
\newunicodechar{Δ}{\ensuremath{\Delta}}
\newunicodechar{Ε}{\ensuremath{\varepsilon}}
\newunicodechar{Θ}{\ensuremath{\Theta}}
\newunicodechar{Λ}{\ensuremath{\Lambda}}
\newunicodechar{Μ}{\ensuremath{\mu}}
\newunicodechar{Τ}{\ensuremath{\tau}}
\newunicodechar{Φ}{\ensuremath{\Phi}}
\newunicodechar{Ψ}{\ensuremath{\Psi}}
\newunicodechar{Ω}{\ensuremath{\Omega}}
\newunicodechar{↦}{\ensuremath{\mapsto}}
\newunicodechar{∈}{\ensuremath{\in}}
\newunicodechar{∉}{\ensuremath{\notin}}
\newunicodechar{∧}{\ensuremath{\wedge}}
\newunicodechar{⇔}{\ensuremath{\Leftrightarrow}}
\newunicodechar{⟺}{\ensuremath{\Longleftrightarrow}}
\newunicodechar{⇒}{\ensuremath{\Rightarrow}}
\newunicodechar{⟹}{\ensuremath{\Longrightarrow}}
\newunicodechar{∘}{\ensuremath{\circ}}
\newunicodechar{∪}{\ensuremath{\cup}}
\newunicodechar{∩}{\ensuremath{\cap}}
\newunicodechar{≡}{\ensuremath{\equiv}}
\newunicodechar{⊂}{\ensuremath{\subset}}
\newunicodechar{⊥}{\ensuremath{\perp}}
\newunicodechar{∃}{\ensuremath{\exists}}
\newunicodechar{∀}{\ensuremath{\forall}}
\newunicodechar{∛}{\ensuremath{\sqrt[3]{\,}}}
\newunicodechar{∜}{\ensuremath{\sqrt[4]{\,}}}
\newunicodechar{⃗}{\ensuremath{{}^{\to}}}
\newunicodechar{ⁿ}{\ifmmode^{n}\else\textsuperscript{\ensuremath{n}}\fi}
\newunicodechar{ˣ}{\ifmmode^{x}\else\textsuperscript{\ensuremath{x}}\fi}
\newunicodechar{ᵇ}{\ifmmode^{b}\else\textsuperscript{\ensuremath{b}}\fi}
\newunicodechar{ʳ}{\ifmmode^{r}\else\textsuperscript{\ensuremath{r}}\fi}
\newunicodechar{ᵘ}{\ifmmode^{u}\else\textsuperscript{\ensuremath{u}}\fi}
\newunicodechar{ʸ}{\ifmmode^{y}\else\textsuperscript{\ensuremath{y}}\fi}
\newunicodechar{ᵃ}{\ifmmode^{a}\else\textsuperscript{\ensuremath{a}}\fi}
\newunicodechar{ᵉ}{\ifmmode^{e}\else\textsuperscript{\ensuremath{e}}\fi}
\newunicodechar{ᵏ}{\ifmmode^{k}\else\textsuperscript{\ensuremath{k}}\fi}
\newunicodechar{ᵖ}{\ifmmode^{p}\else\textsuperscript{\ensuremath{p}}\fi}
\newunicodechar{ᴺ}{\ifmmode^{N}\else\textsuperscript{\ensuremath{N}}\fi}
\newunicodechar{ᶜ}{\ifmmode^{c}\else\textsuperscript{\ensuremath{c}}\fi}
\newunicodechar{ʰ}{\ifmmode^{h}\else\textsuperscript{\ensuremath{h}}\fi}
\newunicodechar{ᵠ}{\ifmmode^{q}\else\textsuperscript{\ensuremath{q}}\fi}
\newunicodechar{ₙ}{\ifmmode_{n}\else\textsubscript{\ensuremath{n}}\fi}
\newunicodechar{ₐ}{\ifmmode_{a}\else\textsubscript{\ensuremath{a}}\fi}
\newunicodechar{ᵢ}{\ifmmode_{i}\else\textsubscript{\ensuremath{i}}\fi}
\newunicodechar{ᵣ}{\ifmmode_{r}\else\textsubscript{\ensuremath{r}}\fi}
\newunicodechar{ₖ}{\ifmmode_{k}\else\textsubscript{\ensuremath{k}}\fi}
\newunicodechar{ᵦ}{\ifmmode_{\beta}\else\textsubscript{\ensuremath{\beta}}\fi}
\newunicodechar{ₓ}{\ifmmode_{x}\else\textsubscript{\ensuremath{x}}\fi}
\newfontfamily\popdisplay{ArchivoBlack-Regular.ttf}[Path=../../../fonts/]
\newfontfamily\poplogo{SpaceGrotesk-Bold.ttf}[Path=../../../fonts/]
\newfontfamily\popsemi{PlusJakartaSans-SemiBold.ttf}[Path=../../../fonts/]
\newfontfamily\popxbold{PlusJakartaSans-ExtraBold.ttf}[Path=../../../fonts/]
\newfontfamily\popxboldls{PlusJakartaSans-ExtraBold.ttf}[Path=../../../fonts/,LetterSpace=3.0]

\setlist[enumerate]{leftmargin=1.7em,itemsep=5pt,topsep=4pt}
\setlist[itemize]{leftmargin=1.7em,itemsep=4pt,topsep=4pt,label={\color{poporange}\textbullet}}
\setlength{\parindent}{0pt}
\setlength{\parskip}{5pt}
\renewcommand{\arraystretch}{1.25}

% pgfplots : style Pop par défaut
\pgfplotsset{
  every axis/.append style={axis line style={popink,line width=1pt},tick style={popink},
    grid style={popline,line width=0.5pt},label style={font=\small},tick label style={font=\footnotesize}},
  popcurve/.style={poporange,line width=1.8pt,no markers,smooth},
}
% Même style disponible dans un \draw TikZ simple (hors pgfplots)
\tikzset{popcurve/.style={poporange,line width=1.8pt}}
\tikzset{popgrid/.style={popline,very thin}}
\colorlet{popcurve}{poporange} % le nom du style est parfois utilisé comme couleur

% ============================================================
% Pill / squiggle
% ============================================================
\newcommand{\poppill}[3]{%
  \tikz[baseline=-0.55ex]\node[draw=popink,line width=1.2pt,rounded corners=2.6pt,%
    fill=#2,inner xsep=6.5pt,inner ysep=3pt]{%
    \popxboldls\fontsize{7}{7}\selectfont\color{#3}\MakeUppercase{#1}};%
}
\newcommand{\popsquiggle}[1]{%
  \tikz[baseline=-0.4ex]{
    \draw[#1,line width=2.6pt,line cap=round]
      plot [smooth] coordinates {(0,0.09) (0.34,0.01) (0.68,0.21) (1.02,0.01) (1.36,0.10)};
  }%
}
% Mot-clé surligné (marqueur orange sous le texte)
\newcommand{\cle}[1]{{\popxbold\color{popdark}#1}}

% ============================================================
% En-tête / pied
% ============================================================
\pagestyle{fancy}
\fancyhf{}
\renewcommand{\headrulewidth}{0pt}
\renewcommand{\footrulewidth}{0pt}
\lhead{\raisebox{-0.15ex}{\poplogo\fontsize{13}{13}\selectfont\color{popink}OnBuch\color{poporange}+}}
\rhead{\poppill{\DOCMATIERE\ \textbullet\ \DOCNIVEAU\ \textbullet\ Leçon \DOCLECON}{white}{popink}}
\lfoot{\popxbold\fontsize{7.6}{7.6}\selectfont\color{popmuted}\MakeUppercase{Cours signé OnBuch+}\ \textbullet\ {\popsemi\DOCTITRE}}
\rfoot{\tikz[baseline=-0.6ex]\node[rounded corners=8.5pt,fill=popink,minimum height=17pt,minimum width=17pt,inner xsep=5pt,inner sep=0]{\popxbold\fontsize{8}{8}\selectfont\color{popcream}\thepage\,/\,\pageref*{LastPage}};}

% ============================================================
% Titres
% ============================================================
\newcommand{\chaptitre}[2]{%
  \begin{tcolorbox}[enhanced,colback=poporange,colframe=popink,boxrule=2.6pt,arc=11pt,
      drop shadow=popink,left=15pt,right=15pt,top=12pt,bottom=11pt,before skip=3pt,after skip=18pt]
    \poppill{#2}{popink}{popcream}\par\vspace{9pt}
    {\raggedright\hyphenpenalty=10000\exhyphenpenalty=10000\popdisplay\fontsize{22}{24}\selectfont\color{popcream}\MakeUppercase{#1}\par}
  \end{tcolorbox}
}
% Section numérotée : pastille orange avec le numéro + titre Archivo
\newcounter{popsec}
\newcommand{\coursec}[1]{%
  \stepcounter{popsec}\setcounter{popsub}{0}%
  \par\vspace{16pt}\noindent
  \begin{minipage}{\linewidth}
  \tikz[baseline=-0.7ex]\node[draw=popink,line width=1.6pt,fill=poporange,rounded corners=4pt,minimum size=22pt,inner sep=0]{\popdisplay\fontsize{12}{12}\selectfont\color{popcream}\thepopsec};%
  \hspace{8pt}{\popdisplay\fontsize{15}{17}\selectfont\color{popink}\MakeUppercase{#1}}\par
  \vspace{4pt}\noindent{\color{poporange}\rule{36pt}{3.6pt}}
  \end{minipage}\par\nopagebreak\vspace{8pt}
}
\newcounter{popsub}
\newcommand{\courssub}[1]{%
  \stepcounter{popsub}%
  \par\vspace{9pt}\noindent{\popxbold\fontsize{11.5}{13}\selectfont\color{popink}{\color{poporange}\thepopsec.\thepopsub}\hspace{6pt}#1}\par\nopagebreak\vspace{4pt}
}

% ============================================================
% Boîtes pédagogiques — même recette : fond blanc, bordure épaisse colorée,
% ombre dure encre, badge plein. Titre optionnel [#1].
% ============================================================
\tcbset{popbase/.style={breakable,enhanced,colback=white,boxrule=2.3pt,arc=9pt,
  shadow={3pt}{-3pt}{0pt}{popink},left=14pt,right=14pt,top=11pt,bottom=11pt,before skip=11pt,after skip=15pt,
  fontupper=\popsemi\fontsize{10.6}{14.5}\selectfont\color{popink},
  fontlower=\popsemi\fontsize{10.6}{14.5}\selectfont\color{popink}}}
% Coche dessinée (le glyphe ✓ n'existe pas dans les polices du document)
\newcommand{\popcheck}[1]{\tikz[baseline=-0.5ex]\draw[#1,line width=1.6pt,line cap=round,line join=round](-0.13,0.0)--(-0.04,-0.09)--(0.13,0.1);}
\providecommand{\coche}{\popcheck{popgreen}}
\providecommand{\resultat}[1]{\par\smallskip\noindent\fcolorbox{popgreen}{popgreenL}{\parbox{\dimexpr\linewidth-2\fboxsep-2\fboxrule\relax}{\textbf{Résultat.} #1}}\par\smallskip}
\newcommand{\popboxhead}[3]{\poppill{#1}{#2}{white}\ifx\relax#3\relax\else\hspace{6pt}{\popxbold\fontsize{10.6}{12}\selectfont\color{popink}#3}\fi\par\vspace{7pt}}

\newtcolorbox{aretenir}[1][]{popbase,colframe=popgreen,before upper={\popboxhead{À retenir}{popgreen}{#1}}}
\newtcolorbox{exemplebox}[1][]{popbase,colframe=poporange,before upper={\popboxhead{Exemple}{poporange}{#1}}}
\newtcolorbox{definition}[1][]{popbase,colframe=popblue,before upper={\popboxhead{Définition}{popblue}{#1}}}
\newtcolorbox{propriete}[1][]{popbase,colframe=poppurple,before upper={\popboxhead{Propriété}{poppurple}{#1}}}
\newtcolorbox{methode}[1][]{popbase,colframe=popgold,before upper={\popboxhead{Méthode}{popgold}{#1}}}
\newtcolorbox{attention}[1][]{popbase,colframe=poppink,before upper={\popboxhead{Attention}{poppink}{#1}}}
\newtcolorbox{experience}[1][]{popbase,colframe=popblue,colback=popblueL,before upper={\popboxhead{Expérience}{popblue}{#1}}}
% « Le savais-tu ? » : carte violette pleine (contexte, culture, Cameroun)
\newtcolorbox{savaistu}[1][]{popbase,colback=poppurple,colframe=popink,
  fontupper=\popsemi\fontsize{10.4}{14.2}\selectfont\color{white},
  before upper={\poppill{Le savais-tu\,?}{white}{poppurple}\ifx\relax#1\relax\else\hspace{6pt}{\popxbold\color{white}#1}\fi\par\vspace{7pt}}}
% Exercice résolu : énoncé en haut, \tcblower puis la solution sur fond crème
\newcounter{popexo}\newcounter{popexr}
\newtcolorbox{exoresolu}[1][]{popbase,colframe=poporange,colbacklower=popcream,
  segmentation style={popink,line width=1.2pt,dashed},
  before upper={\stepcounter{popexr}\popboxhead{Exercice résolu \thepopexr}{poporange}{#1}},
  before lower={\poppill{Solution}{popink}{popcream}\par\vspace{6pt}}}
% Exercice (non résolu en place) — la correction est dans la partie Corrigés
\newtcolorbox{exercice}[1][]{popbase,colframe=popink,
  before upper={\stepcounter{popexo}\popboxhead{Exercice \thepopexo}{popink}{#1}}}
% Corrigé
\newtcolorbox{corrige}[1][]{popbase,colframe=popgreen,colback=popgreenL,
  before upper={\popboxhead{Corrigé}{popgreen}{#1}}}
% Objectifs / prérequis / bilan
\newtcolorbox{objectifs}{popbase,colframe=popink,colback=poporangeL,
  before upper={\popboxhead{Objectifs de la leçon}{popink}{}}}
\newtcolorbox{prerequis}{popbase,colframe=popink,colback=popblueL,
  before upper={\popboxhead{Prérequis}{popblue}{}}}
\newtcolorbox{bilan}{popbase,colframe=popink,colback=white,boxrule=2.8pt,
  before upper={\popboxhead{Fiche bilan}{poporange}{L'essentiel à connaître pour l'examen}}}
% Figure encadrée avec légende
\newtcolorbox{popfigure}[1][]{enhanced,breakable=false,colback=white,colframe=popline,boxrule=1.4pt,arc=8pt,
  left=10pt,right=10pt,top=10pt,bottom=8pt,before skip=10pt,after skip=12pt,halign=center,
  lower separated=false,halign lower=center,
  fontlower=\popsemi\fontsize{9}{11.5}\selectfont\color{popmuted},
  #1}
\newcommand{\legende}[1]{\par\vspace{6pt}{\popsemi\fontsize{9}{11.5}\selectfont\color{popmuted}#1}}

% ============================================================
% Signature OnBuch+
% ============================================================
\newcommand{\poplogoinline}[1]{{\poplogo\fontsize{#1}{#1}\selectfont\color{popink}OnBuch\color{poporange}+}}
\newcommand{\signatureonbuch}{%
  \begin{tcolorbox}[enhanced,colback=popink,colframe=popink,boxrule=2.2pt,arc=10pt,
      drop shadow=poporange,left=14pt,right=14pt,top=10pt,bottom=10pt,before skip=0pt,after skip=0pt]
    \tikz[baseline=-0.7ex]\node[circle,fill=popgreen,draw=popcream,line width=1.3pt,minimum size=22pt,inner sep=0]{\popcheck{white}};%
    \hspace{9pt}\begin{minipage}[c]{0.62\linewidth}
      {\popxbold\fontsize{12}{13}\selectfont\color{popcream}Signé\ \textbullet\ {\poplogo OnBuch\color{poporange}+}}\par\vspace{2pt}
      {\raggedright\popsemi\fontsize{8}{10.5}\selectfont\color{popline}\MakeUppercase{Équipe pédagogique OnBuch+}\\\MakeUppercase{Conforme au programme officiel MINESEC}\par}
    \end{minipage}\hfill
    {\poplogo\fontsize{22}{22}\selectfont\color{popcream}OnBuch\color{poporange}+}
  \end{tcolorbox}%
}

% ============================================================
% Page de garde
% #1 titre · #2 matière · #3 niveau · #4 module · #5 n° leçon · #6 durée ·
% #7 description / objectifs (texte ou itemize)
% ============================================================
\newcommand{\pagegarde}[7]{%
  \thispagestyle{empty}
  \newgeometry{a4paper,margin=1.7cm,top=1.6cm,bottom=1.4cm}%
  \begin{tikzpicture}[remember picture,overlay]
    \fill[poporange!18] ([xshift=-3.2cm,yshift=-3.6cm]current page.north east) circle (4.6cm);
    \fill[poppurple!14] ([xshift=2.4cm,yshift=7.6cm]current page.south west) circle (3.8cm);
  \end{tikzpicture}%
  {\poplogo\fontsize{30}{30}\selectfont\color{popink}OnBuch\color{poporange}+}\hfill
  \raisebox{0.6ex}{\poppill{Cours signé \textbullet\ Premium}{popink}{popcream}}\par
  \vspace{1pt}\popsquiggle{poppurple}\par
  \vspace{8pt}
  \begin{tcolorbox}[enhanced,colback=poporange,colframe=popink,boxrule=2.8pt,arc=13pt,
      drop shadow=popink,left=18pt,right=18pt,top=12pt,bottom=13pt,after skip=10pt]
    \poppill{#2 \textbullet\ #3}{popink}{popcream}\hspace{5pt}\poppill{#4}{white}{popink}\par\vspace{8pt}
    {\popdisplay\fontsize{14}{14}\selectfont\color{popink}LEÇON #5}\par\vspace{4pt}
    {\raggedright\hyphenpenalty=10000\exhyphenpenalty=10000\popdisplay\fontsize{25}{27}\selectfont\color{popcream}\MakeUppercase{#1}\par}
    \vspace{8pt}
    {\popxbold\fontsize{9.5}{9.5}\selectfont\color{popink}\MakeUppercase{Durée conseillée : #6}}
  \end{tcolorbox}
  \begin{tcolorbox}[enhanced,colback=white,colframe=popink,boxrule=2.2pt,arc=10pt,
      drop shadow=popink,left=16pt,right=16pt,top=10pt,bottom=10pt,after skip=8pt]
    \poppill{Dans cette leçon}{poppurple}{white}\par\vspace{5pt}
    \popsemi\fontsize{9.3}{12.6}\selectfont\color{popink}#7
  \end{tcolorbox}
  \noindent\begin{minipage}[t]{0.487\linewidth}
    \begin{tcolorbox}[enhanced,colback=popgreen,colframe=popink,boxrule=2.2pt,arc=10pt,
        drop shadow=popink,left=11pt,right=11pt,top=10pt,bottom=10pt,height=3.1cm,valign=center]
      \tcbox[colback=white,colframe=popink,boxrule=1.3pt,arc=5pt,boxsep=0pt,nobeforeafter,
        left=3pt,right=3pt,top=3pt,bottom=3pt,valign=center]{\qrcode[height=1.9cm]{https://play.google.com/store/apps/details?id=com.luvvix.onbuch&hl=fr}}%
      \hspace{8pt}\begin{minipage}[c]{0.5\linewidth}
        {\popdisplay\fontsize{11}{12.5}\selectfont\color{white}\MakeUppercase{Télécharge OnBuch}}\par\vspace{4pt}
        {\raggedright\popsemi\fontsize{8.4}{11}\selectfont\color{white}Gratuit \textbullet\ Google Play\\Tuteur IA, annales, concours, résultats\par}
      \end{minipage}
    \end{tcolorbox}
  \end{minipage}\hfill
  \begin{minipage}[t]{0.487\linewidth}
    \begin{tcolorbox}[enhanced,colback=poppurple,colframe=popink,boxrule=2.2pt,arc=10pt,
        drop shadow=popink,left=11pt,right=11pt,top=10pt,bottom=10pt,height=3.1cm,valign=center]
      \tikz[baseline=-0.7ex]\node[circle,fill=white,draw=popink,line width=1.3pt,minimum size=1.6cm,inner sep=0]{\popdisplay\fontsize{24}{24}\selectfont\color{poppurple}+};%
      \hspace{8pt}\begin{minipage}[c]{0.6\linewidth}
        {\popdisplay\fontsize{11}{12.5}\selectfont\color{white}\MakeUppercase{Passe à OnBuch+}}\par\vspace{4pt}
        {\raggedright\popsemi\fontsize{8.4}{11}\selectfont\color{white}Tous les cours signés, Tuteur IA illimité, fiches PDF — onglet OnBuch+ de l'app\par}
      \end{minipage}
    \end{tcolorbox}
  \end{minipage}
  \vspace{8pt}
  \popcontenu
  \vfill
  \signatureonbuch
  \restoregeometry
  \newpage
}

% Rangée « Ce cours contient » (page de garde)
\newcommand{\poptile}[3]{% #1 couleur #2 gros texte #3 légende
  \begin{tcolorbox}[enhanced,colback=#1,colframe=popink,boxrule=2pt,arc=9pt,shadow={2.5pt}{-2.5pt}{0pt}{popink},
      left=6pt,right=6pt,top=9pt,bottom=9pt,halign=center,valign=center,height=1.75cm,width=\linewidth,nobeforeafter]
    {\popdisplay\fontsize{12.5}{13}\selectfont\color{white}\MakeUppercase{#2}}\par\vspace{3pt}
    {\centering\popsemi\fontsize{7.8}{9.5}\selectfont\color{white}#3\par}
  \end{tcolorbox}}
\newcommand{\popcontenu}{%
  \par\noindent{\popxboldls\fontsize{7.5}{7.5}\selectfont\color{popmuted}CE COURS SIGNÉ CONTIENT}\par\vspace{7pt}
  \noindent\begin{minipage}[t]{0.235\linewidth}\poptile{popblue}{Cours}{complet, illustré, pas à pas}\end{minipage}\hfill
  \begin{minipage}[t]{0.235\linewidth}\poptile{poporange}{Méthodes}{et exercices résolus}\end{minipage}\hfill
  \begin{minipage}[t]{0.235\linewidth}\poptile{poppink}{Exercices}{type examen + corrigés}\end{minipage}\hfill
  \begin{minipage}[t]{0.235\linewidth}\poptile{popgreen}{Fiche bilan}{l'essentiel à retenir}\end{minipage}\par}

% Sommaire
\newtcolorbox{sommaire}{enhanced,colback=white,colframe=popink,boxrule=2.2pt,arc=10pt,
  drop shadow=popink,left=18pt,right=18pt,top=13pt,bottom=13pt,before skip=6pt,after skip=20pt,
  before upper={\poppill{Sommaire}{poppurple}{white}\par\vspace{9pt}\popsemi\fontsize{10.6}{14.5}\selectfont\color{popink}}}

% Signature de fin de cours — poussée tout en bas de la dernière page
\newcommand{\findecours}{%
  \par\vfill\nopagebreak
  \begin{center}\popsquiggle{poporange}\end{center}\vspace{4pt}
  \signatureonbuch
}
\AtBeginDocument{\def\llbracket{\mathopen{[\mkern-3.5mu[}}\def\rrbracket{\mathclose{]\mkern-3.5mu]}}} % intervalles d'entiers (glyphe ⟦ absent de la police)
% Glyphes absents de Fira Math : redéfinis à partir de glyphes disponibles
\AtBeginDocument{%
  \def\setminus{\mathbin{\backslash}}\let\smallsetminus\setminus
  \def\vdots{\mathord{\vbox{\baselineskip3.2pt\lineskiplimit0pt\kern4pt\hbox{.}\hbox{.}\hbox{.}}}}%
  \def\ddots{\mathinner{\mkern1mu\raise6.5pt\hbox{.}\mkern2mu\raise3.6pt\hbox{.}\mkern2mu\raise0.7pt\hbox{.}\mkern1mu}}%
  \def\triangle{\mathord{\scalebox{0.9}{$\symup{\Delta}$}}}%
  \def\nabla{\mathord{\raisebox{\depth}{\scalebox{1}[-1]{$\symup{\Delta}$}}}}%
  \def\checkmark{\popcheck{popdark}}%
  \def\star{\mathbin{\ast}}\def\bigstar{\mathord{\ast}}%
  \def\diamond{\mathbin{\scalebox{0.6}{$\Diamond$}}}%
  \def\bigcup{\mathop{\mathchoice{\raisebox{-0.3em}{\scalebox{1.7}{$\cup$}}}{\scalebox{1.25}{$\cup$}}{\cup}{\cup}}\displaylimits}%
  \def\bigcap{\mathop{\mathchoice{\raisebox{-0.3em}{\scalebox{1.7}{$\cap$}}}{\scalebox{1.25}{$\cap$}}{\cap}{\cap}}\displaylimits}%
  \def\lesseqqgtr{\mathrel{\lessgtr}}\def\bigoplus{\mathop{\oplus}\displaylimits}%
}
\providecommand{\ssi}{si et seulement si} % abréviation fréquente
\AtBeginDocument{\providecommand{\wideparen}{\overparen}} % arc de cercle

%% FIGFIX-BEGIN v5 — correctifs globaux des figures (docs/onbuch-generation/tools/figfix)
\usepackage{adjustbox}\usepackage{etoolbox}\tcbuselibrary{raster}
% toute figure TikZ/pgfplots plus large que la ligne est réduite à la largeur disponible
\BeforeBeginEnvironment{tikzpicture}{\begin{adjustbox}{max width=\linewidth}}
\AfterEndEnvironment{tikzpicture}{\end{adjustbox}}
% légendes pgfplots lisibles par-dessus les courbes
\pgfplotsset{every axis/.append style={legend style={fill=white,fill opacity=0.92,text opacity=1,draw=black!15,inner sep=3pt}}}
% étiquettes d'arêtes (syntaxe edge["..."]) avec fond blanc
\tikzset{every edge quotes/.append style={fill=white,inner sep=1.2pt}}
%% FIGFIX-END
\begin{document}

\pagegarde{Economic Life and Occupations}{Anglais}{1re Tech}{Anglais – classe de Première technique}{N02}{13 séances (fiche de progression harmonisée)}{Cette leçon plonge l'apprenant dans la vie économique quotidienne du Cameroun : achats et ventes au marché, gestion des finances personnelles et découverte des métiers. À travers des activités de speaking, listening, reading, grammar et writing, l'élève maîtrise le vocabulaire commercial, les requêtes polies, les articles partitifs et les quantifieurs pour communiquer efficacement dans des situations économiques réelles.
\vspace{4pt}
\begin{multicols}{2}\small
\begin{itemize}
\item À la fin de cette leçon, je sais échanger des informations sur l'achat et la vente d'articles et de denrées dans un magasin ou un marché proche
\item À la fin de cette leçon, je sais formuler des requêtes et suggestions polies avec 'let me + infinitif', 'please', 'could you', 'can you'
\item À la fin de cette leçon, je sais utiliser correctement les articles partitifs (a bunch of, a cup of, a loaf of...) pour quantifier les denrées
\item À la fin de cette leçon, je sais comprendre des textes oraux sur la gestion des finances personnelles
\item À la fin de cette leçon, je sais employer les adjectifs nominalisés et les structures 'So do I' / 'Neither did he' pour exprimer l'accord et le désaccord
\item À la fin de cette leçon, je sais utiliser les déterminants et quantifieurs (some, any, much, all, most, few, a few, little, a little, each, every, many, both, no, none)
\end{itemize}\end{multicols}}

\begin{sommaire}
\begin{multicols}{2}
\begin{enumerate}
\item Buying and Selling at the Market (Speaking \& Vocabulary)
\item Polite Requests and Partitive Articles (Grammar)
\item Managing Personal Finances (Listening \& Grammar)
\item Determiners and Quantifiers — Parts 1, 2 \& 3 (Grammar)
\item Developing Interest in a Trade or Profession (Reading \& Vocabulary)
\item Writing about Trades and Professions (Writing \& Integration)
\item Activité d'intégration
\item Méthodes et exercices résolus
\item Exercices
\item Corrigés des exercices
\item Fiche bilan
\end{enumerate}
\end{multicols}
\end{sommaire}

\begin{objectifs}
\begin{itemize}
\item À la fin de cette leçon, je sais échanger des informations sur l'achat et la vente d'articles et de denrées dans un magasin ou un marché proche
\item À la fin de cette leçon, je sais formuler des requêtes et suggestions polies avec 'let me + infinitif', 'please', 'could you', 'can you'
\item À la fin de cette leçon, je sais utiliser correctement les articles partitifs (a bunch of, a cup of, a loaf of...) pour quantifier les denrées
\item À la fin de cette leçon, je sais comprendre des textes oraux sur la gestion des finances personnelles
\item À la fin de cette leçon, je sais employer les adjectifs nominalisés et les structures 'So do I' / 'Neither did he' pour exprimer l'accord et le désaccord
\item À la fin de cette leçon, je sais utiliser les déterminants et quantifieurs (some, any, much, all, most, few, a few, little, a little, each, every, many, both, no, none)
\item À la fin de cette leçon, je sais lire et comprendre des textes sur le développement et le partage d'intérêt pour un métier ou une profession
\item À la fin de cette leçon, je sais rédiger une composition sur le développement et le partage d'intérêt pour un métier, en justifiant ma démarche
\item À la fin de cette leçon, je sais appliquer ces notions à des situations concrètes de la vie économique camerounaise
\end{itemize}
\end{objectifs}

\begin{prerequis}
\begin{itemize}
\item Vocabulaire de base de la vie quotidienne acquis en classe de Seconde (food, money, shops)
\item Structures simples du présent et du passé (simple present, simple past)
\item Notions élémentaires sur les noms dénombrables et indénombrables
\item Expressions de politesse de base (please, thank you, excuse me)
\item Capacité à tenir un dialogue simple appris en début d'année
\end{itemize}
\end{prerequis}

\newpage

\coursec{Buying and Selling at the Market (Speaking \& Vocabulary)}

\textbf{Situation d'accroche.} It is Saturday morning, 7 a.m., at Mokolo Market in Yaoundé. Amina, a student in \emph{Première technique}, is walking behind her mother between the noisy stalls. Her mother gives her a shopping list: \emph{a bunch of plantains}, \emph{a cup of groundnuts}, and \emph{a bottle of palm oil}. But today, there is a challenge: the seller is a woman from Bamenda who speaks only English. Amina must ask the prices, bargain, pay and thank her — all in English! In the evening, her father, a carpenter, will ask her to help him manage the small family budget and to write a short text about his trade.

\textbf{Problématique :} How can Amina use English to buy, sell, manage her money and talk about her future profession? This lesson gives her — and you — all the necessary linguistic tools. We start with the first skill: \cle{buying and selling at the market}.

\courssub{Discovering the Market: Key Vocabulary}

Before you can buy anything, you must know the names of the goods. At a Cameroonian market, you find \cle{foodstuff} (denrées alimentaires) and \cle{articles} (articles, marchandises). Study the following lists carefully.

\begin{definition}[Key words]
A \cle{market} is a place where people buy and sell goods. A \cle{stall} (étal) is the small table or stand where a seller displays goods. The \cle{seller} or \cle{vendor} (vendeur) sells; the \cle{customer} (client) buys. The \cle{goods} (marchandises) are the things sold. A \cle{scale} (balance) is used to weigh goods, and the \cle{cash box} (caisse) is where the seller keeps the money.
\end{definition}

\textbf{Foodstuff at the Cameroonian market (les denrées) :}

\begin{center}
\begin{tabularx}{\linewidth}{|l|X|l|X|}
\hline
\rowcolor{popblueL} \textbf{English} & \textbf{Français} & \textbf{English} & \textbf{Français} \\
\hline
plantains & plantains & tomatoes & tomates \\
\hline
cassava & manioc & palm oil & huile de palme \\
\hline
groundnuts & arachides & dried fish & poisson fumé/séché \\
\hline
maize & maïs & cocoyams & taro/macabo \\
\hline
beans & haricots & pepper & piment \\
\hline
rice & riz & onions & oignons \\
\hline
\end{tabularx}
\end{center}

\textbf{Other articles (autres articles) :} second-hand clothes (vêtements d'occasion), shoes (chaussures), soap (savon), cooking pots (marmites), fabric / \emph{wax} (pagne), school supplies (fournitures scolaires).

\textbf{People and places :} customer (client), seller/vendor (vendeur), hawker (vendeur ambulant), stall (étal), market square (place du marché), price (prix), change (la monnaie), receipt (reçu).

\begin{popfigure}
\begin{center}
\begin{tikzpicture}[font=\small]
% étal (table)
\draw[fill=popgoldL, thick] (0,1.6) rectangle (6,2.0);
\draw[thick] (0.3,1.6) -- (0.3,0);
\draw[thick] (5.7,1.6) -- (5.7,0);
% toit
\draw[thick, fill=poporangeL] (-0.4,4.2) -- (3,5.0) -- (6.4,4.2) -- cycle;
\draw[thick] (0.2,4.2) -- (0.2,2.0);
\draw[thick] (5.8,4.2) -- (5.8,2.0);
% marchandises
\draw[fill=popgreenL] (0.6,2.0) ellipse (0.45 and 0.28);
\draw[fill=popgreenL] (1.5,2.0) ellipse (0.45 and 0.28);
\draw[fill=popgreenL] (1.05,2.45) ellipse (0.45 and 0.28);
\draw[fill=poporange] (2.6,2.0) circle (0.3);
\draw[fill=poporange] (3.2,2.0) circle (0.3);
\draw[fill=poporange] (2.9,2.45) circle (0.3);
\draw[fill=poppinkL] (4.2,2.0) rectangle (5.4,2.6);
% balance
\draw[thick] (4.8,2.6) -- (4.8,3.3);
\draw[thick] (4.3,3.3) -- (5.3,3.3);
\draw[thick] (4.3,3.3) -- (4.3,3.0);
\draw[thick] (5.3,3.3) -- (5.3,3.0);
\draw[fill=popblueL] (4.05,2.85) rectangle (4.55,3.0);
\draw[fill=popblueL] (5.05,2.85) rectangle (5.55,3.0);
% vendeuse
\draw[fill=popcream] (3,3.9) circle (0.32);
\draw[thick, fill=popblueL] (2.65,3.55) -- (3.35,3.55) -- (3.5,2.6) -- (2.5,2.6) -- cycle;
% cliente
\draw[fill=popcream] (-2.2,2.9) circle (0.32);
\draw[thick, fill=poppinkL] (-2.55,2.55) -- (-1.85,2.55) -- (-1.7,1.3) -- (-2.7,1.3) -- cycle;
\draw[thick] (-2.5,1.3) -- (-2.5,0.6);
\draw[thick] (-1.9,1.3) -- (-1.9,0.6);
% cash box
\draw[fill=popgold, thick] (6.6,2.0) rectangle (7.6,2.7);
\draw[thick] (6.6,2.45) -- (7.6,2.45);
% labels
\draw[->, popdark] (1.05,3.3) -- (1.05,2.85); \node[above] at (1.05,3.3) {\textbf{goods} (plantains)};
\draw[->, popdark] (4.8,4.0) -- (4.8,3.4); \node[above] at (4.8,4.0) {\textbf{scale}};
\draw[->, popdark] (3.0,5.6) -- (3.0,4.35); \node[above] at (3.0,5.6) {\textbf{seller}};
\draw[->, popdark] (-2.2,3.9) -- (-2.2,3.25); \node[above] at (-2.2,3.9) {\textbf{customer}};
\draw[->, popdark] (7.1,3.4) -- (7.1,2.75); \node[above] at (7.1,3.4) {\textbf{cash box}};
\node[below] at (3,-0.2) {\textbf{stall}};
\end{tikzpicture}
\end{center}
\legende{A market stall: 1. seller, 2. customer, 3. goods, 4. scale, 5. cash box, 6. stall.}
\end{popfigure}

\courssub{Asking Prices: How Much Is It?}

When you want to know the price of something, English offers three very common questions. Learn them by heart:

\begin{aretenir}[Asking the price]
\begin{itemize}
\item \textbf{How much is...?} — used with a singular or uncountable noun: \emph{How much is a bunch of plantains? How much is the palm oil?}
\item \textbf{How much are...?} — used with a plural noun: \emph{How much are the tomatoes?}
\item \textbf{How much does it cost?} / \textbf{How much do they cost?} — \emph{How much does this dress cost?}
\item \textbf{What's the price of...?} — \emph{What's the price of this dried fish?}
\end{itemize}
Answers: \emph{It's 1,500 FCFA. / They're 500 FCFA. / It costs 2,000 FCFA.}
\end{aretenir}

\begin{attention}[Common mistake]
Never say \emph{``How much costs...?''} or \emph{``How much it costs?''} In English questions, you need the auxiliary \textbf{does/do}: \emph{How much \textbf{does it cost}?} Also remember: \emph{How much \textbf{is}} (singular) but \emph{How much \textbf{are}} (plural).
\end{attention}

\courssub{Polite Requests and Partitive Articles (Grammar)}

\begin{propriete}[Polite Requests: Asking and Offering]
In a market context, being polite helps the transaction. We use three main structures:
\begin{itemize}
    \item \textbf{Could you + bare infinitive...?} (Polite request) — \emph{Could you give me a bunch of plantains?}
    \item \textbf{Can you + bare infinitive...?} (Less formal request) — \emph{Can you reduce the price?}
    \item \textbf{Let me + bare infinitive} (Offering to do something) — \emph{Let me help you with that bag.}
    \item \textbf{May I / Can I + bare infinitive...?} (Asking permission) — \emph{May I have a cup of groundnuts?}
\end{itemize}
\textbf{Note:} \emph{Let me} is followed by the bare infinitive (without \emph{to}). \emph{Could/Can/May} are modals followed by bare infinitive.
\end{propriete}

\begin{definition}[Partitive Articles / Quantifiers (Containers)]
When buying uncountable goods (rice, oil, groundnuts) or items sold in groups, we use a \cle{partitive structure}: \textbf{Container/Quantity word + of + Noun}.
\begin{center}
\begin{tabularx}{\linewidth}{|l|X|l|X|}
\hline
\rowcolor{popblueL} \textbf{Container / Quantity} & \textbf{Example} & \textbf{Container / Quantity} & \textbf{Example} \\
\hline
a \cle{bunch of} & a bunch of plantains / bananas & a \cle{cup of} & a cup of groundnuts / rice / beans \\
\hline
a \cle{bottle of} & a bottle of palm oil / water / soda & a \cle{piece of} & a piece of dried fish / meat / soap \\
\hline
a \cle{kilo of} / \cle{kg of} & a kilo of tomatoes / onions / maize & a \cle{bag of} & a bag of rice / cement / flour \\
\hline
a \cle{tin of} & a tin of tomatoes / milk / sardines & a \cle{packet of} & a packet of biscuits / sugar / spices \\
\hline
\end{tabularx}
\end{center}
\textbf{Plural:} \emph{Two cups of rice, three bottles of oil, five pieces of fish.}
\end{definition}

\begin{attention}[Common Mistake: Partitives]
Do not say \emph{``a groundnuts''} or \emph{``two rices''}. Uncountable nouns need a container: \emph{``a cup of groundnuts'', ``two cups of rice''}. Countable nouns in groups need a collective: \emph{``a bunch of plantains'', ``two bunches of bananas''}.
\end{attention}

\begin{exemplebox}[From Imperative to Polite Request]
\textbf{Customer (Imperative):} Give me two cups of groundnuts. \textbf{(Rude)}\par
\textbf{Customer (Polite):} \emph{Could you give me two cups of groundnuts, please?} \textbf{(Standard)}\par
\textbf{Customer (Polite):} \emph{Can I have two cups of groundnuts, please?} \textbf{(Standard)}\par
\textbf{Seller (Offer):} \emph{Let me weigh them for you.} \textbf{(Helpful)}
\end{exemplebox}

\begin{exercice}[Grammar Practice: Polite Requests and Partitives]
Rewrite the following sentences using a polite request structure (\emph{Could you...}, \emph{Can I have...}, \emph{Let me...}) and the correct partitive expression.
\begin{enumerate}
    \item I want a plantains. (Use: bunch of)
    \item Give me the palm oil. (Use: bottle of / Could you)
    \item I take two rice. (Use: cups of / Can I have)
    \item Help me carry the bag. (Use: Let me)
    \item Reduce the price! (Use: Could you)
\end{enumerate}
\end{exercice}

\begin{corrige}[Exercice: Grammar Practice]
\begin{enumerate}
    \item \textbf{Can I have a bunch of plantains, please?} (or \emph{Could you give me a bunch of plantains?})
    \item \textbf{Could you give me a bottle of palm oil, please?}
    \item \textbf{Can I have two cups of rice, please?}
    \item \textbf{Let me help you carry the bag.} (or \emph{Let me carry the bag for you.})
    \item \textbf{Could you reduce the price, please?}
\end{enumerate}
\end{corrige}

\courssub{Bargaining: Discussing the Price}

\begin{definition}[Bargaining]
\cle{Bargaining} (marchandage) is \textbf{the act of discussing the price of something before buying it}. The buyer proposes a lower price; the seller accepts, refuses, or proposes another price. In French: \emph{négocier le prix}.
\end{definition}

\textbf{The customer's expressions (l'acheteur) :}
\begin{itemize}
\item \emph{That's too expensive!} (C'est trop cher !)
\item \emph{Can you reduce the price?} (Pouvez-vous baisser le prix ?)
\item \emph{I'll give you 500 FCFA.} (Je vous donne 500 FCFA.)
\item \emph{Is that your last price?} (C'est votre dernier prix ?)
\item \emph{Please, add something for me!} (Ajoutez-moi quelque chose, s'il vous plaît !)
\end{itemize}

\textbf{The seller's expressions (le vendeur) :}
\begin{itemize}
\item \emph{Welcome, dear customer!} (Bienvenue, cher client !)
\item \emph{It's fresh! / It's good quality!} (C'est frais ! / C'est de bonne qualité !)
\item \emph{I'm giving you a good price.} (Je vous fais un bon prix.)
\item \emph{Take it or leave it.} (À prendre ou à laisser.)
\item \emph{For you, my friend, 1,200 FCFA — last price!} (Pour vous, 1 200 FCFA, dernier prix !)
\end{itemize}

\textbf{Concluding the transaction (conclure la transaction) :}
\begin{itemize}
\item \emph{Here's your change.} (Voici votre monnaie.)
\item \emph{Thank you, come again!} (Merci, revenez !)
\item \emph{Thank you very much. Have a nice day!} (Merci beaucoup, bonne journée !)
\end{itemize}

\begin{savaistu}[Bargaining in Cameroon]
In Cameroon, bargaining is an expected cultural practice in markets like Mokolo (Yaoundé) or the Central Market of Douala. If you accept the first price immediately, the seller may be surprised! On the contrary, in supermarkets (like Santa Lucia or Mahima), prices are \textbf{fixed} — bargaining is not possible. Knowing when to bargain is part of economic life!
\end{savaistu}

\courssub{Guided Dialogue: At Mokolo Market}

Read this dialogue aloud with a classmate. Amina (A) is the customer; Mrs. Ndi (B) is the seller from Bamenda.

\begin{exemplebox}[Dialogue: Buying plantains and groundnuts]
\textbf{B:} Welcome, dear customer! What do you want today?\par
\textbf{A:} Good morning, Madam. How much is a bunch of plantains?\par
\textbf{B:} It's 1,500 FCFA. It's very fresh!\par
\textbf{A:} That's too expensive! Can you reduce the price? I'll give you 1,200 FCFA.\par
\textbf{B:} No, my dear. But for you, 1,300 FCFA — last price!\par
\textbf{A:} OK, I take it. And how much is a cup of groundnuts?\par
\textbf{B:} A cup is 200 FCFA.\par
\textbf{A:} Fine. \emph{Could you give me two cups, please?}\par
\textbf{B:} Here you are. That's 1,700 FCFA, please.\par
\textbf{A:} Here is 2,000 FCFA.\par
\textbf{B:} Here's your change: 300 FCFA. Thank you, come again!\par
\textbf{A:} Thank you, Madam. Have a nice day!
\end{exemplebox}

\begin{center}
\begin{tabularx}{\linewidth}{|l|X|X|}
\hline
\rowcolor{popblueL} \textbf{Item (article)} & \textbf{First price (FCFA)} & \textbf{Final price (FCFA)} \\
\hline
a bunch of plantains & 1,500 & 1,300 \\
\hline
two cups of groundnuts & 400 & 400 \\
\hline
\textbf{Total paid} & \textbf{1,900} & \textbf{1,700} \\
\hline
Money given / change & 2,000 & 300 \\
\hline
\end{tabularx}
\end{center}

\begin{methode}[How to lead a buying dialogue in 5 steps]
\begin{enumerate}
\item \textbf{Greeting} (saluer) : \emph{Good morning! / Welcome, dear customer!}
\item \textbf{Asking} (demander le prix) : \emph{How much is...? / What's the price of...?}
\item \textbf{Bargaining} (négocier) : \emph{That's too expensive! Can you take 1,200?}
\item \textbf{Paying} (payer) : \emph{Here is 2,000 FCFA. — Here's your change.}
\item \textbf{Thanking} (remercier) : \emph{Thank you, come again! / Have a nice day!}
\end{enumerate}
\end{methode}

\begin{exoresolu}[Mini-dialogue with numbers]
Complete the dialogue with the correct expressions.\par
A: \emph{How much} .............. a bottle of palm oil?\par
B: It's 2,500 FCFA.\par
A: That's .............. expensive! Can you .............. the price? I'll give you 2,000 FCFA.\par
B: OK, for you, 2,200 FCFA. Take it or .............. it!\par
A: Fine, I take it. Here is 2,500 FCFA.\par
B: Here's your .............. : 300 FCFA. Thank you, come again!
\tcblower
\textbf{Solution :} A: How much \textbf{is} a bottle of palm oil? B: It's 2,500 FCFA. A: That's \textbf{too} expensive! Can you \textbf{reduce} the price? I'll give you 2,000 FCFA. B: OK, for you, 2,200 FCFA. Take it or \textbf{leave} it! A: Fine, I take it. Here is 2,500 FCFA. B: Here's your \textbf{change}: 300 FCFA. Thank you, come again!\par
\textbf{Check the calculation :} $2\,500 - 2\,200 = 300$ FCFA of change. Correct!
\end{exoresolu}

\begin{exercice}[Role-play: At the Central Market of Douala]
Work in pairs. Student A is the customer, Student B is the seller. Build a dialogue using the 5-step method. The customer must buy \emph{a cup of groundnuts} (200 FCFA), \emph{three tomatoes} (250 FCFA) and \emph{a piece of dried fish} (1,000 FCFA). The customer must bargain at least once and pay with 2,000 FCFA. Then change roles.
\end{exercice}

\begin{corrige}[Exercice: role-play]
\textbf{Model dialogue :}\par
B: Welcome, dear customer! What do you need?\par
A: Good morning. How much is a cup of groundnuts?\par
B: It's 200 FCFA.\par
A: OK. And how much are the tomatoes?\par
B: Three tomatoes cost 250 FCFA. They're fresh!\par
A: Fine. What's the price of this dried fish?\par
B: It's 1,000 FCFA.\par
A: That's too expensive! Can you reduce the price? I'll give you 800 FCFA.\par
B: For you, 900 FCFA — last price!\par
A: OK, I take everything. How much is the total?\par
B: $200 + 250 + 900 = 1\,350$ FCFA, please.\par
A: Here is 2,000 FCFA.\par
B: Here's your change: 650 FCFA. Thank you, come again!\par
A: Thank you! Have a nice day!
\end{corrige}

\begin{aretenir}[To remember]
At the market, master three blocks of language: (1) \textbf{asking prices} — \emph{How much is/are...? How much does it cost?} ; (2) \textbf{polite requests \& partitives} — \emph{Could you give me...? Can I have a bunch of...? A cup of... A bottle of...} ; (3) \textbf{bargaining} — \emph{That's too expensive! Can you reduce the price?} ; (4) \textbf{closing the deal} — \emph{Here's your change. Thank you, come again!} With these tools, Amina — and you — can buy and sell in English at any Cameroonian market.
\end{aretenir}

\coursec{Polite Requests and Partitive Expressions (Grammar)}

In the previous section, you learned the vocabulary of buying and selling at a nearby store or market: foodstuffs, articles, prices and bargaining expressions. But knowing the words is not enough. When you speak to a customer, a shopkeeper, a teacher or an examiner, \emph{how} you ask is as important as \emph{what} you ask. In English, a direct order like \og Give me bread \fg{} sounds rude and can cost you marks at the Probatoire oral examination. In this section, you will learn two essential grammar tools for the market situation: \cle{polite requests} (\emph{Could you...?}, \emph{Can you...?}, \emph{Please...}, \emph{Let me...}) and \cle{partitive expressions} (\emph{a cup of rice}, \emph{a bunch of plantains}...).

\courssub{Polite Requests: \emph{Could you...?}, \emph{Can you...?}, \emph{Please...}}

\textbf{Intuition: from the market to the classroom}\par

Imagine two customers at Mokolo market in Yaoundé. The first one says: \og Give me two kilos of rice! \fg{} The second one says: \og Could you give me two kilos of rice, please? \fg{} Who will get better service? Obviously the second customer. English, like French, has \emph{degrees of politeness} (des degrés de politesse). The more polite you are, the more you use \emph{question forms} and \emph{modal verbs} instead of direct orders.

\begin{definition}[Polite request]
A \cle{polite request} (requête polie) is a way of asking someone to do something for you without giving a direct order. The main structures are:
\begin{itemize}
\item \textbf{Can you + bare infinitive...?} — informal, friendly (entre amis, registre familier).
\item \textbf{Could you + bare infinitive...?} — polite, standard (forme polie recommandée, registre courant).
\item \textbf{Please + bare infinitive...} or \textbf{Could you..., please?} — adding \emph{please} always increases politeness.
\end{itemize}
The \emph{bare infinitive} (infinitif sans \emph{to}) is the verb in its simple form: \emph{give}, \emph{carry}, \emph{help}, \emph{show}.
\end{definition}

\begin{propriete}[The politeness scale]
From the least polite to the most polite:
\begin{enumerate}
\item \textbf{Direct order:} \emph{Give me a bag of flour.} (rude, to avoid)
\item \textbf{Neutral:} \emph{Can you give me a bag of flour?}
\item \textbf{Polite:} \emph{Could you give me a bag of flour?}
\item \textbf{Very polite:} \emph{Could you give me a bag of flour, please?}
\end{enumerate}
\emph{Could} is the past form of \emph{can}, but here it does not express the past: it expresses \textbf{politeness and hesitation} (politesse et hésitation), which softens the request (adoucit la demande).
\end{propriete}

\begin{popfigure}
\begin{center}
\begin{tikzpicture}[scale=1]
\draw[-{Stealth},very thick,popmuted] (0,0) -- (11,0) node[below left=2pt and 0pt, font=\small\itshape]{increasing politeness};
\fill[popink] (1,0) circle (3pt);
\node[above=6pt, align=center, font=\small, text width=2.6cm] at (1,0) {\textbf{Direct}\\ \emph{Give me rice.}};
\node[below=6pt, font=\footnotesize\itshape, popmuted] at (1,0) {rude};
\fill[poporange] (4.2,0) circle (3pt);
\node[above=6pt, align=center, font=\small, text width=2.8cm] at (4.2,0) {\textbf{Neutral}\\ \emph{Can you help me?}};
\node[below=6pt, font=\footnotesize\itshape, popmuted] at (4.2,0) {friendly};
\fill[popblue] (7.4,0) circle (3pt);
\node[above=6pt, align=center, font=\small, text width=2.8cm] at (7.4,0) {\textbf{Polite}\\ \emph{Could you help me?}};
\node[below=6pt, font=\footnotesize\itshape, popmuted] at (7.4,0) {standard};
\fill[popgreen] (10.2,0) circle (3pt);
\node[above=6pt, align=center, font=\small, text width=2.9cm] at (10.2,0) {\textbf{Very polite}\\ \emph{Could you..., please?}};
\node[below=6pt, font=\footnotesize\itshape, popmuted] at (10.2,0) {examination level};
\end{tikzpicture}
\end{center}
\legende{The politeness scale: from a direct order (rude) to a very polite request with \emph{Could you... please?}. At the Probatoire oral exam, always aim at the green level.}
\end{popfigure}

\begin{exemplebox}[At the store]
\begin{itemize}
\item \emph{Can you show me the price of this shirt?} (to a friend)
\item \emph{Could you weigh this fish for me?} (to the fishmonger)
\item \emph{Could you reduce the price, please?} (bargaining politely)
\item \emph{Please, wrap the meat in paper.} (short polite form)
\end{itemize}
\end{exemplebox}

\courssub{Suggestions and Offers with \emph{Let me + bare infinitive}}

When you want to \emph{offer help} (proposer de l'aide) or \emph{suggest doing something yourself} (suggérer de faire quelque chose soi-même), English uses a very simple and elegant structure.

\begin{definition}[\emph{Let me + bare infinitive}]
\cle{Let me + bare infinitive} is used to offer to do something for someone, or to suggest an action that you will do yourself. It is polite and friendly.
\begin{center}
\textbf{Let me + verb (bare infinitive) + complement.}
\end{center}
Examples: \emph{Let me help you.} / \emph{Let me carry your bag.} / \emph{Let me show you the way to the market.}
\end{definition}

\begin{exemplebox}[Offering help at the market]
\begin{itemize}
\item \emph{Your basket is heavy. \textbf{Let me carry} it for you.}
\item \emph{\textbf{Let me help} you choose the best plantains.}
\item \emph{\textbf{Let me give} you a discount: two tins for 1\,500 FCFA.}
\end{itemize}
Note: to suggest an action for a group (you and others), use \emph{Let's} (= \emph{let us}): \emph{Let's go to the market!} Do not confuse \emph{let me} (je) with \emph{let's} (nous).
\end{exemplebox}

\courssub{Responding to Requests}

A request always needs an answer. Here are the standard responses you must master.

\begin{center}
\begin{tabularx}{\linewidth}{|l|X|X|}
\hline
\rowcolor{popblueL}
\textbf{Type} & \textbf{Expression} & \textbf{Example} \\
\hline
Accepting & \emph{Sure! / Of course! / Certainly! / No problem!} & \emph{— Could you wrap the fish? — Sure!} \\
\hline
Accepting (offer) & \emph{Yes, please. / Thank you.} & \emph{— Let me carry your bag. — Thank you!} \\
\hline
Refusing politely & \emph{I'm sorry, I can't. / I'm afraid not.} & \emph{— Could you reduce the price? — I'm sorry, I can't; it is already cheap.} \\
\hline
\end{tabularx}
\end{center}

\begin{attention}[Refuse politely]
Never answer a request with a dry \emph{No.} Always soften a refusal: \emph{I'm sorry, I can't} + a reason if possible. Example: \emph{I'm sorry, I can't. The shop is closing.} A polite refusal keeps the conversation friendly — and earns marks in speaking tests.
\end{attention}

\begin{exoresolu}[Building a polite dialogue at the market]
Énoncé — Complete this dialogue between a customer (A) and a shopkeeper (B) using: \emph{Could you}, \emph{Let me}, \emph{Sure}, \emph{I'm sorry, I can't}, \emph{please}.

A: \ldots\ldots\ldots\ldots give me a loaf of bread, \ldots\ldots\ldots\ldots?

B: \ldots\ldots\ldots\ldots! Here you are. Anything else?

A: Yes. \ldots\ldots\ldots\ldots reduce the price of these tomatoes?

B: \ldots\ldots\ldots\ldots, \ldots\ldots\ldots\ldots. They are very fresh. But \ldots\ldots\ldots\ldots add one more tomato for free.

A: Thank you very much!
\tcblower
Solution détaillée —

A: \textbf{Could you} give me a loaf of bread, \textbf{please}? (very polite request: question form + \emph{please})

B: \textbf{Sure}! Here you are. Anything else? (accepting a request)

A: Yes. \textbf{Could you} reduce the price of these tomatoes? (polite bargaining)

B: \textbf{I'm sorry}, \textbf{I can't}. They are very fresh. But \textbf{let me} add one more tomato for free. (polite refusal + offer with \emph{let me})

A: Thank you very much!
\end{exoresolu}

\courssub{Partitive Expressions: Quantifying Uncountable Nouns}

\textbf{Intuition: the problem of \emph{rice} and \emph{bread}}\par

Look at these two sentences: \og I bought \emph{a rice} \fg{} and \og I bought \emph{a cup of rice} \fg{}. The first one is wrong. Why? Because \emph{rice}, \emph{bread}, \emph{oil}, \emph{flour} and \emph{meat} are \cle{uncountable nouns} (noms indénombrables / non comptables): you cannot count them directly with \emph{a/an} or with numbers (one, two...). To quantify them, English uses a \emph{container}, a \emph{unit of measurement} or a \emph{shape classifier} + \emph{of}.

\begin{definition}[Partitive expression]
A \cle{partitive expression} (expression partisitive / quantificateur) is an expression used to quantify uncountable nouns. Its structure is:
\begin{center}
\textbf{a / number + container / unit / classifier + \emph{of} + uncountable noun}
\end{center}
Examples: \emph{a cup of rice}, \emph{a loaf of bread}, \emph{two bottles of palm oil}, \emph{a piece of meat}.
\end{definition}

\begin{methode}[Choosing the right partitive]
To choose the correct partitive, follow these steps:
\begin{enumerate}
\item \textbf{Identify the noun:} is it countable (\emph{egg, tomato, orange}) or uncountable (\emph{rice, bread, oil, meat, flour})? If it is countable, simply use \emph{a/an} or a number.
\item \textbf{Look at the shape, the container or the standard unit} of the foodstuff: a liquid usually comes in a \emph{bottle}, \emph{cup}, \emph{glass} or \emph{litre}; bread comes in a \emph{loaf}; fruits on a stem come in a \emph{bunch}; tinned food comes in a \emph{tin}; powder in a \emph{bag} or \emph{packet}.
\item \textbf{Build the expression:} a/number + partitive + \emph{of} + noun. Check the plural: \emph{two bottle\textbf{s} of oil} (the partitive takes the plural, not the uncountable noun).
\end{enumerate}
\end{methode}

\begin{aretenir}[Essential partitives for the market]
\begin{center}
\begin{tabularx}{\linewidth}{|l|X|}
\hline
\rowcolor{popgreenL}
\textbf{Partitive} & \textbf{Used with (examples)} \\
\hline
a cup of & rice, beans, gari, groundnuts, sugar \\
\hline
a bunch of & plantains, bananas, carrots, grapes \\
\hline
a loaf of & bread (\emph{two loaves of bread}) \\
\hline
a bottle of & palm oil, groundnut oil, water, soda \\
\hline
a piece of & meat, fish, soap, chalk, cake \\
\hline
a tin of & sardines, tomato paste, milk, corned beef \\
\hline
a bag of & flour, rice, cement, fertiliser \\
\hline
a packet of & spaghetti, biscuits, tea, sugar, maggi \\
\hline
a kilo of & meat, onions, tomatoes, rice, fish \\
\hline
a litre of & palm oil, kerosene, milk, petrol \\
\hline
\end{tabularx}
\end{center}
\end{aretenir}

\begin{popfigure}
\begin{center}
\begin{tikzpicture}[font=\small, node distance=1.2cm and 3cm]
% Left column: Foodstuffs
\node[draw=popblue, fill=popblueL, rounded corners=2pt, minimum width=3cm, inner sep=4pt] (rice) {rice, beans, gari};
\node[draw=popblue, fill=popblueL, rounded corners=2pt, minimum width=3cm, inner sep=4pt, below=of rice] (plantains) {plantains, bananas};
\node[draw=popblue, fill=popblueL, rounded corners=2pt, minimum width=3cm, inner sep=4pt, below=of plantains] (bread) {bread};
\node[draw=popblue, fill=popblueL, rounded corners=2pt, minimum width=3cm, inner sep=4pt, below=of bread] (oil) {palm oil, water};
\node[draw=popblue, fill=popblueL, rounded corners=2pt, minimum width=3cm, inner sep=4pt, below=of oil] (meat) {meat, fish};
\node[draw=popblue, fill=popblueL, rounded corners=2pt, minimum width=3cm, inner sep=4pt, below=of meat] (sardines) {sardines, tomato paste};
\node[draw=popblue, fill=popblueL, rounded corners=2pt, minimum width=3cm, inner sep=4pt, below=of sardines] (flour) {flour, cement};
% Right column: Partitives
\node[draw=popgreen, fill=popgreenL, rounded corners=2pt, minimum width=3cm, inner sep=4pt, right=of rice] (cup) {a cup of / a bag of};
\node[draw=popgreen, fill=popgreenL, rounded corners=2pt, minimum width=3cm, inner sep=4pt, right=of plantains] (bunch) {a bunch of};
\node[draw=popgreen, fill=popgreenL, rounded corners=2pt, minimum width=3cm, inner sep=4pt, right=of bread] (loaf) {a loaf of};
\node[draw=popgreen, fill=popgreenL, rounded corners=2pt, minimum width=3cm, inner sep=4pt, right=of oil] (bottle) {a bottle of / a litre of};
\node[draw=popgreen, fill=popgreenL, rounded corners=2pt, minimum width=3cm, inner sep=4pt, right=of meat] (piece) {a piece of / a kilo of};
\node[draw=popgreen, fill=popgreenL, rounded corners=2pt, minimum width=3cm, inner sep=4pt, right=of sardines] (tin) {a tin of};
\node[draw=popgreen, fill=popgreenL, rounded corners=2pt, minimum width=3cm, inner sep=4pt, right=of flour] (bag) {a bag of / a packet of};
% Arrows
\foreach \l/\r in {rice/cup, plantains/bunch, bread/loaf, oil/bottle, meat/piece, sardines/tin, flour/bag}{
  \draw[thick, popdark, -{Stealth[length=3mm]}] (\l.east) -- (\r.west);
}
% Labels
\node[font=\bfseries, popblue, above=0.3cm of rice] {Foodstuff (Uncountable)};
\node[font=\bfseries, popgreen, above=0.3cm of cup] {Partitive Expression};
\end{tikzpicture}
\end{center}
\legende{Mapping diagram: each uncountable foodstuff (left) connects to its correct partitive expression(s) (right). The partitive reflects the \emph{shape}, \emph{container} or \emph{standard unit} of the product.}
\end{popfigure}

\begin{exemplebox}[A complete shopping sentence]
\emph{Could you give me \textbf{two bottles of palm oil} and \textbf{three tins of sardines}, please?}

Observe the numbers: the \textbf{partitive} becomes plural (\emph{bottles}, \emph{tins}), but the uncountable nouns (\emph{palm oil}, \emph{sardines} as a product category) never take \emph{-s} in this structure. Compare: \emph{a bag of flour} $\rightarrow$ \emph{five bags of flour} (never \emph{flours}).
\end{exemplebox}

\begin{attention}[The most frequent mistake]
Do not say \og \emph{a rice} \fg{}, \og \emph{a bread} \fg{} or \og \emph{an oil} \fg{}. Uncountable nouns never take \emph{a/an} directly. Always use a partitive expression:
\begin{itemize}
\item Wrong: \emph{I want a bread.} $\rightarrow$ Right: \emph{I want \textbf{a loaf of bread}.}
\item Wrong: \emph{She bought a rice.} $\rightarrow$ Right: \emph{She bought \textbf{a cup of rice} / \textbf{a bag of rice}.}
\item Wrong: \emph{Give me an oil.} $\rightarrow$ Right: \emph{Give me \textbf{a bottle of oil} / \textbf{a litre of oil}.}
\end{itemize}
Also remember: \emph{bread} has no plural (\emph{breads} does not exist in standard English); count \emph{loaves}: \emph{three loaves of bread}.
\end{attention}

\begin{savaistu}[\emph{Could you} vs \emph{Can you} at the exam]
Did you know? \emph{Could you} is more polite than \emph{Can you}, and examiners at the Probatoire and Baccalauréat technique oral examinations pay close attention to politeness. A candidate who says \emph{Could you repeat the question, please?} immediately shows mastery of polite interaction. In Cameroonian markets, politeness also has economic value: a polite customer often gets a \emph{bonus} (a small extra quantity for free, called \og \emph{un petit plus} \fg{} or \og \emph{arrondi} \fg{})!
\end{savaistu}

\courssub{Practice: From Impolite to Polite}

\begin{exercice}[Make it polite]
Rewrite these impolite sentences as polite requests, using \emph{Could you...?} or \emph{Could you..., please?}.
\begin{enumerate}
\item Give me a kilo of meat.
\item Show me your tomatoes.
\item Wrap the fish in paper.
\item Carry my bag.
\item Reduce the price of this shirt.
\end{enumerate}
\end{exercice}

\begin{corrige}[Exercice 1]
\begin{enumerate}
\item \emph{Could you give me a kilo of meat, please?}
\item \emph{Could you show me your tomatoes?}
\item \emph{Could you wrap the fish in paper, please?}
\item \emph{Could you carry my bag, please?} (or, to offer help: \emph{Let me carry your bag.})
\item \emph{Could you reduce the price of this shirt, please?}
\end{enumerate}
\end{corrige}

\begin{exercice}[Choose the right partitive]
Complete with: \emph{a bunch of}, \emph{a loaf of}, \emph{a bottle of}, \emph{a tin of}, \emph{a bag of}, \emph{a cup of}, \emph{a piece of}, \emph{a litre of}.
\begin{enumerate}
\item My mother bought \ldots\ldots\ldots plantains at the market.
\item Please give me \ldots\ldots\ldots bread for breakfast.
\item We need \ldots\ldots\ldots palm oil to cook eru.
\item He drank \ldots\ldots\ldots water after the football match.
\item The carpenter bought \ldots\ldots\ldots flour for his wife.
\item I ate \ldots\ldots\ldots meat with my fufu.
\item She sells \ldots\ldots\ldots rice for 200 FCFA.
\item Open \ldots\ldots\ldots sardines for the sandwich.
\end{enumerate}
\end{exercice}

\begin{corrige}[Exercice 2]
\begin{enumerate}
\item \emph{a bunch of} plantains
\item \emph{a loaf of} bread
\item \emph{a bottle of} palm oil (or \emph{a litre of})
\item \emph{a litre of} water (or \emph{a bottle of})
\item \emph{a bag of} flour
\item \emph{a piece of} meat
\item \emph{a cup of} rice (or \emph{a bag of})
\item \emph{a tin of} sardines
\end{enumerate}
\end{corrige}

\begin{exercice}[Numbers and partitives]
Put the partitive in the plural where necessary.
\begin{enumerate}
\item three (tin) \ldots\ldots\ldots of sardines
\item two (loaf) \ldots\ldots\ldots of bread
\item five (cup) \ldots\ldots\ldots of rice
\item four (bottle) \ldots\ldots\ldots of palm oil
\end{enumerate}
\end{exercice}

\begin{corrige}[Exercice 3]
\begin{enumerate}
\item three \emph{tins} of sardines
\item two \emph{loaves} of bread (irregular plural: loaf $\rightarrow$ loaves)
\item five \emph{cups} of rice
\item four \emph{bottles} of palm oil
\end{enumerate}
Remember: the uncountable noun never changes; only the partitive takes the plural mark.
\end{corrige}

\begin{aretenir}[Key points of this section]
\begin{itemize}
\item Politeness scale: \emph{Give me...} (rude) $<$ \emph{Can you...?} $<$ \emph{Could you...?} $<$ \emph{Could you..., please?}
\item Offer help with \emph{Let me + bare infinitive}: \emph{Let me help you.}
\item Accept: \emph{Sure! / Of course!} — Refuse politely: \emph{I'm sorry, I can't.}
\item Uncountable nouns (\emph{rice, bread, oil, meat, flour}) need a partitive expression: \emph{a cup of rice}, \emph{a loaf of bread}, \emph{two bottles of palm oil}.
\item The plural mark goes on the partitive, never on the uncountable noun.
\end{itemize}
\end{aretenir}

You can now ask for anything at the market politely and quantify any foodstuff correctly. In the next section, \emph{Managing Personal Finances}, you will use these skills in listening activities about money, and discover a new grammar point: adjectives used as nouns and the short answers \emph{So do I} and \emph{Neither did he}.

\coursec{Managing Personal Finances (Listening \& Grammar)}

\courssub{Transition from Polite Requests to Financial Conversations}

In the previous section, we learned how to make \textbf{polite requests} at the market using structures like \emph{``Could you give me a kilo of rice?''} or \emph{``Let me pay now, please.''} We also studied \textbf{partitive articles} to talk about quantities: \emph{a bunch of bananas, a cup of tea, a loaf of bread}.

Now we move from \emph{buying} to \emph{managing} the money that makes those purchases possible. In daily life -- whether you are a student, an apprentice, a shop owner or a farmer -- you must \textbf{plan your income, control your expenses, and save for the future}. This section gives you the English tools to \emph{listen to}, \emph{talk about}, and \emph{write about} personal and small-business finances in the Cameroonian context.

\begin{definition}[Budget]
A \textbf{budget} is a plan that shows how much money you earn (\textbf{income}) and how you will spend it (\textbf{expenses}) over a given period (week, month, year). It helps you avoid \textbf{debt} and build \textbf{savings}.
\end{definition}

\begin{savaistu}[La tontine (njangi) au Cameroun]
La \textbf{tontine} (appelée \emph{njangi} en pidgin/anglais camerounais) est un système d'épargne rotative très répandu. Un groupe de personnes (collègues, voisins, membres d'une église) cotise chaque semaine ou chaque mois. À tour de rôle, chaque membre reçoit la totalité de la cagnotte. Cela permet d'accéder à des sommes importantes sans passer par la banque. On trouve des tontines aussi bien en milieu rural qu'urbain, et de plus en plus via \textbf{Mobile Money} (MTN MoMo, Orange Money).
\end{savaistu}

\courssub{Key Vocabulary: Personal and Small-Business Finances}

\begin{center}
\begin{tabularx}{\linewidth}{|l|X|X|}\hline
\textbf{English term} & \textbf{Meaning / French gloss} & \textbf{Example in context} \\ \hline
income & revenus (salaire, gains, recettes) & \emph{My monthly income is 150,000 FCFA.} \\ \hline
expenses & dépenses (loyer, nourriture, transport, frais scolaires) & \emph{She writes down all her expenses in a notebook.} \\ \hline
budget & budget (plan de revenus et de dépenses) & \emph{We must stick to our budget this month.} \\ \hline
savings & épargne (argent mis de côté) & \emph{His savings helped him buy a new sewing machine.} \\ \hline
debt & dette (argent dû) & \emph{He is trying to pay off his debt to the supplier.} \\ \hline
loan & prêt (emprunt bancaire ou entre particuliers) & \emph{She took a loan to expand her shop.} \\ \hline
tontine / njangi & tontine (épargne rotative communautaire) & \emph{Our tontine meets every Friday at 6 p.m.} \\ \hline
mobile money / MoMo & argent mobile (MTN MoMo, Orange Money) & \emph{I sent the money via MoMo.} \\ \hline
interest & intérêt (taux d'intérêt) & \emph{The loan has a 5\% interest rate.} \\ \hline
repay & rembourser & \emph{He repays 10,000 FCFA every month.} \\ \hline
\emph{to stick to} & respecter, s'en tenir à (un budget) & \emph{You must stick to your budget.} \\ \hline
\emph{to make ends meet} & joindre les deux bouts & \emph{With rising prices, it's hard to make ends meet.} \\ \hline
\end{tabularx}
\end{center}

\courssub{Listening: Managing a Family Budget}

\textbf{Before you listen:} Look at the bar chart below. It shows the typical monthly budget of a Cameroonian family earning \textbf{200,000 FCFA}. Try to guess the percentage for each category.

\begin{popfigure}
\begin{tikzpicture}
\begin{axis}[
    ybar,
    bar width=18pt,
    width=\linewidth,
    height=9cm,
    ymin=0, ymax=50,
    ylabel={Percentage (\% of total income)},
    symbolic x coords={Food, School fees, Transport, Savings, Others},
    xtick=data,
    nodes near coords={\pgfmathprintnumber\pgfplotspointmeta\%},
    nodes near coords align={vertical},
    ymajorgrids=true,
    grid style=dashed,
    title={Monthly Family Budget Breakdown (200,000 FCFA)},
    title style={yshift=0.3cm},
    x tick label style={rotate=0, anchor=north},
    enlarge x limits=0.25,
    legend style={at={(0.5,-0.15)},anchor=north,legend columns=-1},
    popbar1/.style={fill=popblue},
    popbar2/.style={fill=poporange},
    popbar3/.style={fill=popgreen},
    popbar4/.style={fill=poppurple},
    popbar5/.style={fill=poppink},
]
\addplot[popbar1] coordinates {(Food,40)};
\addplot[popbar2] coordinates {(School fees,25)};
\addplot[popbar3] coordinates {(Transport,15)};
\addplot[popbar4] coordinates {(Savings,10)};
\addplot[popbar5] coordinates {(Others,10)};
\end{axis}
\end{tikzpicture}
\legende{Répartition type d'un budget familial mensuel de 200\,000 FCFA. Source : estimation basée sur enquêtes ménages Cameroun (INS, données publiques).}
\end{popfigure}

\textbf{Audio script (Teacher reads aloud / played from device):}

\begin{quote}
\textbf{Dialogue 1: At the kitchen table}
\textsc{Mother}: ``Paul, let's look at this month's budget. Our \textbf{income} is 200,000 FCFA: my salary 120,000 and your father's 80,000.''
\textsc{Father}: ``Right. \textbf{Food} takes 40\% -- that's 80,000 FCFA. \textbf{School fees} for the three children are 25\%, so 50,000 FCFA.''
\textsc{Mother}: ``\textbf{Transport} to work and school is 15\% -- 30,000 FCFA. We put 10\% into \textbf{savings} -- 20,000 FCFA -- through our \textbf{tontine}.''
\textsc{Father}: ``The remaining 10\% (20,000 FCFA) covers \textbf{others}: electricity, water, medicine, airtime. We have no \textbf{debt} this month because we repaid the \textbf{loan} for the roof last month.''
\textsc{Mother}: ``Good. But prices are rising. We must \textbf{stick to the budget}.''
\end{quote}

\begin{quote}
\textbf{Dialogue 2: Small business owner}
\textsc{Alice (tailor)}: ``I earn about 180,000 FCFA a month from sewing. My \textbf{expenses}: rent 40,000, electricity 10,000, fabric and thread 50,000, transport 15,000. That's 115,000 FCFA. I save 30,000 in my \textbf{MoMo} account and put 20,000 in the \textbf{tontine}. The rest (15,000) is for food and emergencies. I took a small \textbf{loan} of 100,000 FCFA to buy a new machine; I \textbf{repay} 10,000 monthly with 2\% \textbf{interest}.''
\end{quote}

\textbf{Listening comprehension exercises:}

\begin{exercice}[Extracting numerical information]
Listen to Dialogue 1 and complete the table.
\begin{center}
\begin{tabularx}{\linewidth}{|l|X|X|}\hline
\textbf{Category} & \textbf{Percentage} & \textbf{Amount (FCFA)} \\ \hline
Income (total) & 100\% & \_\_\_\_\_\_\_\_\_\_ \\ \hline
Food & 40\% & \_\_\_\_\_\_\_\_\_\_ \\ \hline
School fees & \_\_\_\_\_\_\_\_\_\_ & 50,000 \\ \hline
Transport & 15\% & \_\_\_\_\_\_\_\_\_\_ \\ \hline
Savings (tontine) & \_\_\_\_\_\_\_\_\_\_ & 20,000 \\ \hline
Others & 10\% & \_\_\_\_\_\_\_\_\_\_ \\ \hline
Debt / Loan repayment & 0\% & 0 \\ \hline
\end{tabularx}
\end{center}
\end{exercice}

\begin{corrige}[Exercice 1]
\begin{center}
\begin{tabularx}{\linewidth}{|l|X|X|}\hline
\textbf{Category} & \textbf{Percentage} & \textbf{Amount (FCFA)} \\ \hline
Income (total) & 100\% & 200,000 \\ \hline
Food & 40\% & 80,000 \\ \hline
School fees & 25\% & 50,000 \\ \hline
Transport & 15\% & 30,000 \\ \hline
Savings (tontine) & 10\% & 20,000 \\ \hline
Others & 10\% & 20,000 \\ \hline
Debt / Loan repayment & 0\% & 0 \\ \hline
\end{tabularx}
\end{center}
\end{corrige}

\begin{exercice}[True / False / Not Mentioned -- Dialogue 2]
Listen to Dialogue 2 and decide.
\begin{enumerate}
\item Alice earns 200,000 FCFA a month. \hfill [True / False / Not Mentioned]
\item Her biggest expense is rent. \hfill [True / False / Not Mentioned]
\item She saves money in two different ways. \hfill [True / False / Not Mentioned]
\item The loan interest rate is 5\%. \hfill [True / False / Not Mentioned]
\item She repays the loan over 10 months. \hfill [True / False / Not Mentioned]
\end{enumerate}
\end{exercice}

\begin{corrige}[Exercice 2]
\begin{enumerate}
\item False (she earns 180,000 FCFA)
\item False (fabric and thread 50,000 > rent 40,000)
\item True (MoMo account + tontine)
\item False (2\% interest)
\item True (100,000 ÷ 10,000 = 10 months)
\end{enumerate}
\end{corrige}

\courssub{Grammar 1: Adjectives Used as Nouns (The + Adjective)}

When we talk about \textbf{groups of people} in society, we often use \textbf{the + adjective} as a plural noun. The verb is \textbf{plural}.

\begin{propriete}[The + adjective = a group of people]
Structure: \textbf{the} + \textbf{adjective} (no noun after it) $\rightarrow$ plural verb. \\
Examples: \emph{the rich, the poor, the young, the old / the elderly, the unemployed, the homeless, the disabled, the sick.}
\end{propriete}

\begin{exemplebox}[Adjectives as nouns in financial contexts]
\begin{itemize}
\item \emph{The \textbf{rich} \textbf{invest} in land and businesses.} (Les riches investissent...)
\item \emph{The \textbf{poor} \textbf{struggle} to make ends meet.} (Les pauvres peinent...)
\item \emph{The \textbf{unemployed} \textbf{need} access to micro-loans.} (Les chômeurs ont besoin...)
\item \emph{The \textbf{elderly} \textbf{receive} a small pension.} (Les personnes âgées touchent...)
\item \emph{The \textbf{young} \textbf{prefer} Mobile Money for savings.} (Les jeunes préfèrent...)
\end{itemize}
\end{exemplebox}

\begin{attention}[Piège fréquent]
Ne dites pas \emph{``The poors struggle...''} ou \emph{``The riches invest...''}. L'adjectif ne prend \textbf{jamais de -s}. C'est \emph{the poor}, \emph{the rich}, \emph{the unemployed} (pluriel de sens, singulier de forme).
\end{attention}

\courssub{Grammar 2: Agreement Structures -- \emph{So do I} / \emph{Neither did he}}

In conversations about money habits, we often want to say \emph{``Me too''} or \emph{``Me neither''} in a grammatically correct way. English uses \textbf{inversion} with the \textbf{auxiliary verb}.

\begin{propriete}[So / Neither + auxiliary + subject]
\begin{itemize}
\item \textbf{Affirmative agreement:} \textbf{So} + \textbf{auxiliary} + \textbf{subject} \\
  \emph{``I save every month.'' $\rightarrow$ ``\textbf{So do I.}''} (Moi aussi.)
\item \textbf{Negative agreement:} \textbf{Neither} (or \textbf{Nor}) + \textbf{auxiliary} + \textbf{subject} \\
  \emph{``He didn't repay the loan.'' $\rightarrow$ ``\textbf{Neither did she.}''} (Elle non plus.)
\end{itemize}
The auxiliary \textbf{matches the tense and polarity} of the first sentence.
\end{propriete}

\begin{center}
\begin{tabularx}{\linewidth}{|l|X|X|X|}\hline
\textbf{First speaker (tense)} & \textbf{Auxiliary used} & \textbf{So + aux + subject} & \textbf{Neither + aux + subject} \\ \hline
Present Simple (I save) & \textbf{do / does} & So \textbf{do} I. / So \textbf{does} she. & Neither \textbf{do} I. / Neither \textbf{does} he. \\ \hline
Past Simple (I saved) & \textbf{did} & So \textbf{did} I. & Neither \textbf{did} I. \\ \hline
Present Continuous (I am saving) & \textbf{am / is / are} & So \textbf{am} I. & Neither \textbf{am} I. \\ \hline
Past Continuous (I was saving) & \textbf{was / were} & So \textbf{was} I. & Neither \textbf{was} I. \\ \hline
Present Perfect (I have saved) & \textbf{have / has} & So \textbf{have} I. & Neither \textbf{have} I. \\ \hline
Future with \emph{will} (I will save) & \textbf{will} & So \textbf{will} I. & Neither \textbf{will} I. \\ \hline
Modal \emph{can} (I can save) & \textbf{can} & So \textbf{can} I. & Neither \textbf{can} I. \\ \hline
Modal \emph{should} (I should save) & \textbf{should} & So \textbf{should} I. & Neither \textbf{should} I. \\ \hline
\emph{be} as main verb (I am careful) & \textbf{am / is / are} & So \textbf{am} I. & Neither \textbf{am} I. \\ \hline
\end{tabularx}
\end{center}

\begin{exemplebox}[Concrete examples from financial dialogues]
\begin{itemize}
\item A: \emph{``I \textbf{save} 5,000 FCFA every month.''} \\
      B: \emph{``\textbf{So do I.}''} (Present Simple $\rightarrow$ \textbf{do})
\item A: \emph{``My sister \textbf{uses} MoMo for her tontine.''} \\
      B: \emph{``\textbf{So does mine.}''} (3\textsuperscript{e} personne sg $\rightarrow$ \textbf{does})
\item A: \emph{``We \textbf{didn't borrow} money last year.''} \\
      B: \emph{``\textbf{Neither did we.}''} (Past Simple négatif $\rightarrow$ \textbf{did})
\item A: \emph{``He \textbf{hasn't repaid} his debt yet.''} \\
      B: \emph{``\textbf{Neither have I.}''} (Present Perfect négatif $\rightarrow$ \textbf{have})
\item A: \emph{``I \textbf{can't afford} a new machine now.''} \\
      B: \emph{``\textbf{Neither can I.}''} (Modal \emph{can't} $\rightarrow$ \textbf{can})
\item A: \emph{``She \textbf{is} very careful with money.''} \\
      B: \emph{``\textbf{So am I.}''} (\emph{be} verbe principal $\rightarrow$ \textbf{am})
\end{itemize}
\end{exemplebox}

\begin{exoresolu}[Choosing the correct agreement]
Complete each response with \emph{So} or \emph{Neither} + the correct auxiliary + subject.

1. ``I \textbf{keep} a daily expense notebook.'' $\rightarrow$ \_\_\_\_\_\_\_\_\_\_ \textbf{I}. \\
2. ``My brother \textbf{doesn't have} a bank account.'' $\rightarrow$ \_\_\_\_\_\_\_\_\_\_ \textbf{I}. \\
3. ``We \textbf{paid} our tontine contribution yesterday.'' $\rightarrow$ \_\_\_\_\_\_\_\_\_\_ \textbf{we}. \\
4. ``She \textbf{hasn't received} her salary yet.'' $\rightarrow$ \_\_\_\_\_\_\_\_\_\_ \textbf{he}. \\
5. ``They \textbf{will invest} in a new freezer.'' $\rightarrow$ \_\_\_\_\_\_\_\_\_\_ \textbf{we}. \\
6. ``He \textbf{can't read} the contract.'' $\rightarrow$ \_\_\_\_\_\_\_\_\_\_ \textbf{she}. \\
7. ``I \textbf{am} worried about inflation.'' $\rightarrow$ \_\_\_\_\_\_\_\_\_\_ \textbf{I}. \\
8. ``We \textbf{weren't} able to save last month.'' $\rightarrow$ \_\_\_\_\_\_\_\_\_\_ \textbf{they}.

\medskip
\textbf{Solution:}
\begin{enumerate}
\item \textbf{So do I} (Present Simple affirmatif)
\item \textbf{Neither do I} (Present Simple négatif \emph{doesn't} $\rightarrow$ \emph{do})
\item \textbf{So did we} (Past Simple affirmatif)
\item \textbf{Neither has he} (Present Perfect négatif \emph{hasn't} $\rightarrow$ \emph{has})
\item \textbf{So will we} (Future \emph{will})
\item \textbf{Neither can she} (Modal \emph{can't})
\item \textbf{So am I} (\emph{be} verbe principal)
\item \textbf{Neither were they} (Past \emph{be} négatif \emph{weren't})
\end{enumerate}
\end{exoresolu}

\begin{attention}[Erreur très fréquente : \emph{``So I do''} vs \emph{``So do I''}]
\begin{itemize}
\item \emph{``\textbf{So do I}.''} = \textbf{Accord} : ``Moi aussi, j'épargne.'' (Structure inversée : auxiliaire + sujet)
\item \emph{``\textbf{So I do}.''} = \textbf{Prise de conscience / confirmation} : ``Effectivement, c'est vrai !'' (Ordre normal : sujet + auxiliaire). \\
  Exemple : A: ``You save 10,000 a month!'' B: ``\emph{So I do!}'' (Ah oui, c'est vrai ! Je ne m'en rendais pas compte.)
\end{itemize}
\textbf{Au Probatoire / Bac :} Pour exprimer l'accord, l'inversion \textbf{est obligatoire}. \emph{``So I do''} sera considéré comme une erreur.
\end{attention}

\courssub{Grammar 3: Reacting to Statements about Cameroonian Financial Habits}

\textbf{Pair work / Whole class activity.} Read each statement. Student A says the statement. Student B responds with \emph{So...} or \emph{Neither...} using the correct auxiliary, \textbf{according to their own reality}. Then swap roles.

\begin{enumerate}
\item ``Most young people in my neighbourhood \textbf{use} Mobile Money for savings.''
\item ``The \textbf{unemployed} \textbf{find} it hard to get a bank loan.''
\item ``My parents \textbf{didn't have} a tontine when they were young.''
\item ``I \textbf{have never borrowed} money from a microfinance institution.''
\item ``The \textbf{elderly} \textbf{prefer} keeping cash at home.''
\item ``We \textbf{should} teach budgeting in secondary school.''
\item ``Inflation \textbf{has increased} food prices a lot this year.''
\item ``I \textbf{am} good at sticking to my budget.''
\item ``My sister \textbf{can} manage a small business account.''
\item ``We \textbf{weren't} taught personal finance at primary school.''
\end{enumerate}

\begin{exemplebox}[Model exchange]
\begin{quote}
\textsc{Student A}: ``Most young people in my neighbourhood \textbf{use} Mobile Money for savings.'' \\
\textsc{Student B}: ``\textbf{So do my friends.}'' (ou ``\textbf{Neither do mine.}'' si ce n'est pas le cas) \\
\textsc{Student A}: ``The \textbf{unemployed} \textbf{find} it hard to get a bank loan.'' \\
\textsc{Student B}: ``\textbf{So do the young graduates I know.}'' \\
\textsc{Student A}: ``My parents \textbf{didn't have} a tontine when they were young.'' \\
\textsc{Student B}: ``\textbf{Neither did mine.}''
\end{quote}
\end{exemplebox}

\courssub{Integrated Practice: Listening + Grammar + Vocabulary}

\textbf{Task:} Listen to a short monologue (teacher reads). Take notes on \textbf{numbers} and \textbf{habits}. Then write \textbf{4 sentences} reacting with \emph{So / Neither}.

\textbf{Monologue (Teacher reads):}
\begin{quote}
``My name is Ngono. I am a mechanic in Douala. I earn 150,000 FCFA per month. I spend 60,000 on food, 30,000 on rent, 15,000 on transport, and 10,000 on airtime and electricity. I save 20,000 in my tontine every month. I have a loan of 200,000 FCFA for my tools; I repay 15,000 monthly. I don't have a bank account. I use Orange Money. The rich in my area buy land; the poor struggle daily. I am careful with money. I will open a bank account next year.''
\end{quote}

\begin{exercice}[Note-taking + Grammar reaction]
\begin{enumerate}
\item Fill the table while listening.
\begin{center}
\begin{tabularx}{\linewidth}{|l|X|}\hline
\textbf{Item} & \textbf{Amount (FCFA)} \\ \hline
Income & \\ \hline
Food & \\ \hline
Rent & \\ \hline
Transport & \\ \hline
Airtime / Electricity & \\ \hline
Savings (tontine) & \\ \hline
Loan repayment & \\ \hline
Total expenses & \\ \hline
Remaining? & \\ \hline
\end{tabularx}
\end{center}
\item Write 4 reactions using \emph{So / Neither} + auxiliary + \emph{I / we / my brother / etc.} based on \textbf{your own situation}.
\end{enumerate}
\end{exercice}

\begin{corrige}[Exercice 3 -- Note-taking]
\begin{center}
\begin{tabularx}{\linewidth}{|l|X|}\hline
\textbf{Item} & \textbf{Amount (FCFA)} \\ \hline
Income & 150,000 \\ \hline
Food & 60,000 \\ \hline
Rent & 30,000 \\ \hline
Transport & 15,000 \\ \hline
Airtime / Electricity & 10,000 \\ \hline
Savings (tontine) & 20,000 \\ \hline
Loan repayment & 15,000 \\ \hline
Total expenses & 150,000 \\ \hline
Remaining & 0 \\ \hline
\end{tabularx}
\end{center}
\textit{Reactions (examples):}
\begin{enumerate}
\item ``I earn 150,000 FCFA.'' $\rightarrow$ \textbf{Neither do I.} (ou \textbf{So do I.} selon votre cas)
\item ``I save 20,000 in my tontine.'' $\rightarrow$ \textbf{So do I.}
\item ``I don't have a bank account.'' $\rightarrow$ \textbf{Neither do I.}
\item ``I am careful with money.'' $\rightarrow$ \textbf{So am I.}
\end{enumerate}
\end{corrige}

\courssub{Summary of Key Points}

\begin{aretenir}[Ce qu'il faut retenir de cette section]
\begin{itemize}
\item \textbf{Vocabulaire finances} : income, expenses, budget, savings, debt, loan, tontine/njangi, mobile money (MoMo), interest, repay, stick to a budget, make ends meet.
\item \textbf{Adjectifs nominalisés} : \emph{the rich, the poor, the unemployed, the elderly, the young} + verbe au \textbf{pluriel}. Pas de -s à l'adjectif.
\item \textbf{Accord affirmatif} : \textbf{So} + auxiliaire + sujet (inversion). Ex: \emph{So do I. / So does she. / So did we.}
\item \textbf{Accord négatif} : \textbf{Neither} (ou \emph{Nor}) + auxiliaire + sujet (inversion). Ex: \emph{Neither did he. / Neither can I. / Neither have they.}
\item \textbf{Auxiliaire} = celui de la phrase originale (do/does, did, be, have/has, will, can, should...).
\item \textbf{Piège} : \emph{So I do} $\neq$ accord. C'est une confirmation surprise. Pour l'accord : \textbf{inversion obligatoire}.
\item \textbf{Contexte camerounais} : tontine, Mobile Money, petits commerces, budget familial, inflation.
\end{itemize}
\end{aretenir}

\courssub{Homework / Autonomous Practice}

\begin{exercice}[Written production]
Write a short paragraph (80--100 words) about \textbf{your personal budget or that of someone you know}. Use at least \textbf{6 vocabulary items} from the list, \textbf{one adjective-as-noun} (the poor, the young...), and \textbf{two So/Neither structures} (invent a friend's statement and your reaction). Example start: ``My friend Pierre says he saves 10,000 FCFA a month. \textbf{So do I.} But \textbf{the unemployed} in our area...''
\end{exercice}

\begin{exercice}[Grammar consolidation]
Transform each sentence into an agreement response for the subject in brackets.
\begin{enumerate}
\item ``I \textbf{track} my expenses daily.'' (she) $\rightarrow$ \_\_\_\_\_\_\_\_\_\_
\item ``He \textbf{hasn't paid} his tontine.'' (they) $\rightarrow$ \_\_\_\_\_\_\_\_\_\_
\item ``We \textbf{will use} MoMo for the next contribution.'' (I) $\rightarrow$ \_\_\_\_\_\_\_\_\_\_
\item ``She \textbf{was} surprised by the interest rate.'' (he) $\rightarrow$ \_\_\_\_\_\_\_\_\_\_
\item ``They \textbf{can't afford} a new fridge.'' (we) $\rightarrow$ \_\_\_\_\_\_\_\_\_\_
\item ``My brother \textbf{doesn't know} how to make a budget.'' (I) $\rightarrow$ \_\_\_\_\_\_\_\_\_\_
\end{enumerate}
\end{exercice}

\begin{corrige}[Exercice 5]
\begin{enumerate}
\item \textbf{So does she.}
\item \textbf{Neither have they.}
\item \textbf{So will I.}
\item \textbf{So was he.}
\item \textbf{Neither can we.}
\item \textbf{Neither do I.}
\end{enumerate}
\end{corrige}

\coursec{Determiners and Quantifiers — Parts 1, 2 \& 3 (Grammar)}

In the previous section, we listened to texts about managing personal finances. We heard sentences like: \emph{``I spend some money on food, but I don't have any money left at the end of the month.''} and \emph{``How much money do you save every week?''} Words like \cle{some}, \cle{any}, \cle{much} and \cle{every} are called \textbf{quantifiers}. They tell us \emph{how much} or \emph{how many} of something there is. In this section, we will master them step by step, because they are essential for talking about money, prices, goods at the market, and trades.

\courssub{What is a Quantifier?}

\begin{definition}[Quantifier]
A \cle{quantifier} is a word or group of words placed before a noun to express its \textbf{quantity} (how much or how many), without giving an exact number. Examples: \emph{some, any, much, many, few, little, all, most, each, every, both, no, none}.
\end{definition}

Before choosing a quantifier, you must know whether the noun is:
\begin{itemize}
\item \textbf{Countable} (dénombrable): it has a plural form and can be counted — \emph{one trader, two traders; a mango, three mangoes}.
\item \textbf{Uncountable} (indénombrable): it has no plural form and cannot be counted directly — \emph{money, rice, water, information, furniture, advice}. We count them with partitive expressions: \emph{a bag of rice, a glass of water}.
\end{itemize}

\begin{attention}[The golden rule]
Never say \emph{``many money''} or \emph{``much mangoes''}. First ask yourself: \emph{Can I count this noun?} Then choose the correct quantifier.
\end{attention}

\courssub{Part 1: some (of), any (of), much, all (of), most (of)}

\textbf{1. Some and any.} Both are used with plural countable nouns and with uncountable nouns. The difference depends on the type of sentence:

\begin{center}
\begin{tabularx}{\linewidth}{|l|X|X|}
\hline
\rowcolor{popblueL} \textbf{Quantifier} & \textbf{Use} & \textbf{Example} \\
\hline
\cle{some} & Affirmative sentences; offers and requests (when we expect ``yes'') & \emph{I bought some tomatoes. / Would you like some tea? / Can I have some change?} \\
\hline
\cle{any} & Negative sentences and most questions & \emph{She doesn't have any customers today. / Do you have any groundnuts?} \\
\hline
\cle{much} & Uncountable nouns, mainly in negatives and questions & \emph{We don't have much time. / How much money do you need?} \\
\hline
\cle{all (of)} & The whole quantity: 100 \% & \emph{All the traders pay taxes. / He spent all of his salary.} \\
\hline
\cle{most (of)} & The majority: more than half & \emph{Most shopkeepers open early. / Most of the rice is sold.} \\
\hline
\end{tabularx}
\end{center}

\textbf{2. Some vs. some of.} Use \emph{some + noun} for a general quantity: \emph{Some students like accounting.} Use \emph{some of + the/these/my + noun} for a part of a specific group: \emph{Some of the students in my class want to become electricians.} The same logic applies to \emph{any of, much of, all of, most of}: the word \cle{of} appears when the noun is made specific by \emph{the, this, that, these, those} or a possessive (\emph{my, his, our}).

\begin{exemplebox}[At the market in Douala]
\emph{``Madam, do you have \textbf{any} fresh fish today?'' — ``Yes, I have \textbf{some} mackerel, but I don't have \textbf{much} left. \textbf{Most of} it was sold this morning. Would you like \textbf{some}?''}
\end{exemplebox}

\courssub{Part 2: few / a few (of) and little / a little (of)}

This is one of the most delicate points of English grammar, because the small word \cle{a} completely changes the meaning.

\begin{propriete}[Negative vs. positive meaning]
\begin{itemize}
\item \cle{few} + plural countable and \cle{little} + uncountable have a \textbf{negative} meaning: ``presque pas, très peu'' — the quantity is insufficient or disappointing.
\item \cle{a few} + plural countable and \cle{a little} + uncountable have a \textbf{positive} meaning: ``quelques-uns, un peu'' — the quantity is small but enough.
\end{itemize}
\end{propriete}

Compare these sentences carefully:
\begin{itemize}
\item \emph{\textbf{Few} customers came to my shop today, so I am worried.} (presque pas de clients $\rightarrow$ situation négative)
\item \emph{\textbf{A few} customers came to my shop today, so I made some profit.} (quelques clients $\rightarrow$ situation positive)
\item \emph{There is \textbf{little} money in my account; I cannot buy the machine.} (presque pas d'argent)
\item \emph{There is \textbf{a little} money in my account; I can buy some tools.} (un peu d'argent, c'est suffisant)
\end{itemize}

\begin{exoresolu}[Choosing between few, a few, little, a little]
Complete: \emph{``The young carpenter has \ldots{} experience, so his father still supervises his work.''} Then: \emph{``Fortunately, he has \ldots{} good tools, so he can start working.''}
\tcblower
\textbf{Solution.} \emph{Experience} is uncountable and the meaning is negative (he still needs supervision), so we use \textbf{little}: \emph{``The young carpenter has \textbf{little} experience\ldots''}. \emph{Tools} is countable plural and the meaning is positive (``fortunately''), so we use \textbf{a few}: \emph{``\ldots he has \textbf{a few} good tools\ldots''}.
\end{exoresolu}

\courssub{Part 3: each, every, many, much, how many/how much, both (of), no, none (of)}

\textbf{1. Each and every.} Both mean ``chaque'' and both are followed by a \textbf{singular noun and a singular verb}. \emph{Each} insists on the individuals one by one (often for a small, precise group); \emph{every} insists on the whole group (often for large numbers or habits): \emph{Each apprentice received a certificate. / Every shop in this street opens at 7 a.m.}

\begin{attention}[Frequent error]
\cle{Each} and \cle{every} are \textbf{always} followed by a singular noun and a singular verb: \emph{Each student \textbf{has} a uniform.} and NOT \emph{Each students have a uniform.} Similarly: \emph{Every trader \textbf{pays} a tax.}
\end{attention}

\textbf{2. Many and much.} \cle{Many} goes with countable plurals; \cle{much} goes with uncountables. Both are common in questions and negatives; in affirmative sentences, prefer \emph{a lot of}: \emph{There are a lot of customers today.} (more natural than \emph{There are many customers today.})

\textbf{3. How many vs. how much.} This distinction is essential at the market.

\begin{methode}[Choosing between How many and How much]
\begin{enumerate}
\item Identify the noun after the question.
\item Ask: can I count it? (\emph{one egg, two eggs?})
\item If \textbf{yes} (countable plural) $\rightarrow$ use \cle{How many}: \emph{How many eggs do you want?}
\item If \textbf{no} (uncountable) $\rightarrow$ use \cle{How much}: \emph{How much rice do you need?}
\item Special case: for \textbf{prices}, always use \emph{How much}: \emph{How much is this bag of rice? — It costs 25,000 FCFA.}
\end{enumerate}
\end{methode}

\textbf{4. Both (of), no, none (of).}
\begin{itemize}
\item \cle{Both (of)} + plural noun + plural verb = ``les deux'': \emph{Both (of the) mechanics work in Bonabéri.}
\item \cle{No} + noun = ``aucun, pas de'': \emph{There is no bread left. / No customer came.}
\item \cle{None (of)} replaces a noun or introduces a group: \emph{``How many apprentices came?'' — ``None.''} / \emph{None of the traders accepted less than 2,000 FCFA for the bag of rice.} Note: \emph{none of + plural noun} can take a singular or plural verb in modern English (\emph{None of the shops is/are open}), but the singular is preferred in formal writing.
\end{itemize}

\begin{attention}[No vs. none]
\cle{No} is always followed by a noun (\emph{no money}); \cle{none} is never followed directly by a noun — it stands alone or uses \emph{none of}: \emph{None of the money was stolen.} Never say \emph{``none money''}.
\end{attention}

\courssub{Summary Table and the Scale of Quantity}

The following table summarises all the quantifiers of this lesson according to the type of noun.

\begin{popfigure}
\begin{center}
\begin{tabular}{|l|c|c|l|}
\hline
\rowcolor{popblueL} \textbf{Quantifier} & \textbf{Countable (plural)} & \textbf{Uncountable} & \textbf{Meaning / use} \\
\hline
some (of) & yes & yes & affirmative, offers \\
\hline
any (of) & yes & yes & negatives, questions \\
\hline
much & no & yes & negatives, questions \\
\hline
many & yes & no & ``beaucoup de'' (countable) \\
\hline
all (of) & yes & yes & 100 \% \\
\hline
most (of) & yes & yes & the majority \\
\hline
a few (of) & yes & no & ``quelques-uns'' (positive) \\
\hline
few & yes & no & ``presque pas'' (negative) \\
\hline
a little (of) & no & yes & ``un peu'' (positive) \\
\hline
little & no & yes & ``presque pas'' (negative) \\
\hline
each / every & singular only & no & ``chaque'' + singular verb \\
\hline
both (of) & yes (two items) & no & ``les deux'' \\
\hline
no & yes & yes & + noun, ``aucun'' \\
\hline
none (of) & yes & yes & without noun or + of \\
\hline
\end{tabular}
\end{center}
\legende{Summary table: quantifiers classified by type of noun (countable / uncountable) and meaning.}
\end{popfigure}

The scale below shows the progression of quantity from zero to totality. It will help you choose the right word when you speak about stock, customers or money.

\begin{popfigure}
\begin{center}
\begin{tikzpicture}[scale=1]
\draw[-{Stealth}, very thick, popline] (0,0) -- (12.5,0);
\node[below, popdark, font=\small] at (6.25,-0.25) {increasing quantity};
\foreach \x/\lab/\col in {0.4/{none (0\,\%)}/popink, 2.8/{few / little}/poporange, 5.4/{a few / a little}/popgold, 7.6/{some}/popgreen, 9.6/{many / much}/popblue, 11.2/{most}/poppurple, 12.3/{all (100\,\%)}/popdark}{
  \fill[\col] (\x,0) circle (3pt);
  \node[above, popdark, font=\small] at (\x,0.05) {\lab};
}
\foreach \x in {0.4,2.8,5.4,7.6,9.6,11.2,12.3} \draw[popline] (\x,-0.12) -- (\x,0.12);
\end{tikzpicture}
\end{center}
\legende{The scale of quantity: from \emph{none} (nothing) to \emph{all} (the totality). Remember: \emph{few/little} feel negative, \emph{a few/a little} feel positive.}
\end{popfigure}

\courssub{Guided Practice in Context}

\begin{exoresolu}[At Nkoulouloun Market, Douala]
Complete the dialogue with: \emph{some, any, much, many, a few, little, every, most of, how much, both}.\\[4pt]
A: Good morning, madam. Do you have (1) \ldots{} fresh tomatoes?\\
B: Yes, I have (2) \ldots{}, but not (3) \ldots{} — (4) \ldots{} my stock was sold this morning.\\
A: (5) \ldots{} is a kilo?\\
B: 500 FCFA. (6) \ldots{} trader here sells at the same price.\\
A: I only have (7) \ldots{} money left, so I will take (8) \ldots{} kilos.\\
B: (9) \ldots{} my tomatoes and my peppers are fresh. Take (10) \ldots{} peppers too!
\tcblower
\textbf{Solution.} (1) \textbf{any} (question) — (2) \textbf{some} (affirmative) — (3) \textbf{many} (countable plural, negative) — (4) \textbf{Most of} (the majority of a specific stock: \emph{most of + possessive}) — (5) \textbf{How much} (price) — (6) \textbf{Every} (singular noun \emph{trader}) — (7) \textbf{little} (uncountable, negative meaning: almost no money) — (8) \textbf{a few} (countable plural, positive) — (9) \textbf{Both} (two categories) — (10) \textbf{some} (offer).
\end{exoresolu}

\begin{exercice}[Quantifiers in the world of trades]
Choose the correct quantifier.
\begin{enumerate}
\item \emph{Each / Every} apprentice in the workshop \emph{has / have} his own toolbox.
\item There is \emph{few / little} electricity in the workshop today, so the machines are off.
\item \emph{How many / How much} does this second-hand sewing machine cost?
\item \emph{None of / No} the customers complained about the price.
\item The tailor has \emph{a few / a little} regular clients, so his business is doing well.
\item I don't have \emph{some / any} change. Can you pay with mobile money?
\end{enumerate}
\end{exercice}

\begin{corrige}[Exercice : Quantifiers in the world of trades]
\begin{enumerate}
\item \textbf{Each} (or \textbf{Every}) apprentice \ldots \textbf{has} — singular noun and singular verb after \emph{each/every}.
\item \textbf{little} — \emph{electricity} is uncountable, and the meaning is negative (machines are off).
\item \textbf{How much} — for prices we always use \emph{how much}.
\item \textbf{None of} — \emph{none of + the + plural noun}; \emph{no} would need a noun directly (\emph{No customer complained}).
\item \textbf{a few} — \emph{clients} is countable plural and the meaning is positive (business is doing well).
\item \textbf{any} — negative sentence.
\end{enumerate}
\end{corrige}

\begin{savaistu}[Quantifiers at the Probatoire technique]
At the Probatoire and Baccalauréat technique, quantifiers appear \textbf{every year} in the grammar (structure) section of the English paper: gap-fill exercises with \emph{some/any}, \emph{much/many}, \emph{few/a few}, and error-correction items with \emph{each/every + singular}. Mastering the summary table of this lesson guarantees you easy marks!
\end{savaistu}

\begin{aretenir}[Key points to remember]
\begin{itemize}
\item A \cle{quantifier} expresses the quantity of a noun; first check if the noun is countable or uncountable.
\item \cle{some} in affirmatives and offers; \cle{any} in negatives and questions.
\item \cle{few/little} = negative (``presque pas''); \cle{a few/a little} = positive (``quelques-uns / un peu'').
\item \cle{Each} and \cle{every} + singular noun + singular verb.
\item \cle{How many} + countable; \cle{How much} + uncountable and for prices.
\item \cle{No} + noun; \cle{none} alone or \emph{none of + noun}.
\end{itemize}
\end{aretenir}

\coursec{Developing Interest in a Trade or Profession (Reading \& Vocabulary)}

In the previous section, we mastered \textbf{determiners and quantifiers} (\emph{some, any, much, many, few, little, each, every, both, none}) to speak accurately about quantities of materials, tools, money or time. We can now say \emph{``I have \textbf{a few} tools''} or \emph{``There isn't \textbf{much} money left''}. This precision is vital in technical fields. We now move to the heart of this unit, \textbf{Economic Life and Occupations}: we will learn to \textbf{read about}, \textbf{talk about} and \textbf{write about} the path from a young apprentice to a qualified professional in Cameroon. You will acquire the vocabulary of trades and training, practise essential reading techniques (skimming and scanning), and learn structures to express your own professional interests and ambitions.

\courssub{Vocabulary: Trades, Professions and Technical Training}

Technical education in Cameroon prepares you for specific \cle{trades} (métiers manuels ou techniques) and \cle{professions} (professions). Mastering this vocabulary allows you to describe your own training, understand job offers, and discuss the economic life of your community.

\begin{center}
\begin{tabularx}{\linewidth}{|l|X|l|X|}
\hline
\rowcolor{popblueL}
\textbf{Trade / Profession} & \textbf{Description (FR)} & \textbf{Trade / Profession} & \textbf{Description (FR)} \\
\hline
carpenter & menuisier / charpentier & mechanic & mécanicien(ne) \\
tailor & couturier(ère) & electrician & électricien(ne) \\
nurse & infirmier(ère) & trader & commerçant(e) \\
farmer & agriculteur(trice) & hairdresser & coiffeur(euse) \\
plumber & plombier(ère) & welder & soudeur(euse) \\
\hline
\end{tabularx}
\end{center}

\begin{center}
\begin{tabularx}{\linewidth}{|l|X|l|X|}
\hline
\rowcolor{poporangeL}
\textbf{Training Term} & \textbf{Meaning / Context (FR)} & \textbf{Training Term} & \textbf{Meaning / Context (FR)} \\
\hline
apprenticeship & apprentissage (chez un patron) & vocational training & formation professionnelle (école/centre) \\
workshop & atelier (lieu de travail/formation) & skills & compétences, savoir-faire \\
certificate & certificat, diplôme & CAPIET & Certificat d'Aptitude Professionnelle (enseignement technique) \\
BT & Brevet de Technicien & BTS & Brevet de Technicien Supérieur \\
master craftsman & maître artisan (patron formateur) & skilled worker & ouvrier(ère) qualifié(e) \\
\hline
\end{tabularx}
\end{center}

\begin{definition}[Apprenticeship]
An \textbf{apprenticeship} is a period of learning a trade by working with a skilled person (a \emph{master craftsman} or \emph{mentor}) for a fixed time, often combined with theoretical lessons. In Cameroon, it is the main pathway into the informal sector and into many technical careers.
\begin{itemize}
    \item \textbf{Apprentice} (n.): the learner (l'apprenti).
    \item \textbf{Master craftsman / Mentor} (n.): the trainer (le patron / le maître d'apprentissage).
    \item \textbf{Contract} (n.): the agreement defining duration, duties, and pay (le contrat d'apprentissage).
\end{itemize}
\end{definition}

\begin{savaistu}[L'apprentissage informel au Cameroun]
Saviez-vous qu'une \textbf{grande majorité des artisans} au Cameroun (menuisiers, soudeurs, couturiers, mécaniciens) sont formés par la \textbf{voie informelle} ? Le jeune signe un contrat chez un \emph{patron} dans un \emph{workshop} (atelier) pendant 2 à 5 ans. Il apprend \emph{on the job} (sur le tas), souvent sans passer par un lycée technique. Depuis les années 2000, l'État, à travers le MINEFOP (Ministère de l'Emploi et de la Formation Professionnelle), encourage la \emph{certification} de ces acquis (CAPIET, BT) pour valoriser ces parcours et faciliter l'insertion professionnelle.
\end{savaistu}

\begin{popfigure}
\begin{tikzpicture}[
    node style/.style={draw=popdark, thick, fill=popcream, rounded corners, minimum width=2.8cm, minimum height=1.2cm, align=center, font=\small},
    arrow style/.style={-Stealth, thick, popblue},
    label style/.style={font=\tiny, text=popmuted, align=center}
]
% Nodes
\node[node style, align=center] (a) at (0,0){APPRENTICE\\(Apprenti)\\(2--5 ans)};
\node[node style, align=center] (b) at (5,0){SKILLED WORKER\\(Ouvrier Qualifié)\\(CAPIET / BT)};
\node[node style, align=center] (c) at (10,0){MASTER CRAFTSMAN\\(Maître Artisan / Patron)\\(Expérience + BTS / BP)};
\node[node style, align=center] (d) at (15,0){TRAINER / TEACHER\\(Formateur / Enseignant)\\(Pédagogie + Expertise)};

% Arrows
\draw[arrow style] (a.east) -- node[above, label style, align=center]{Learning \\ \emph{on the job} \\ + Theory} (b.west);
\draw[arrow style] (b.east) -- node[above, label style, align=center]{Experience \\ Certification \\ Management} (c.west);
\draw[arrow style] (c.east) -- node[above, label style, align=center]{Transmission \\ Mentoring \\ Pedagogy} (d.west);

% Side annotations
\node[label style, anchor=north, align=center] at (a.south){Entry: BEPC / Probatoire \\ or Primary Cert.};
\node[label style, anchor=north, align=center] at (b.south){Salary \\ Autonomy};
\node[label style, anchor=north, align=center] at (c.south){Own workshop \\ Employs apprentices};
\node[label style, anchor=north, align=center] at (d.south){Technical Lycée \\ CFP};
\end{tikzpicture}
\legende{Parcours type de formation professionnelle au Cameroun : de l'apprenti au formateur. Chaque étape valide des compétences (skills) et permet l'accès à la certification (CAPIET, BT, BTS).}
\end{popfigure}

\begin{popfigure}
\begin{tikzpicture}[
    scale=0.9,
    every node/.style={font=\small},
    tool/.style={draw=popdark, fill=poporangeL, rounded corners=2pt, inner sep=2pt},
    label/.style={font=\tiny, text=popdark, anchor=north}
]
% Workshop layout
\draw[thick, fill=popcream] (0,0) rectangle (14,8); % Room
% Workbenches
\draw[fill=popmuted!30] (1,1) rectangle (6,3);
\draw[fill=popmuted!30] (8,1) rectangle (13,3);
\draw[fill=popmuted!30] (1,5) rectangle (6,7);
\draw[fill=popmuted!30] (8,5) rectangle (13,7);
% Tools on benches
\node[tool] at (2,2) {Plane (Rabot)};
\node[tool] at (4,2) {Saw (Scie)};
\node[tool] at (9,2) {Hammer (Marteau)};
\node[tool] at (11,2) {Chisel (Ciseau à bois)};
\node[tool] at (2,6) {Wood stock (Bois)};
\node[tool] at (4,6) {Measuring tape (Mètre)};
\node[tool] at (9,6) {Sandpaper (Papier de verre)};
\node[tool] at (11,6) {Varnish (Vernis)};
% People
\node[draw, circle, fill=popgreenL, minimum size=12pt] (appr) at (1,4) {};
\node[draw, circle, fill=popblueL, minimum size=12pt] (mast) at (13,4) {};
\node[label] at (1,3.5) {Apprentice};
\node[label] at (13,3.5) {Master};
% Machines
\draw[fill=poppurpleL, thick] (6.5, 2) rectangle (7.5, 4) node[midway, align=center, font=\tiny, text=white] {Circular \\ Saw \\ (Scie circulaire)};
\draw[fill=poppurpleL, thick] (6.5, 5) rectangle (7.5, 7) node[midway, align=center, font=\tiny, text=white] {Planer \\ (Raboteuse)};
% Annotations
\node[label, rotate=90, anchor=south] at (0,4) {Entrance};
\node[label] at (7,0.5) {Carpentry Workshop Layout (Vue de dessus)};
\end{tikzpicture}
\legende{Atelier de menuiserie annoté : postes de travail (workbenches), outillage à main (hand tools : plane, saw, chisel, hammer, measuring tape), machines fixes (circular saw, planer), matières premières (wood stock, sandpaper, varnish) et acteurs (apprentice, master).}
\end{popfigure}

\courssub{Reading Skills: Skimming and Scanning}

In technical professional life, you must read quickly: safety manuals, invoices, technical sheets, or professional profiles. Two techniques are essential:

\begin{methode}[Reading Comprehension Strategy: Skimming \& Scanning]
\textbf{Step 1: Read the Questions First.} Identify \textbf{keywords} (names, dates, places, specific terms like \emph{salary, duration, tool, dream}).
\textbf{Step 2: Skim for Gist (General Idea).} Read the \textbf{title}, the \textbf{first sentence} of each paragraph, and the \textbf{last sentence}. Do not read every word. Ask: \emph{``What is this text mainly about?''}
\textbf{Step 3: Scan for Details.} Move your eyes fast over the text to locate the \textbf{keywords} from Step 1. Read \emph{only} the sentences around those keywords.
\textbf{Step 4: Answer Precisely.} Write short answers based \emph{only} on the text. Use your own words if possible.
\textbf{Step 5: Check.} Does the answer match the question type (Who? Where? How long? Why?)?
\end{methode}

\textbf{Text 1: Ngong's Journey in Bafoussam} \hfill (Skimming Practice)

\begin{exemplebox}[Reading Text 1]
Ngong is an eighteen-year-old apprentice carpenter in Bafoussam, the capital of the West Region. He wakes up at 5:30 a.m. every weekday to reach Master Talla's workshop before 7:00 a.m. The workshop is located near the main market, a busy place full of noise and dust. Ngong loves the smell of fresh wood and the sound of the planer.

His day starts with sweeping the floor and sharpening tools --- chisels, planes, and saws. Then, he watches Master Talla measure and cut expensive mahogany for a client's bedroom set. Ngong is \textbf{keen on} learning how to make dovetail joints, a difficult but strong technique. He says: ``I am \textbf{interested in} traditional designs, but I also want to learn modern CNC machines.''

In the afternoon, he practises cutting small pieces of scrap wood. He makes mistakes, but Master Talla corrects him patiently. ``Patience and precision are the keys,'' the master often repeats. Ngong dreams of obtaining his \textbf{CAPIET} next year and later his \textbf{BT} at the Technical High School. His ultimate goal is to open his own workshop in Foumban, the city of arts, and train other young apprentices.
\end{exemplebox}

\begin{exoresolu}[Skimming Question]
\textbf{Question:} What is the main idea of the text? \\
\textbf{Options:}
A) The history of Bafoussam market.
B) The daily routine and ambitions of an apprentice carpenter.
C) How to use a planer machine.
D) The price of mahogany wood.
\tcblower
\textbf{Solution:}
\textbf{Step 1:} Keywords in options: \emph{history, market, daily routine, ambitions, planer, price}. \\
\textbf{Step 2:} Skim the text. Title: ``Ngong's Journey''. First paragraph: Ngong, 18, apprentice, Bafoussam, workshop. Second paragraph: day starts, sweeping, watching the master, \emph{keen on} learning, \emph{interested in} designs. Third paragraph: practises, makes mistakes, dreams of CAPIET, BT, own workshop. \\
\textbf{Step 3:} The text covers his \textbf{routine} (sweeping, watching, practising) and his \textbf{ambitions} (CAPIET, BT, own workshop, train others). Options A, C and D are only small details, not the main idea. \\
\textbf{Answer:} \textbf{B) The daily routine and ambitions of an apprentice carpenter.}
\end{exoresolu}

\textbf{Text 2: Marie's Tailoring Workshop in Douala} \hfill (Scanning Practice)

\begin{exemplebox}[Reading Text 2]
Marie, 22, works in a busy tailoring workshop in Douala, Akwa district. She completed her \textbf{apprenticeship} three years ago and passed her \textbf{CAPIET} with distinction. Now she is a \textbf{skilled worker} (couturière qualifiée). She earns a good salary and supports her younger brother's school fees.

Marie is \textbf{passionate about} African prints (pagnes) and modern styles. She spends her free time sketching new designs on her phone. She is currently preparing for the \textbf{BT} exam at the Centre de Formation Professionnelle (CFP) in Bonanjo. She attends evening classes three times a week.

Her boss, Madame Ngo, is a \textbf{master craftsman} with 25 years of experience. She encourages Marie to take responsibilities. Last month, Marie managed a big order for a school uniform contract (500 shirts). She organized the cutting, stitching, and finishing teams. The client was satisfied. Marie says: ``I \textbf{dream of becoming} a fashion designer and opening a boutique in Bonapriso. I want to \textbf{share my experience} with future apprentices.''
\end{exemplebox}

\begin{exercice}[Scanning Questions --- Text 2]
Answer these questions by scanning Text 2. Find the keyword first.
\begin{enumerate}
    \item What qualification did Marie obtain three years ago? (Keyword: \emph{three years ago / qualification})
    \item Where does she attend evening classes? (Keyword: \emph{evening classes / where})
    \item How many shirts were in the school uniform order? (Keyword: \emph{shirts / order / number})
    \item What is Marie's dream job? (Keyword: \emph{dream / becoming})
    \item Who is Madame Ngo? (Keyword: \emph{Madame Ngo / who})
\end{enumerate}
\end{exercice}

\begin{corrige}[Exercice Scanning]
\begin{enumerate}
    \item \textbf{The CAPIET} (Sentence: ``...passed her \textbf{CAPIET} with distinction.'')
    \item \textbf{At the CFP in Bonanjo} (Sentence: ``...preparing for the BT exam at the \textbf{Centre de Formation Professionnelle (CFP) in Bonanjo}.'')
    \item \textbf{500 shirts} (Sentence: ``...big order for a school uniform contract (\textbf{500 shirts}).'')
    \item \textbf{A fashion designer} (Sentence: ``I \textbf{dream of becoming} a \textbf{fashion designer}...'')
    \item \textbf{Marie's boss, a master craftsman} (Sentence: ``Her boss, Madame Ngo, is a \textbf{master craftsman} with 25 years of experience.'')
\end{enumerate}
\end{corrige}

\courssub{Expressing and Sharing Interest in a Trade}

To develop a professional network and succeed in interviews or oral exams (Probatoire/Bac technique), you must express your motivation clearly. We use specific structures with gerunds (\textbf{V-ing}) or infinitives.

\begin{propriete}[Structures for Expressing Interest]
\begin{center}
\begin{tabularx}{\linewidth}{|l|X|X|}
\hline
\rowcolor{popgreenL}
\textbf{Structure} & \textbf{Example} & \textbf{Traduction} \\
\hline
\textbf{be interested in} + \textbf{V-ing} / Noun & I am \textbf{interested in learning} welding. & Je m'intéresse à la soudure / Je veux l'apprendre. \\
\textbf{be keen on} + \textbf{V-ing} / Noun & She is \textbf{keen on repairing} engines. & Elle adore réparer des moteurs. \\
\textbf{dream of} + \textbf{V-ing} / Noun & We \textbf{dream of opening} a garage. & Nous rêvons d'ouvrir un garage. \\
\textbf{be passionate about} + \textbf{V-ing} / Noun & He is \textbf{passionate about} coding. & Il est passionné par le codage. \\
\textbf{want / would like / plan} + \textbf{to} + \textbf{Infinitive} & I \textbf{want to become} an electrician. & Je veux / Je voudrais / Je prévois de devenir électricien. \\
\hline
\end{tabularx}
\end{center}
\end{propriete}

\begin{attention}[Piège Majeur : \texttt{interested in} + V-ing]
\textbf{Erreur fréquente :} \emph{``I am interested \textbf{to learn} mechanics.''} (Incorrect !) \\
\textbf{Correct :} \emph{``I am interested \textbf{in learning} mechanics.''} \\
\textbf{Règle :} Après une préposition (\textbf{in, on, of, about, at, to} dans \emph{look forward to}), le verbe prend la forme \textbf{-ing} (gérondif). C'est valable pour \emph{keen on, dream of, passionate about, good at, bad at, look forward to}.
\begin{itemize}
    \item Correct : I am \textbf{keen on driving} trucks.
    \item Correct : He \textbf{dreams of becoming} a nurse.
    \item Correct : We look \textbf{forward to starting} the internship.
\end{itemize}
\end{attention}

\begin{exemplebox}[Profil professionnel type]
Ngong, 18, is an \textbf{apprentice mechanic} in Bamenda. He is \textbf{keen on repairing} motorbikes and \textbf{dreams of opening} his own garage. He is \textbf{passionate about} two-stroke engines and \textbf{interested in} taking a BT course next year.
\end{exemplebox}

\courssub{Sharing Interest: Encouraging Others and Transmitting Skills}

A professional does not only work; they \textbf{share} their passion. This creates the next generation of skilled workers.

\begin{aretenir}[Language for Sharing \& Mentoring]
\begin{itemize}
    \item \textbf{Encourage others:} ``You should try \textbf{carpentry}, it's creative.'' / ``\textbf{Why don't you} join the vocational centre?'' / ``\textbf{Let me show you} how to use this tool.'' (\emph{Let me + bare infinitive} = polite offer/suggestion, seen in Section 2).
    \item \textbf{Share experience:} ``When I was an apprentice, I \textbf{learned that} patience is key.'' / ``\textbf{My advice is to} practise every day.''
    \item \textbf{Train apprentices:} ``I \textbf{teach} them safety rules first.'' / ``I \textbf{supervise} their work on the lathe.''
\end{itemize}
\end{aretenir}

\courssub{Integration Activities: Reading, Vocabulary \& Writing}

\begin{exercice}[Comprehensive Practice]
\textbf{Part A: Vocabulary Matching.} Match the terms (1--6) with definitions (a--f).
\begin{center}
\begin{tabular}{|c|l|c|l|}
\hline
1. Apprenticeship & \hfill & a. & A qualified worker who has finished training. \\
2. Master craftsman & & b. & Period of learning a trade with a skilled person. \\
3. Skilled worker & & c. & Official paper proving a qualification (ex: BT). \\
4. Certificate & & d. & Place where practical work is done. \\
5. Workshop & & e. & Expert who trains apprentices. \\
6. Vocational training & & f. & Organized courses to learn a trade. \\
\hline
\end{tabular}
\end{center}

\textbf{Part B: Grammar in Context.} Complete with the correct form (\emph{V-ing} or \emph{to + infinitive}).
\begin{enumerate}
    \item My brother is interested in \_\_\_\_\_\_\_\_\_\_ (become) a plumber.
    \item We plan \_\_\_\_\_\_\_\_\_\_ (open) a new workshop next month.
    \item She is keen on \_\_\_\_\_\_\_\_\_\_ (design) clothes.
    \item The master asked the apprentice \_\_\_\_\_\_\_\_\_\_ (clean) the machines.
    \item I dream of \_\_\_\_\_\_\_\_\_\_ (get) my BTS in Electrotechnics.
\end{enumerate}

\textbf{Part C: Writing Task (Probatoire Style).}
\textbf{Topic:} ``Write a composition of about 150--180 words on the following subject: \emph{You are a young apprentice in a technical trade (carpentry, tailoring, mechanics, etc.). Describe your daily routine, what you like about the trade, your current qualifications (CAPIET, BT preparation), and your professional dream for the future. Use expressions like \textbf{I am interested in..., I am keen on..., I dream of..., I am passionate about...}.}''
\end{exercice}

\begin{corrige}[Exercice Comprehensive Practice]
\textbf{Part A:}
1--b, 2--e, 3--a, 4--c, 5--d, 6--f.

\textbf{Part B:}
\begin{enumerate}
    \item \textbf{becoming} (interested \textbf{in} + V-ing)
    \item \textbf{to open} (plan \textbf{to} + infinitive)
    \item \textbf{designing} (keen \textbf{on} + V-ing)
    \item \textbf{to clean} (asked somebody \textbf{to} + infinitive)
    \item \textbf{getting} (dream \textbf{of} + V-ing)
\end{enumerate}

\textbf{Part C: Model Composition (Guide)}
\textit{Title: My Path as an Apprentice Electrician}

My name is Paul, I am 19 years old, and I am an apprentice electrician in Yaoundé. I have been in this apprenticeship for two years now, working under Master Etoundi in his workshop at Mvog-Ada.

My daily routine starts early. I arrive at 7:00 a.m. to prepare the tools: testers, pliers, cables, and conduits. I am \textbf{keen on installing} electrical panels because it requires precision and logic. I am also \textbf{interested in} solar energy systems, a growing field in Cameroon. However, the work is hard; sometimes we work late to finish a wiring contract.

Last year, I obtained my \textbf{CAPIET} with a good grade. Currently, I attend evening classes at the CFP to prepare for the \textbf{BT} exam. I \textbf{dream of becoming} a certified electrical technician specialized in renewable energies. I \textbf{am passionate about} bringing electricity to rural villages.

In the future, I want to open my own company and \textbf{train young apprentices} myself. I believe technical skills are the key to development. As my master says: ``A trade in hand finds gold in every land.''
\end{corrige}

\coursec{Writing about Trades and Professions (Writing \& Integration)}

In the previous section, you read texts about young people who developed a strong interest in a trade or profession: a tailor in Yaoundé, a nurse in Bamenda, a mechanic in Douala. You discovered new words such as \cle{apprenticeship} (apprentissage), \cle{vocation} (vocation) and \cle{workshop} (atelier). Now it is your turn to \emph{produce} language. At the Probatoire technique, writing a clear, well-organised composition is a key skill. This final section teaches you how to write about your dream profession, how to correct a classmate's work, and how to use everything you have learnt in this lesson --- buying and selling, polite requests, partitive articles, quantifiers and finance vocabulary --- in one complete real-life situation.

\courssub{The Structure of a Composition}

\textbf{Why structure matters}\par
Imagine you enter a carpenter's workshop. If the tools are everywhere on the floor, you cannot work. A good composition is like a tidy workshop: every idea has its place. Examiners at the Probatoire read hundreds of copies. A composition with a clear plan is easy to read and immediately gets better marks.

\begin{definition}[The three parts of a composition]
A composition always has \textbf{three parts}:
\begin{enumerate}
\item The \cle{introduction}: you \emph{present} the topic (the trade or profession) and announce your plan.
\item The \cle{body} (development): two or three paragraphs with your \emph{arguments} --- why you like the job, the qualities required, the training needed.
\item The \cle{conclusion}: you \emph{summarise} and present your personal project for the future.
\end{enumerate}
\end{definition}

\begin{popfigure}
\begin{tikzpicture}[font=\small]
% Entonnoir : trois trapèzes superposés
\fill[popblueL] (-3.2,4.4) -- (3.2,4.4) -- (2.2,2.9) -- (-2.2,2.9) -- cycle;
\fill[poporangeL] (-2.2,2.6) -- (2.2,2.6) -- (1.3,0.6) -- (-1.3,0.6) -- cycle;
\fill[popgreenL] (-1.3,0.3) -- (1.3,0.3) -- (0.5,-1.4) -- (-0.5,-1.4) -- cycle;
\draw[popdark,thick] (-3.2,4.4) -- (3.2,4.4) -- (2.2,2.9) -- (-2.2,2.9) -- cycle;
\draw[popdark,thick] (-2.2,2.6) -- (2.2,2.6) -- (1.3,0.6) -- (-1.3,0.6) -- cycle;
\draw[popdark,thick] (-1.3,0.3) -- (1.3,0.3) -- (0.5,-1.4) -- (-0.5,-1.4) -- cycle;
\node[popdark] at (0,3.65) {\textbf{INTRODUCTION: present the job}};
\node[popdark] at (0,1.6) {\textbf{BODY: reasons, qualities, training}};
\node[popdark] at (0,-0.55) {\textbf{CONCLUSION}};
\node[popdark] at (0,-1.05) {\textbf{my project}};
% Flèches latérales
\draw[-{Stealth},popblue,very thick] (3.8,4.4) -- (3.8,2.9);
\draw[-{Stealth},poporange,very thick] (2.8,2.6) -- (2.8,0.6);
\draw[-{Stealth},popgreen,very thick] (1.9,0.3) -- (1.9,-1.4);
\node[right,popmuted,align=left] at (4.0,3.65) {general\\idea};
\node[right,popmuted,align=left] at (3.0,1.6) {details and\\arguments};
\node[right,popmuted,align=left] at (2.1,-0.55) {personal\\decision};
\end{tikzpicture}
\legende{The funnel structure of a composition: you move from the general presentation of the profession (wide) to your precise personal project (narrow).}
\end{popfigure}

\textbf{Linking words: the cement of your text}\par
Bricks without cement fall down. Sentences without \cle{linking words} (connecteurs logiques) are difficult to follow. Here are the essential connectors for your compositions:

\begin{center}
\begin{tabularx}{\linewidth}{|l|X|}
\hline
\rowcolor{popblueL} \textbf{Function} & \textbf{Linking words} \\
\hline
Beginning a list of arguments & \cle{first}, \cle{first of all}, \cle{to begin with} \\
\hline
Adding an idea & \cle{moreover}, \cle{furthermore}, \cle{also}, \cle{in addition} \\
\hline
Expressing a contrast & \cle{however}, \cle{but}, \cle{nevertheless} \\
\hline
Expressing a result & \cle{therefore}, \cle{so}, \cle{that is why} \\
\hline
Concluding & \cle{in conclusion}, \cle{to sum up}, \cle{finally} \\
\hline
\end{tabularx}
\end{center}

\begin{exemplebox}[Linking words in action]
\emph{First}, I want to become an electrician because I like working with my hands. \emph{Moreover}, electricians earn a good salary in Cameroon. \emph{However}, the training is long and difficult. \emph{Therefore}, I must work hard at school. \emph{In conclusion}, electricity is my passion and my future.
\end{exemplebox}

\begin{attention}[Agreement: the third person trap]
The most frequent error at the Probatoire is forgetting the \textbf{-s} of the verb at the third person singular (he, she, it, or a singular noun):
\begin{itemize}
\item Correct: \emph{A carpenter \textbf{works} with wood.} / \emph{My sister \textbf{wants} to be a nurse.}
\item Wrong: \emph{A carpenter work with wood.} / \emph{She want to be a nurse.}
\end{itemize}
Rule: singular subject + present simple $\Rightarrow$ verb + \textbf{s}. Check every verb when you revise your copy!
\end{attention}

\courssub{Reusing the Grammar and Vocabulary of the Lesson}

A good composition at the Probatoire shows the examiner that you master the lesson. Deliberately reuse:

\begin{itemize}
\item \textbf{Quantifiers}: \emph{\cle{many} young people choose this trade; \cle{a few} of them succeed immediately; \cle{most} carpenters learn in a workshop; \cle{each} apprentice needs patience; there is \cle{little} unemployment in this sector.}
\item \textbf{Partitive articles}: \emph{a \cle{bunch} of tools, a \cle{cup} of garri, a \cle{piece} of advice.}
\item \textbf{Polite requests}: \emph{Could you show me your workshop, please? Let me explain my project.}
\item \textbf{Trade vocabulary}: \cle{apprentice} (apprenti), \cle{customer} (client), \cle{salary} (salaire), \cle{skill} (compétence), \cle{training} (formation), \cle{workshop} (atelier), \cle{profit} (bénéfice), \cle{budget} (budget).
\end{itemize}

\begin{exemplebox}[Model plan for ``My dream profession'']
\begin{itemize}
\item \textbf{Paragraph 1 --- The job:} name the profession, say what the person does every day. (\emph{A mechanic repairs cars and motorbikes...})
\item \textbf{Paragraph 2 --- Why I like it:} give two or three reasons with \emph{first}, \emph{moreover}, \emph{furthermore}.
\item \textbf{Paragraph 3 --- Training needed:} school, apprenticeship, qualities required (\emph{patient, honest, hard-working}).
\item \textbf{Paragraph 4 --- My project:} your concrete plan for the future (\emph{After the Baccalauréat, I will...}).
\end{itemize}
\end{exemplebox}

\begin{methode}[Five steps to write a composition at the Probatoire]
\begin{enumerate}
\item \textbf{Brainstorm}: write all your ideas and vocabulary on a draft paper, in any order (5 minutes).
\item \textbf{Plan}: organise the ideas into 3 or 4 paragraphs (introduction, body, conclusion) (5 minutes).
\item \textbf{Draft}: write the composition with linking words; do not stop for small mistakes (15--20 minutes).
\item \textbf{Revise}: check verb agreements (-s!), quantifiers, spelling and punctuation (5 minutes).
\item \textbf{Final copy}: copy neatly on the exam sheet. A clean copy gives a good impression.
\end{enumerate}
\end{methode}

\begin{savaistu}[The composition at the Probatoire technique]
At the Probatoire and Baccalauréat technique, the composition carries an important part of the English mark. Examiners use criteria: content, organisation, grammar and vocabulary. A clear plan with an introduction, a body and a conclusion already guarantees about half of the organisation points --- even before the examiner reads your ideas! Never write without a plan.
\end{savaistu}

\courssub{Guided Writing: ``My Dream Profession''}

\begin{exoresolu}[A model composition]
Write a composition of 120--150 words entitled \emph{``Why I want to become a tailor''}.
\tcblower
\textbf{Solution (model answer, 140 words):}\par
\emph{Many young people in my town dream of office jobs, but my dream profession is tailoring. A tailor measures customers, cuts cloth and sews beautiful clothes with a machine.}\par
\emph{First, I love this trade because I am creative and I like working with my hands. Moreover, a good tailor has many customers, especially before Christmas and weddings. Furthermore, tailoring does not need much money to start: a few tools, a machine and a small shop are enough. However, the work requires patience, because each customer wants a perfect dress or suit.}\par
\emph{To succeed, I need some training. After my Baccalauréat, I will spend two years as an apprentice with Mr Ndi, the best tailor in our quarter. Both of my parents encourage me, and I save a little money every month for my future workshop.}\par
\emph{In conclusion, tailoring is a useful and profitable profession. Therefore, I will work hard to open my own shop one day.}
\par\medskip
\textbf{Notice:} the -s agreements (\emph{a tailor measures, cuts, sews}), the quantifiers (\emph{many, a few, much, each, some, both of, a little}) and the connectors (\emph{first, moreover, furthermore, however, therefore, in conclusion}).
\end{exoresolu}

\courssub{Peer Correction}

\begin{definition}[Peer correction]
\cle{Peer correction} means \emph{exchanging and correcting each other's compositions using a checklist} (correction entre pairs). You read your classmate's text, you tick the criteria on the grid, and you write one positive comment and one piece of advice. It helps both students: the writer receives feedback, and the reader learns to spot mistakes.
\end{definition}

\begin{center}
\begin{tabularx}{\linewidth}{|l|X|c|}
\hline
\rowcolor{popblueL} \textbf{Criterion} & \textbf{Checklist question} & \textbf{Mark /5} \\
\hline
Content & Are there clear ideas and reasons about the profession? & \\
\hline
Organisation & Is there an introduction, a body and a conclusion? Are linking words used? & \\
\hline
Grammar & Are verbs correct (-s at the 3rd person)? Are quantifiers used correctly? & \\
\hline
Vocabulary & Are there at least five words about trades and money? & \\
\hline
\rowcolor{popcream} \textbf{Total} & & \textbf{/20} \\
\hline
\end{tabularx}
\end{center}

\begin{exercice}[Peer correction practice]
Exchange your composition \emph{``My dream profession''} with your neighbour.
\begin{enumerate}
\item Read the text twice. The first time, read for pleasure; the second time, use the grid above.
\item Underline in red every verb without -s at the third person.
\item Circle the linking words. If there are fewer than four, give advice.
\item Give a mark out of 20 and write one positive comment and one suggestion.
\end{enumerate}
\end{exercice}

\begin{corrige}[Exercice 1]
There is no single answer: the correction depends on each composition. A good peer corrector: (a) ticks each criterion honestly; (b) underlines errors such as \emph{``she want''} $\rightarrow$ \emph{``she wants''}; (c) checks the presence of at least four connectors (\emph{first, moreover, however, therefore, in conclusion}); (d) writes a kind, precise comment, for example: \emph{``Good ideas and rich vocabulary, but add a conclusion and check the verb agreements.''}
\end{corrige}

\courssub{Integration Activities, Evaluation and Remediation}

Now let us combine \emph{everything} from this lesson in one complete situation: buying and selling, polite requests, partitive articles, quantifiers, personal finances and professions.

\begin{exercice}[Integration: A day at the market with a budget]
\textbf{Situation:} You are an apprentice tailor. Your master gives you \SI{10000}{FCFA} and sends you to the market to buy supplies. Write a short dialogue (8--10 lines) with the seller, then a short paragraph (5 lines) explaining how you managed the money.
\begin{enumerate}
\item In the dialogue, use at least: two polite requests (\emph{Could you...? / Let me...}), three partitive articles (\emph{a bunch of, a cup of, a piece of...}) and two quantifiers (\emph{a few, some, much, many...}).
\item In the paragraph, say what you bought, how much you spent, and how much you saved. Use \emph{first, moreover, therefore}.
\end{enumerate}
\end{exercice}

\begin{corrige}[Exercice 2 (integration)]
\textbf{Possible dialogue:}\par
--- Good morning, madam. \emph{Could you} show me \emph{a piece of} white cloth, please?\par
--- Of course. This one costs \SI{3000}{FCFA} a metre.\par
--- It is expensive. \emph{Let me} see \emph{a few} other pieces, please.\par
--- Here you are. This one is \SI{2000}{FCFA}.\par
--- Fine. I also need \emph{a bunch of} bananas and \emph{a cup of} rice for lunch. How much is everything?\par
--- \SI{6500}{FCFA}, my dear.\par
--- Thank you very much, madam.\par
\textbf{Possible paragraph:} \emph{First}, I bought two metres of cloth for \SI{4000}{FCFA}. \emph{Moreover}, I bought food for \SI{2500}{FCFA}. I spent \SI{6500}{FCFA} in total; \emph{therefore}, I saved \SI{3500}{FCFA} for my master. I did not spend \emph{much} money because I compared the prices.
\end{corrige}

\begin{exercice}[Evaluation: final composition]
Write a composition of 120--150 words: \emph{``Why I want to become a \ldots''} (choose: mechanic, nurse, electrician, farmer, teacher, tailor, computer technician). Follow the five-step method (brainstorm, plan, draft, revise, final copy). Your teacher will mark it with the grid: content /5, organisation /5, grammar /5, vocabulary /5.
\end{exercice}

\begin{corrige}[Exercice 3 (evaluation)]
Marking guide: \textbf{Content} (clear profession, real reasons, personal project). \textbf{Organisation} (introduction, 2--3 body paragraphs, conclusion; at least four linking words). \textbf{Grammar} (correct -s agreements; correct quantifiers: \emph{much} with uncountables, \emph{many} with countables, \emph{a few} $=$ some, \emph{few} $=$ almost none). \textbf{Vocabulary} (at least five lesson words: \emph{apprentice, salary, customer, training, workshop, budget, profit, skill}). Remove one point per category for repeated errors.
\end{corrige}

\begin{methode}[Remediation: if your mark is low]
\begin{enumerate}
\item \textbf{Grammar problem?} Rewrite five sentences about professions with correct -s agreements (\emph{A nurse helps\ldots}).
\item \textbf{Vocabulary problem?} Make a personal word list of 15 trade words with French glosses and learn five per day.
\item \textbf{Organisation problem?} Copy the model composition and label each paragraph: introduction, body 1, body 2, conclusion.
\item \textbf{Then rewrite} the same composition and ask a classmate for a new peer correction. Compare your two marks.
\end{enumerate}
\end{methode}

\begin{aretenir}[What to remember]
\begin{itemize}
\item A composition = \textbf{introduction} (present the job) + \textbf{body} (reasons, qualities, training) + \textbf{conclusion} (personal project).
\item Use linking words: \emph{first, moreover, furthermore, however, therefore, in conclusion}.
\item Show the lesson: quantifiers (\emph{many, a few, much, each, both of}), partitive articles (\emph{a bunch of, a cup of}) and trade vocabulary.
\item Never forget the \textbf{-s} at the third person singular: \emph{A carpenter works}.
\item Method: brainstorm $\rightarrow$ plan $\rightarrow$ draft $\rightarrow$ revise $\rightarrow$ final copy.
\item Peer correction with a checklist improves both the writer and the reader.
\end{itemize}
\end{aretenir}

\newpage

\coursec{Activité d'intégration}

\coursec{Economic Life and Occupations}

\courssub{Objectifs de l'activité}
\begin{itemize}
    \item \textbf{Mobiliser} l'ensemble des savoirs de la leçon \emph{Economic Life and Occupations} dans une situation unique et authentique : le marché Mokolo à Yaoundé.
    \item \textbf{Produire} un dialogue oral d'achat/vente intégrant les \cle{polite requests} (\emph{requêtes polies}) et les \cle{partitive articles} (\emph{articles partitifs}).
    \item \textbf{Analyser} un budget familial réel en utilisant les \cle{determiners and quantifiers} (\emph{déterminants et quantificateurs}) et réagir avec \cle{So do I / Neither do I}.
    \item \textbf{Comprendre} un texte sur l'orientation professionnelle (\emph{apprentissage d'un métier}) et en extraire le vocabulaire clé.
    \item \textbf{Rédiger} une composition personnelle argumentée : \emph{The trade I would like to learn and why}.
\end{itemize}

\courssub{Situation}
\begin{aretenir}[Contexte : Une journée au Marché Mokolo]
Imaginez que vous êtes un groupe d'élèves de Première Technique en stage d'observation à Yaoundé. Votre tuteur vous donne trois missions pour la journée. Vous disposez d'un dossier contenant trois documents authentiques.
\end{aretenir}

\textbf{Document 1 : Liste de courses et budget (10\,000 FCFA)}
Vous devez acheter les articles suivants pour le déjeuner de l'équipe. Les prix sont ceux affichés chez \emph{Maman Simone}, une vendeuse bien connue au marché Mokolo. Le total exact est de 10\,000 FCFA après négociation sur les tomates.

\begin{center}
\begin{tabularx}{\linewidth}{|l|X|c|c|}\hline
\rowcolor{popblueL}
\textbf{Article} & \textbf{Description / Unités (Partitives)} & \textbf{Prix unitaire (FCFA)} & \textbf{Qté} \\ \hline
Riz & \emph{a bag of} 5\,kg & 2\,500 & 1 \\ \hline
Tomates & \emph{a heap of} (tas) & 500 & 3 \\ \hline
Poisson (bar) & \emph{a piece of} & 1\,200 & 2 \\ \hline
Huile & \emph{a bottle of} 1\,L & 1\,000 & 1 \\ \hline
Oignons & \emph{a net of} (filet) 2\,kg & 800 & 1 \\ \hline
Piment & \emph{a bunch of} (botte) & 200 & 2 \\ \hline
Pain & \emph{a loaf of} & 300 & 2 \\ \hline
Boisson (jus) & \emph{a carton of} 6 bouteilles & 1\,000 & 1 \\ \hline
\end{tabularx}
\end{center}
\textbf{Total après négociation (3 tas de tomates à 1\,300 FCFA) :} 10\,000 FCFA exactement.

\textbf{Document 2 : Relevé des dépenses mensuelles de la famille Ngono (Yaoundé, 4 personnes)}
\begin{center}
\begin{tabular}{|l|r|c|}\hline
\rowcolor{popblueL}
\textbf{Poste de dépense} & \textbf{Montant (FCFA)} & \textbf{\% du total} \\ \hline
Loyer & 150\,000 & 30\% \\ \hline
Alimentation & 200\,000 & 40\% \\ \hline
Transport & 50\,000 & 10\% \\ \hline
Éducation (frais scolaires) & 60\,000 & 12\% \\ \hline
Santé & 20\,000 & 4\% \\ \hline
Épargne & 20\,000 & 4\% \\ \hline
\textbf{Total} & \textbf{500\,000} & \textbf{100\%} \\ \hline
\end{tabular}
\end{center}

\begin{popfigure}
\begin{tabular}{|l|r|}\hline
Rent (Loyer) & 30\,\% \\ \hline Food (Alimentation) & 40\,\% \\ \hline Transport & 10\,\% \\ \hline Education & 12\,\% \\ \hline Health (Santé) & 4\,\% \\ \hline Savings (Épargne) & 4\,\% \\ \hline
\end{tabular}
\legende{Figure 1 -- Répartition du budget mensuel de la famille Ngono (Total : 500\,000 FCFA). Source : Document 2.}
\end{popfigure}

\textbf{Document 3 : Témoignage — \emph{« From Apprentice to Electrician »}}
\begin{exemplebox}[Texte de compréhension]
My name is Paul. I am 19 years old and I live in Douala. Two years ago, I failed my Probatoire. I was very sad. My uncle, who is a master electrician, told me: \emph{``Let me teach you the trade. It is better than staying at home.''} So, I became his apprentice.

At first, it was difficult. I had to wake up at 5:30 a.m. every day. I carried heavy tools and crawled under houses to fix wires. \emph{Few} of my friends understood my choice; \emph{most} of them wanted office jobs. But I liked the work. \emph{Each} day, I learned something new: how to read a plan, how to test a circuit, how to respect safety rules.

Today, I earn a small salary, but I save \emph{a little} money every month. I want to pass my CAP (Certificat d'Aptitude Professionnelle) next year. \emph{Both} my uncle and my father are proud of me. \emph{None} of us regret this decision.

\textbf{Key vocabulary (Glossaire)} : \cle{apprentice} (apprenti), \cle{master electrician} (électricien qualifié/maître d'apprentissage), \cle{crawl} (ramper), \cle{wires} (fils électriques), \cle{circuit} (circuit), \cle{safety rules} (règles de sécurité), \cle{CAP} (Certificat d'Aptitude Professionnelle), \cle{earn} (gagner), \cle{save} (épargner).
\end{exemplebox}

\begin{savaistu}[L'apprentissage et le CAP au Cameroun]
Au Cameroun, la formation professionnelle est pilotée par le \textbf{MINEFOP} (Ministère de l'Emploi et de la Formation Professionnelle). Le \cle{CAP} (Certificat d'Aptitude Professionnelle) s'obtient après 2 à 3 ans d'apprentissage chez un \cle{maître artisan} (maître d'apprentissage) validé, complété par des cours théoriques dans un \cle{CFA} (Centre de Formation d'Apprentis) ou un \cle{Lycée Technique}. L'examen est organisé par l'\cle{OBC} (Office du Baccalauréat du Cameroun). C'est une voie d'excellence pour l'insertion professionnelle rapide, très valorisée dans les secteurs du BTP, de l'industrie, des services et du numérique.
\end{savaistu}

\courssub{Consignes}
Travaillez en groupes de 4 élèves. Chaque groupe rend un rapport écrit et joue une saynète.

\begin{enumerate}
    \item \textbf{Tâche 1 — Role-play (Speaking / Grammar / Vocabulary) :} Rédigez et jouez un dialogue de 2 minutes entre \emph{l'acheteur (vous)} et \emph{Maman Simone (vendeuse)}. Utilisez \textbf{au moins} :
    \begin{itemize}
        \item 2 requêtes polies différentes (\emph{Let me...}, \emph{Could you...?}, \emph{Can you...?}).
        \item 4 articles partitifs différents (\emph{a bag of}, \emph{a heap of}, \emph{a piece of}, \emph{a bottle of}, \emph{a net of}, \emph{a bunch of}, \emph{a loaf of}, \emph{a carton of}).
        \item Une négociation de prix pour un article (ex: \emph{``Could you make it 400 FCFA for the heap?''}).
    \end{itemize}
    
    \item \textbf{Tâche 2 — Analyse budgétaire (Grammar / Speaking) :} À partir du Document 2, écrivez 5 phrases d'analyse en utilisant \textbf{chaque type de quantifieur} vu en cours (Part 1, 2, 3). Exemple : \emph{Most of the money goes to food.} Ensuite, en groupe, réagissez à chaque phrase avec \textbf{So do I} ou \textbf{Neither do we} selon votre avis personnel.
    
    \item \textbf{Tâche 3 — Compréhension de texte (Reading / Vocabulary) :} Répondez par écrit aux questions suivantes sur le Document 3 :
    \begin{enumerate}
        \item Why did Paul become an apprentice? (2 reasons)
        \item Find in the text: one example of \emph{few}, \emph{most}, \emph{each}, \emph{a little}, \emph{both}, \emph{none}. Explain the meaning in context.
        \item What are Paul's future plans?
        \item Give the English definitions for: \emph{apprentice, crawl, wires, safety rules}.
    \end{enumerate}
    
    \item \textbf{Tâche 4 — Production écrite individuelle (Writing) :} Rédigez une composition de 150--180 mots : \textbf{``The trade I would like to learn and why''}. Structurez votre texte :
    \begin{itemize}
        \item Introduction : Présentez le métier choisi (électricien, menuisier, plombier, informaticien, coiffeur, etc.).
        \item Corps : Expliquez \emph{why} (intérêt personnel, demande sur le marché, revenu, fierté). Utilisez \emph{determiners/quantifiers} (\emph{each step, many opportunities, a little patience, both theory and practice}).
        \item Conclusion : Projet concret (CAP, stage, création d'entreprise).
    \end{itemize}
\end{enumerate}

\courssub{Ressources à mobiliser}
\begin{definition}[Polite Requests -- Requêtes polies]
\begin{itemize}
    \item \textbf{Let me + bare infinitive} : \emph{Let me help you with that bag.} (Proposition d'aide / offre de service).
    \item \textbf{Could you + bare infinitive?} : \emph{Could you give me a better price?} (Politesse standard, demande courtoise).
    \item \textbf{Can you + bare infinitive?} : \emph{Can you weigh the tomatoes?} (Moins formel, familier, demande de service simple).
\end{itemize}
\end{definition}

\begin{attention}[Règle d'or]
Après \emph{Let me}, \emph{Could you}, \emph{Can you}, le verbe est \textbf{toujours à l'infinitif sans \emph{to}} (bare infinitive). \emph{Never say: ``Let me to help'' or ``Could you to give?''}
\end{attention}

\begin{definition}[Partitive Articles -- Articles partitifs (Quantités / Conteneurs)]
Structure : \textbf{a/an + noun of quantity + of + noun}.
Utilisés pour dénombrer des noms \emph{indénombrables} (uncountable) ou des ensembles.
\begin{center}
\begin{tabular}{|l|l|}\hline
\rowcolor{popblueL}
\textbf{Contenant / Mesure} & \textbf{Exemple (Document 1)} \\ \hline
a bag of & a bag of rice \\ \hline
a heap of & a heap of tomatoes \\ \hline
a piece of & a piece of fish \\ \hline
a bottle of & a bottle of oil \\ \hline
a net of & a net of onions \\ \hline
a bunch of & a bunch of pepper \\ \hline
a loaf of & a loaf of bread \\ \hline
a carton of & a carton of juice \\ \hline
\end{tabular}
\end{center}
\end{definition}

\begin{attention}[Accord au pluriel]
Le nom de quantité (container) prend la marque du pluriel, pas le contenu : \emph{three \textbf{heaps} of tomatoes}, \emph{two \textbf{pieces} of fish}, \emph{two \textbf{loaves} of bread}.
\end{attention}

\begin{propriete}[Determiners and Quantifiers -- Résumé des 3 parties (Programme MINESEC)]
\begin{itemize}
    \item \textbf{Part 1 -- Quantité générale (Affirmation / Négation / Question) :}
    \cle{some} (affirmatif, dén. pl. \& indén.), \cle{any} (négatif / interrogatif), \cle{much} (indénombrable, surtout ? / -), \cle{all (of)} (totalité), \cle{most (of)} (majorité).
    \item \textbf{Part 2 -- Quantité faible (Nuance positive / négative) :}
    \cle{few} / \cle{a few} (dénombrables pluriel) : \emph{few} = presque aucun (négatif) ; \emph{a few} = quelques-uns (positif). \\
    \cle{little} / \cle{a little} (indénombrables singulier) : \emph{little} = presque rien (négatif) ; \emph{a little} = un peu (positif).
    \item \textbf{Part 3 -- Distribution, Dualité, Négation totale :}
    \cle{each} (chacun, focus individuel), \cle{every} (tous, focus collectif), \cle{many} (dénombrables, ? / +), \cle{how many} (question), \cle{both (of)} (les deux), \cle{no} (aucun, + nom), \cle{none (of)} (aucun d'entre eux, pronom).
\end{itemize}
\end{propriete}

\begin{definition}[So / Neither -- Accord / Désaccord (Short Answers)]
\begin{itemize}
    \item \textbf{So + auxiliaire + sujet} : Accord avec une phrase \textbf{affirmative}. \emph{``I like this job.'' -- ``So do I.''}
    \item \textbf{Neither + auxiliaire + sujet} : Accord avec une phrase \textbf{négative}. \emph{``I don't like office work.'' -- ``Neither do I.''}
    \item L'auxiliaire \textbf{doit correspondre} au temps et au verbe de la phrase originale (\emph{do/does/did/will/can/be/have}).
\end{itemize}
\end{definition}

\begin{exemplebox}[Exemples variés]
\begin{tabular}{ll}
\emph{``I am tired.''} & $\rightarrow$ \emph{``So am I.''} (auxiliaire \emph{be}) \\
\emph{``I have finished.''} & $\rightarrow$ \emph{``So have I.''} (auxiliaire \emph{have}) \\
\emph{``I can drive.''} & $\rightarrow$ \emph{``So can I.''} (modal \emph{can}) \\
\emph{``I didn't go.''} & $\rightarrow$ \emph{``Neither did I.''} (passé \emph{did}) \\
\end{tabular}
\end{exemplebox}

\courssub{Démarche et corrigé commenté}
\textbf{Corrigé guidé de l'activité d'intégration}\par
\textbf{Étape 1 -- Rédaction et jeu du dialogue (Tâche 1)}
\begin{exemplebox}[Dialogue modèle : At Maman Simone's Stall]
\textbf{Buyer:} Good morning, Maman Simone. \textbf{Let me} see your tomatoes. They look fresh.
\textbf{Seller:} Good morning, my son. Yes, very fresh today. \textbf{A heap} is 500 FCFA.
\textbf{Buyer:} \textbf{Could you} give me three heaps? \textbf{Can you} make it 1\,300 FCFA for the three?
\textbf{Seller:} Hmm, okay, 1\,300 FCFA for three heaps. Anything else?
\textbf{Buyer:} Yes. I need \textbf{a bag of} rice (5kg), \textbf{two pieces of} fish, \textbf{a bottle of} oil, \textbf{a net of} onions, \textbf{two bunches of} pepper, \textbf{two loaves of} bread and \textbf{a carton of} juice.
\textbf{Seller:} That makes exactly 10\,000 FCFA with the discount on tomatoes. \textbf{Let me} pack everything for you.
\textbf{Buyer:} Thank you very much. Here is the money.
\textbf{Seller:} Thank you. Good bye!
\end{exemplebox}
\begin{attention}[Points de vigilance pour l'évaluation]
\begin{itemize}
    \item Vérifier l'usage de \emph{bare infinitive} après \emph{Let me / Could you / Can you} (pas de \emph{to}).
    \item Vérifier l'accord du partitif : \emph{a heap} (singulier) vs \emph{three heaps} (pluriel). \emph{Two pieces of fish} (pas \emph{fishes}).
    \item Prononciation de \emph{Could you} /kʊdʒuː/ et \emph{Let me} /lemmiː/.
    \item Intonation montante pour les questions (\emph{Could you...?}), descendante pour les offres (\emph{Let me...}).
\end{itemize}
\end{attention}

\textbf{Étape 2 -- Analyse du budget familial (Tâche 2)}
\textit{Exemples de phrases produites par les élèves (une par type de quantifieur, couvrant les 3 parties) :}
\begin{enumerate}
    \item \textbf{Most of the money} goes to food and rent (70\%). \rightarrow \textit{Réaction :} \textbf{So do we.} (C'est notre cas aussi).
    \item \textbf{Some of the income} is saved, but \textbf{not much}. \rightarrow \textit{Réaction :} \textbf{Neither do we.} (Nous n'épargnons pas beaucoup non plus).
    \item \textbf{Few families} can save 10\% in Yaoundé today. \rightarrow \textit{Réaction :} \textbf{So do I.} (Je pense la même chose).
    \item \textbf{A little} money is left for leisure after expenses. \rightarrow \textit{Réaction :} \textbf{So do we.}
    \item \textbf{Each} family member contributes to the chores. \rightarrow \textit{Réaction :} \textbf{So do we.}
    \item \textbf{Both} parents work to pay the fees. \rightarrow \textit{Réaction :} \textbf{So do we.}
    \item \textbf{None of the children} pay rent. \rightarrow \textit{Réaction :} \textbf{Neither do ours.}
\end{enumerate}
\begin{methode}[Comment choisir le quantifieur ?]
\begin{enumerate}
    \item Identifier le nom : dénombrable (pluriel : \emph{friends, families}) ou indénombrable (singulier : \emph{money, time}) ?
    \item Identifier le sens visé : quantité grande (\emph{many, much, most, all}), petite positive (\emph{a few, a little}), petite négative (\emph{few, little}), nulle (\emph{no, none}), distribution (\emph{each, every}), dualité (\emph{both}).
    \item Vérifier la structure syntaxique :
    \begin{itemize}
        \item \emph{Quantifier + of + déterminant + nom} : \emph{most of the money, both of the parents, none of the children}.
        \item \emph{Quantifier + nom} (sans \emph{of}) : \emph{much money, many friends, each day, every student, no money}.
    \end{itemize}
\end{enumerate}
\end{methode}

\textbf{Étape 3 -- Compréhension du texte (Tâche 3)}
\begin{enumerate}
    \item \textbf{Why apprentice?} (1) He failed Probatoire. (2) His uncle proposed to teach him the trade (\emph{Let me teach you}).
    \item \textbf{Quantifieurs dans le texte :}
    \begin{itemize}
        \item \emph{Few of my friends understood...} = \emph{Peu de mes amis} (presque aucun, connotation négative / déception).
        \item \emph{Most of them wanted office jobs.} = \emph{La plupart d'entre eux} (la majorité, > 50\%).
        \item \emph{Each day, I learned...} = \emph{Chaque jour} (individuellement, jour après jour, insistance sur la régularité).
        \item \emph{I save a little money...} = \emph{Un peu d'argent} (petite quantité, connotation positive : il y en a, ça compte).
        \item \emph{Both my uncle and my father...} = \emph{Les deux} (mon oncle ET mon père, dualité).
        \item \emph{None of us regret...} = \emph{Aucun de nous} (négation totale du groupe, pronom).
    \end{itemize}
    \item \textbf{Future plans:} Pass CAP next year, become a certified electrician, eventually own a workshop.
    \item \textbf{Définitions :} \emph{Apprentice} = someone learning a trade from a skilled employer. \emph{Crawl} = move on hands and knees. \emph{Wires} = metal threads conducting electricity. \emph{Safety rules} = regulations to prevent accidents.
\end{enumerate}

\textbf{Étape 4 -- Composition type (Tâche 4)}
\begin{exemplebox}[Modèle de composition : The trade I would like to learn and why]
\textbf{The Trade I Would Like to Learn and Why}

I would like to become a \textbf{computer maintenance technician}. In our modern world, \textbf{every} office, school and shop uses computers. \textbf{Many} young people in Cameroon have a phone or a laptop, but \textbf{few} know how to repair them when they break down. This is a real opportunity.

\textbf{Each} day, technology changes. I like solving problems. When a screen goes black or a virus attacks a system, I want to be the one who finds the solution. It requires \textbf{a little} patience and \textbf{both} theoretical knowledge (software, operating systems) and practical skills (changing a motherboard, cleaning a fan).

\textbf{Most of} my friends want to work in an office, but I prefer a technical trade. \textbf{None of} them can fix a hardware issue. After my Baccalauréat Technique, I plan to do a \textbf{BTS} in Computer Maintenance or a specialized \textbf{CAP}. I want to open my own small workshop in my neighbourhood. It is a trade that gives dignity, a steady income, and the chance to help \textbf{all} members of my community stay connected.
\end{exemplebox}
\begin{aretenir}[Grille d'évaluation de la composition (sur 20)]
\begin{itemize}
    \item Contenu / Idées (pertinence, arguments, exemples) : 5 pts
    \item Organisation (paragraphes, connecteurs logiques : \emph{however, furthermore, therefore}) : 4 pts
    \item Grammaire (Quantifieurs variés, temps verbaux, structures polies) : 5 pts
    \item Vocabulaire (métier, outils, verbes d'action, adjectifs qualitatifs) : 3 pts
    \item Orthographe / Ponctuation / Lisibilité : 3 pts
\end{itemize}
\end{aretenir}

\courssub{Bilan de l'activité}
\begin{aretenir}[Ce que l'activité a permis de mettre en évidence]
\begin{itemize}
    \item \textbf{Transfert des acquis :} Les élèves ont dû passer de l'exercice de grammaire isolé (remplir des trous) à l'usage fonctionnel : \emph{acheter, analyser, argumenter, écrire}.
    \item \textbf{Contextualisation camerounaise :} L'utilisation du FCFA, du marché Mokolo, du système éducatif (Probatoire, CAP, BTS, MINEFOP) et des noms locaux (Maman Simone, famille Ngono, Douala/Yaoundé) ancre l'anglais dans la réalité vécue, augmentant la motivation.
    \item \textbf{Intégration des 4 compétences :} \emph{Speaking} (dialogue), \emph{Listening} (écoute des autres groupes), \emph{Reading} (texte Paul), \emph{Writing} (analyse + composition).
    \item \textbf{Difficultés récurrentes observées (pour remédiation) :}
    \begin{itemize}
        \item Confusion \emph{few / a few} et \emph{little / a little} (connotation négative vs positive).
        \item Oubli du \emph{of} après \emph{all / most / both / none} quand il y a un déterminant (\emph{most of the money} vs \emph{most people}).
        \item Mauvaise formation du \emph{bare infinitive} après \emph{Let me / Could you} (ajout de \emph{to}).
        \item Difficulté à choisir l'auxiliaire pour \emph{So / Neither} (ex: \emph{So am I} au lieu de \emph{So do I} avec verbe d'action simple).
        \item Pluriel des partitifs : \emph{two loaves} (irrégulier), \emph{two pieces} (pas \emph{fishes}).
    \end{itemize}
    \item \textbf{Pistes de remédiation immédiate :}
    \begin{itemize}
        \item Fiche exercices ciblée : \emph{Quantifiers Mini-Test} (10 phrases à trous mélangés Part 1, 2, 3).
        \item Atelier oral : \emph{Market Price Negotiation} (cartes prix + cartes budget, rotation rapide par paires).
        \item Réécriture guidée de la composition pour les notes < 10/20, avec cadre de phrases (\emph{skeleton essay} : \emph{I would like to be a... Because... Each day... Both... Finally...}).
    \end{itemize}
\end{itemize}
Cette activité de synthèse valide la compétence \emph{« Communiquer dans des situations de la vie économique »} et prépare efficacement les épreuves du Probatoire (Expression orale, Compréhension, Expression écrite).\end{aretenir}


\newpage

\coursec{Méthodes et exercices résolus}

\coursec{Economic Life and Occupations}

\courssub{Buying and Selling at the Market (Speaking \& Vocabulary)}

\begin{methode}[Réaliser un dialogue d'achat/vente au marché avec politesse et articles partitifs]
    \textbf{Objectif :} Échanger des informations pour acheter ou vendre des articles/denrées alimentaires en utilisant des formules de politesse et des quantificateurs partitifs.
    \begin{enumerate}
        \item \textbf{Identifier les interlocuteurs et le contexte :} Vendeur vs Acheteur, marché vs boutique. Adapter le registre (formel/poli).
        \item \textbf{Utiliser les formules de requêtes polies (Polite Requests) :}
              \begin{itemize}
                  \item \cle{Let me + bare infinitive} (Proposition d'aide : \emph{Let me help you with that bag.})
                  \item \cle{Could you / Can you + bare infinitive + please ?} (Demande : \emph{Could you give me a kilo of rice, please?})
                  \item \cle{Would you like + noun / to + infinitive ?} (Offre : \emph{Would you like some tomatoes?})
              \end{itemize}
        \item \textbf{Employer les articles partitifs (Partitive Articles) pour exprimer des quantités indéfinies :}
              \begin{center}
                  \begin{tabular}{|l|l|l|}\hline
                  \textbf{Structure} & \textbf{Exemple} & \textbf{Traduction} \\ \hline
                  a \cle{bunch of} + nom pluriel & a bunch of bananas & une grappe/régime de bananes \\
                  a \cle{cup of} + nom non comptable & a cup of tea & une tasse de thé \\
                  a \cle{loaf of} + pain & a loaf of bread & une miche de pain \\
                  a \cle{kilo of} / \cle{packet of} + nom & a kilo of sugar / a packet of biscuits & un kilo de sucre / un paquet de biscuits \\
                  a \cle{piece of} / \cle{slice of} & a piece of meat / a slice of cake & un morceau de viande / une part de gâteau \\ \hline
                  \end{tabular}
              \end{center}
        \item \textbf{Structurer le dialogue :} Salutation $\rightarrow$ Demande/Offre (quantité + article) $\rightarrow$ Négociation prix (optionnel) $\rightarrow$ Paiement $\rightarrow$ Remerciements/Au revoir.
        \item \textbf{Vérifier :} Accord nombre (a bunch \emph{of bananas}), politesse (\emph{please}, \emph{could}), vocabulaire spécifique (marché : \emph{stall, price, change, receipt}).
    \end{enumerate}
\end{methode}

\begin{exoresolu}[Dialogue au marché : Complétion et réécriture]
    \textbf{Énoncé :} Complétez le dialogue ci-dessous entre un vendeur (V) et un acheteur (A) au marché Mfoundi. Utilisez les mots entre parenthèses sous la forme correcte (articles partitifs, forme polie). Traduisez les phrases en français.
    \begin{enumerate}
        \item V: Good morning, Madam. \_\_\_\_\_\_\_\_\_\_\_\_\_\_\_\_ (help / you) ?
        \item A: Good morning. Yes, \_\_\_\_\_\_\_\_\_\_\_\_\_\_\_\_ (like / buy) \_\_\_\_\_\_\_\_\_\_\_\_\_\_\_\_ (a bunch / banana) and \_\_\_\_\_\_\_\_\_\_\_\_\_\_\_\_ (a kilo / tomato).
        \item V: Here you are. \_\_\_\_\_\_\_\_\_\_\_\_\_\_\_\_ (anything / else) ?
        \item A: \_\_\_\_\_\_\_\_\_\_\_\_\_\_\_\_ (let / me / see) ... Oh, \_\_\_\_\_\_\_\_\_\_\_\_\_\_\_\_ (could / you / give / me) \_\_\_\_\_\_\_\_\_\_\_\_\_\_\_\_ (a packet / sugar), please?
        \item V: Certainly. That will be 3,500 FCFA.
        \item A: Here is 5,000 FCFA. \_\_\_\_\_\_\_\_\_\_\_\_\_\_\_\_ (thank / you).
        \item V: Here is your change. Have a nice day!
    \end{enumerate}
    \tcblower
    \textbf{Solution détaillée :}
    \begin{enumerate}
        \item \textbf{V: Good morning, Madam. \cle{Can I help you} ? / \cle{Could I help you} ?} \\
              \textit{Raisonnement :} Offre de service standard. "Can I" ou "Could I" + bare infinitive. "Let me help you" est possible mais plus direct (souvent après avoir vu un besoin). \\
              \textit{Français :} Bonjour Madame. Puis-je vous aider ?
        \item \textbf{A: Good morning. Yes, \cle{I would like to buy} \cle{a bunch of bananas} and \cle{a kilo of tomatoes}.} \\
              \textit{Raisonnement :} "Would like to" pour exprimer un désir poli. "A bunch of" + pluriel (bananas). "A kilo of" + pluriel (tomatoes). Attention : \emph{tomato} $\rightarrow$ \emph{tomatoes} (pluriel après kilo of). \\
              \textit{Français :} Bonjour. Oui, je voudrais acheter une grappe de bananes et un kilo de tomates.
        \item \textbf{V: Here you are. \cle{Would you like anything else} ? / \cle{Is there anything else} ?} \\
              \textit{Raisonnement :} Proposition polie pour continuer l'achat. "Anything else" est l'expression consacrée. \\
              \textit{Français :} Voici. Désirez-vous autre chose ?
        \item \textbf{A: \cle{Let me see} ... Oh, \cle{could you give me} \cle{a packet of sugar}, please?} \\
              \textit{Raisonnement :} "Let me see" (réflexion, bare infinitive sans "to"). Demande polie : "Could you give me" + objet. Article partitif : "a packet of" + nom non comptable (sugar). \\
              \textit{Français :} Laissez-moi voir... Oh, pourriez-vous me donner un paquet de sucre, s'il vous plaît ?
        \item \textbf{V: Certainly. That will be 3,500 FCFA.} (Phrase fournie).
        \item \textbf{A: Here is 5,000 FCFA. \cle{Thank you} / \cle{Thanks}.} \\
              \textit{Raisonnement :} Remerciement standard. \\
              \textit{Français :} Voici 5000 FCFA. Merci.
        \item \textbf{V: Here is your change. Have a nice day!} (Phrase fournie).
    \end{enumerate}
    \textbf{Conclusion :} Le dialogue respecte les conventions de politesse (\emph{could, would like, please}) et utilise correctement les articles partitifs (\emph{a bunch of, a kilo of, a packet of}) pour quantifier des denrées non comptables ou en vrac.
\end{exoresolu}

\courssub{Polite Requests and Partitive Articles (Grammar)}

\begin{methode}[Utiliser "So do I / Neither did he" et les adjectifs comme noms]
    \textbf{Objectif :} Exprimer l'accord ou le désaccord avec une affirmation (positive ou négative) et utiliser des adjectifs comme noms pour désigner des groupes de personnes.
    \begin{enumerate}
        \item \textbf{Structure "So + Auxiliaire + Sujet" (Accord positif) :}
              \begin{itemize}
                  \item Reprend l'auxiliaire de la phrase originale (do/does/did/is/are/was/were/have/has/can/will/must...).
                  \item Ex: \emph{I like finance.} $\rightarrow$ \cle{So do I.} (Moi aussi).
                  \item Ex: \emph{She saved money.} $\rightarrow$ \cle{So did she.}
              \end{itemize}
        \item \textbf{Structure "Neither / Nor + Auxiliaire + Sujet" (Accord négatif) :}
              \begin{itemize}
                  \item Phrase originale négative (not, never, no...). Auxiliaire affirmatif dans la réponse.
                  \item Ex: \emph{I don't like debt.} $\rightarrow$ \cle{Neither do I.} (Moi non plus).
                  \item Ex: \emph{He didn't spend much.} $\rightarrow$ \cle{Neither did he.}
              \end{itemize}
        \item \textbf{Adjectifs comme noms (The + Adjectif = Groupe de personnes) :}
              \begin{itemize}
                  \item Toujours précédés de \cle{the}. Verbe au pluriel.
                  \item \cle{The rich} (les riches), \cle{the poor} (les pauvres), \cle{the unemployed} (les chômeurs), \cle{the young} (les jeunes), \cle{the elderly} (les personnes âgées).
                  \item Ex: \emph{\cle{The rich} should help \cle{the poor}.}
              \end{itemize}
        \item \textbf{Piège :} Ne pas confondre "Neither" (accord négatif) et "Either" (ajout négatif en fin de phrase : \emph{I don't like it either}). Au Bac, "Neither/Nor" déclenche l'inversion sujet-auxiliaire.
    \end{enumerate}
\end{methode}

\begin{exoresolu}[Finances personnelles : Accord et groupes sociaux]
    \textbf{Énoncé :} Lisez les phrases sur la gestion financière. Réécrivez-les en utilisant la structure indiquée entre parenthèses. Mettez le verbe à la forme correcte.
    \begin{enumerate}
        \item "I track my expenses every week." $\rightarrow$ My brother \_\_\_\_\_\_\_\_\_\_\_\_\_\_\_\_. (So)
        \item "She doesn't have a credit card." $\rightarrow$ Her husband \_\_\_\_\_\_\_\_\_\_\_\_\_\_\_\_. (Neither)
        \item "They saved 50,000 FCFA last year." $\rightarrow$ We \_\_\_\_\_\_\_\_\_\_\_\_\_\_\_\_. (So)
        \item "He cannot invest in stocks now." $\rightarrow$ I \_\_\_\_\_\_\_\_\_\_\_\_\_\_\_\_. (Neither)
        \item Complete: \_\_\_\_\_\_\_\_\_\_\_\_\_\_\_\_ (rich) often invest in real estate, while \_\_\_\_\_\_\_\_\_\_\_\_\_\_\_\_ (poor) struggle to save.
        \item "The government helps \_\_\_\_\_\_\_\_\_\_\_\_\_\_\_\_ (unemployed) with training programs."
    \end{enumerate}
    \tcblower
    \textbf{Solution détaillée :}
    \begin{enumerate}
        \item \textbf{My brother \cle{so does he}.} \\
              \textit{Analyse :} Phrase affirmative, présent simple (\emph{track} = \emph{do track}). Auxiliaire \emph{do}. Sujet \emph{he}. Inversion : \emph{So does he}. (La structure canonique "So + aux + sujet" est attendue).
        \item \textbf{Her husband \cle{neither does he}.} \\
              \textit{Analyse :} Phrase négative (\emph{doesn't have}). Accord négatif $\rightarrow$ \emph{Neither} + auxiliaire affirmatif \emph{does} + sujet \emph{he}. \emph{Neither does he}.
        \item \textbf{We \cle{so did we}.} \\
              \textit{Analyse :} Passé simple (\emph{saved} = \emph{did save}). Auxiliaire \emph{did}. Sujet \emph{we}. \emph{So did we}.
        \item \textbf{I \cle{neither can I}.} \\
              \textit{Analyse :} Modale \emph{cannot} (négatif). Accord négatif $\rightarrow$ \emph{Neither} + modale \emph{can} + sujet \emph{I}. \emph{Neither can I}.
        \item \textbf{\cle{The rich} often invest in real estate, while \cle{the poor} struggle to save.} \\
              \textit{Analyse :} Adjectifs utilisés comme noms $\rightarrow$ \emph{The} + adjectif. Verbe au pluriel (\emph{invest}, \emph{struggle}).
        \item \textbf{The government helps \cle{the unemployed} with training programs.} \\
              \textit{Analyse :} \emph{Unemployed} (chômeurs) $\rightarrow$ \emph{the unemployed}. Groupe de personnes.
    \end{enumerate}
    \textbf{Conclusion :} Maîtriser l'auxiliaire (temps, mode) est la clé pour "So/Neither". Pour les adjectifs-noms, ne jamais oublier \textbf{the} et le verbe au \textbf{pluriel}.
\end{exoresolu}

\courssub{Managing Personal Finances (Listening \& Grammar)}

\begin{definition}[Vocabulaire clé : Gestion des finances personnelles]
    \begin{itemize}
        \item \cle{Income} (revenus) / \cle{Expenses} (dépenses) / \cle{Budget} (budget)
        \item \cle{Savings} (épargne) / \cle{Investment} (investissement) / \cle{Loan} (prêt/emprunt)
        \item \cle{Interest rate} (taux d'intérêt) / \cle{Collateral} (garantie/hypothèque)
        \item \cle{Track expenses} (suivre ses dépenses) / \cle{Make ends meet} (joindre les deux bouts)
        \item \cle{Financial literacy} (éducation financière) / \cle{Emergency fund} (fonds d'urgence)
    \end{itemize}
\end{definition}

\courssub{Determiners and Quantifiers — Parts 1, 2 \& 3 (Grammar)}

\begin{methode}[Maîtriser les déterminants et quantificateurs (Parties 1, 2 et 3)]
    \textbf{Objectif :} Choisir le bon quantificateur selon la nature du nom (comptable/non comptable), la polarité de la phrase (affirmative/négative/interrogative) et le sens (quantité grande/petite/nulle/totale).
    \begin{enumerate}
        \item \textbf{Partie 1 : \cle{some} / \cle{any} / \cle{much} / \cle{all} / \cle{most}}
              \begin{center}
                  \begin{tabular}{|c|c|c|c|}\hline
                  \textbf{Quantificateur} & \textbf{Nom} & \textbf{Contexte} & \textbf{Exemple} \\ \hline
                  \cle{Some} / \cle{Some of} & Comptable pluriel / Non comptable & Affirmatif ; Offres/Demandes polies (interrogatif) & \emph{I have \cle{some} ideas.} / \emph{Would you like \cle{some} water?} \\ \hline
                  \cle{Any} / \cle{Any of} & Comptable pluriel / Non comptable & Négatif / Interrogatif standard / Conditionnel & \emph{I don't have \cle{any} money.} / \emph{Do you have \cle{any} questions?} \\ \hline
                  \cle{Much} / \cle{Much of} & Non comptable seulement & Négatif / Interrogatif (quantité) & \emph{Not \cle{much} time.} / \emph{How \cle{much} does it cost?} \\ \hline
                  \cle{All} / \cle{All of} & Comptable pluriel / Non comptable & Totalité (100\%) & \emph{\cle{All} students passed.} / \emph{\cle{All of} the money is gone.} \\ \hline
                  \cle{Most} / \cle{Most of} & Comptable pluriel / Non comptable & Majorité ($>$ 50\%) & \emph{\cle{Most} people agree.} / \emph{\cle{Most of} the work is done.} \\ \hline
                  \end{tabular}
              \end{center}
              \textbf{Règle "of" :} Obligatoire devant déterminant (the, my, this) + nom : \emph{Some \cle{of the} students}. Interdit devant nom nu : \emph{Some students}.
        \item \textbf{Partie 2 : \cle{few} / \cle{a few} / \cle{little} / \cle{a little}}
              \begin{center}
                  \begin{tabular}{|c|c|c|}\hline
                  \textbf{Forme} & \textbf{Nom} & \textbf{Sens (Nuance cruciale)} \\ \hline
                  \cle{Few} (peu, presque pas) & Comptable pluriel & \textbf{Négatif} : quantité insuffisante. \emph{Few people came} (peu de gens sont venus = déception). \\ \hline
                  \cle{A few} (quelques-uns) & Comptable pluriel & \textbf{Positif} : existence d'une petite quantité. \emph{A few people came} (quelques gens sont venus = c'est déjà ça). \\ \hline
                  \cle{Little} (peu, presque pas) & Non comptable & \textbf{Négatif} : \emph{We have \cle{little} time} (presque pas de temps = urgence). \\ \hline
                  \cle{A little} (un peu) & Non comptable & \textbf{Positif} : \emph{We have \cle{a little} time} (nous avons un peu de temps = on peut faire qqch). \\ \hline
                  \end{tabular}
              \end{center}
        \item \textbf{Partie 3 : \cle{each} / \cle{every} / \cle{many} / \cle{both} / \cle{no} / \cle{none}}
              \begin{itemize}
                  \item \cle{Each} / \cle{Every} : Singulier (distributif). \emph{Each} = individu par individu (2+). \emph{Every} = groupe (3+), $\approx$ All. \emph{Each student has a book. Every day I work.}
                  \item \cle{Many} / \cle{Many of} : Comptable pluriel. Affirmatif (formel), Négatif, Interrogatif. \emph{Many students}. \emph{How many?}
                  \item \cle{Both} / \cle{Both of} : Deux éléments, pluriel. \emph{Both options are good.}
                  \item \cle{No} + nom (nul, zéro) : \emph{No student failed.}
                  \item \cle{None} / \cle{None of} : Pronom (remplace nom). Verbe singulier ou pluriel (souvent pluriel si "of" + pluriel). \emph{None of the students \cle{have} arrived.}
              \end{itemize}
    \end{enumerate}
\end{methode}

\begin{exoresolu}[Choix des quantificateurs : Texte à trous (Niveau Bac)]
    \textbf{Énoncé :} Complétez le texte suivant en choisissant le quantificateur approprié parmi la liste. Attention à la nature des noms (C = comptable, NC = non comptable) et au sens.
    \textit{Liste :} \cle{some} / \cle{any} / \cle{much} / \cle{all} / \cle{most} / \cle{a few} / \cle{few} / \cle{little} / \cle{a little} / \cle{each} / \cle{every} / \cle{many} / \cle{both} / \cle{no} / \cle{none of}.
    \begin{quote}
        Managing a small business requires discipline. (1) \_\_\_\_\_\_\_\_\_\_ entrepreneur must keep records. (2) \_\_\_\_\_\_\_\_\_\_ businesses fail because they have (3) \_\_\_\_\_\_\_\_\_\_ capital to start with. They borrow (4) \_\_\_\_\_\_\_\_\_\_ money from banks but (5) \_\_\_\_\_\_\_\_\_\_ banks require collateral. If you have (6) \_\_\_\_\_\_\_\_\_\_ savings, you can invest. (7) \_\_\_\_\_\_\_\_\_\_ my friends started businesses last year; (8) \_\_\_\_\_\_\_\_\_\_ them succeeded. (9) \_\_\_\_\_\_\_\_\_\_ the profit was reinvested. (10) \_\_\_\_\_\_\_\_\_\_ people think it's easy, but (11) \_\_\_\_\_\_\_\_\_\_ have the patience needed. You need (12) \_\_\_\_\_\_\_\_\_\_ luck too, but hard work matters (13) \_\_\_\_\_\_\_\_\_\_ day.
    \end{quote}
    \tcblower
    \textbf{Solution détaillée et justifications :}
    \begin{enumerate}
        \item \textbf{\cle{Every}} (ou \cle{Each}). \textit{Justification :} "Entrepreneur" singulier comptable, sens distributif (chacun). \emph{Every} pour généralité (3+).
        \item \textbf{\cle{Many}}. \textit{Justification :} "Businesses" pluriel comptable. Phrase affirmative (style formel/écrit) $\rightarrow$ \emph{Many}. (\emph{A lot of} possible à l'oral, mais \emph{Many} attendu ici).
        \item \textbf{\cle{little}}. \textit{Justification :} "Capital" NC. Sens négatif (insuffisant pour démarrer) $\rightarrow$ \emph{little} (peu, presque pas). \emph{A little} serait positif (ils ont un petit capital, ça irait).
        \item \textbf{\cle{some}}. \textit{Justification :} "Money" NC. Phrase affirmative (action réelle : ils empruntent) $\rightarrow$ \emph{some}.
        \item \textbf{\cle{most}} (ou \cle{most of the}). \textit{Justification :} "Banks" pluriel. Majorité. Devant nom nu $\rightarrow$ \emph{Most banks}. (Si "the banks" $\rightarrow$ \emph{Most of the banks}).
        \item \textbf{\cle{some}} (ou \cle{a little}). \textit{Justification :} "Savings" pluriel (souvent traité comme pluriel comptable en anglais : \emph{savings are}). Conditionnel "If you have..." $\rightarrow$ \emph{some} (quantité indéfinie positive). \emph{A little savings} incorrect (savings pluriel).
        \item \textbf{\cle{Both}}. \textit{Justification :} "friends" (impliqué 2 personnes car "Both" suit). "Both my friends" = mes deux amis. Verbe pluriel \emph{started}.
        \item \textbf{\cle{Both of}} (ou \cle{Both}). \textit{Justification :} "Them" pronom $\rightarrow$ obligatoire \emph{of}. \emph{Both of them}. Verbe pluriel \emph{succeeded}.
        \item \textbf{\cle{All}} (ou \cle{All of the}). \textit{Justification :} "Profit" NC. Totalité. \emph{All the profit} ou \emph{All of the profit}. Devant "the" $\rightarrow$ \emph{of} recommandé.
        \item \textbf{\cle{Many}}. \textit{Justification :} "People" pluriel. Phrase affirmative (concession) $\rightarrow$ \emph{Many}.
        \item \textbf{\cle{Few}}. \textit{Justification :} "People" (sous-entendu). Sens négatif (peu de gens ont la patience = critique) $\rightarrow$ \emph{Few}. \emph{A few} signifierait "quelques-uns ont la patience" (positif, contredit "but").
        \item \textbf{\cle{a little}}. \textit{Justification :} "Luck" NC. Sens positif (il en faut un peu, c'est utile) $\rightarrow$ \emph{a little}.
        \item \textbf{\cle{Every}}. \textit{Justification :} "Day" singulier. Fréquence quotidienne $\rightarrow$ \emph{Every day} (locution adverbiale). \emph{Each day} possible mais \emph{Every} plus naturel pour routine.
    \end{enumerate}
    \textbf{Conclusion :} La distinction \textbf{Few/A Few} et \textbf{Little/A Little} (valence négative vs positive) et le choix \textbf{Some/Affirmatif vs Any/Négatif-Interrogatif} sont les pièges n°1 au Bac. Vérifiez toujours : Nom Comptable ? Non Comptable ? Phrase + / - / ? Sens "presque rien" vs "un peu" ?
\end{exoresolu}

\courssub{Developing Interest in a Trade or Profession (Reading \& Vocabulary)}

\begin{methode}[Lire et comprendre un texte sur un métier/une profession (Reading Comprehension)]
    \textbf{Objectif :} Extraire des informations explicites et implicites, identifier le vocabulaire technique du métier, et analyser la structure d'un texte descriptif/argumentatif sur une carrière.
    \begin{enumerate}
        \item \textbf{Étape 1 : Anticipation (Before reading).} Lire le titre, les sous-titres, les images. Prédire le thème (ex: \emph{A Day in the Life of a Welder}, \emph{Why I Chose Nursing}). Activer le vocabulaire connu (\cle{tools}, \cle{skills}, \cle{training}, \cle{apprenticeship}, \cle{shift}, \cle{salary}, \cle{challenges}, \cle{rewards}).
        \item \textbf{Étape 2 : Lecture globale (Skimming).} Lire vite pour l'idée générale (Gist). Qui ? Quoi ? Où ? Ton (positif/négatif/neutre) ? Structure (Introduction $\rightarrow$ Formation $\rightarrow$ Quotidien $\rightarrow$ Avantages/Inconvénients $\rightarrow$ Conseil/Conclusion).
        \item \textbf{Étape 3 : Lecture détaillée (Scanning).} Chercher des infos précises : dates, noms, chiffres, mots-clés de la question. Surligner les \cle{connecteurs logiques} (\emph{however, therefore, moreover, because, although}) pour suivre l'argumentation.
        \item \textbf{Étape 4 : Inférence lexicale (Guessing meaning).} Deviner le sens des mots inconnus par le contexte (définition dans la phrase, opposition, exemple, famille de mots). Ex: \emph{"The job is \cle{grueling}, \emph{it requires long hours of standing and heavy lifting}.} $\rightarrow$ \emph{grueling} = épuisant/dur.
        \item \textbf{Étape 5 : Réponse aux questions.}
              \begin{itemize}
                  \item \textbf{Right/Wrong/Not mentioned :} Se baser \textbf{strictement} sur le texte.
                  \item \textbf{Wh-questions :} Réponse complète, en reprenant les mots du texte si possible (reformulation simple).
                  \item \textbf{Vocabulary :} Synonymes/antonymes trouvés dans le texte.
              \end{itemize}
    \end{enumerate}
\end{methode}

\begin{exoresolu}[Compréhension de texte : "Devenir électricien industriel"]
    \textbf{Énoncé :} Lisez l'extrait et répondez aux questions.
    \begin{quote}
        \textbf{Powering the Future: My Path as an Industrial Electrician}
        \par
        When I finished my Probatoire in Electrical Engineering at the Lycée Technique de Douala, I had a clear goal: become an industrial electrician. Many of my classmates chose office jobs, but \cle{few} of them understood the satisfaction of seeing a complex machine start up because of your wiring.
        \par
        My \cle{apprenticeship} at SONARA lasted three years. It was \cle{grueling} at times — night shifts, high voltage risks, and strict safety protocols. However, I earned \cle{a little} money which helped my family, and I learned \cle{a great deal} from senior technicians. \cle{Most of} the theory I learned at school became real when I had to troubleshoot a faulty PLC (Programmable Logic Controller).
        \par
        Today, I work for a manufacturing plant in Bassa. \cle{Every} day brings a new challenge: preventive maintenance, reading blueprints, or programming automation systems. The salary is good, but \cle{the real reward} is solving problems. \cle{No} two days are alike. I advise young students: \cle{do not} fear technical trades. They offer stability and pride. \cle{Both} the country and your future need skilled hands.
    \end{quote}
    \textbf{Questions :}
    \begin{enumerate}
        \item Are these statements Right (R), Wrong (W), or Not Mentioned (NM)? Justify with a quote.
              \begin{enumerate}
                  \item The writer studied at a technical high school.
                  \item He earned a lot of money during his apprenticeship.
                  \item He works in an office now.
              \end{enumerate}
        \item Answer in your own words (based on the text):
              \begin{enumerate}
                  \item Why did the writer choose this trade over an office job?
                  \item What does "grueling" mean in paragraph 2? Find a clue in the text.
              \end{enumerate}
        \item Vocabulary: Find in the text words/expressions meaning:
              \begin{enumerate}
                  \item A period of training working for a skilled person (para 2).
                  \item Documents showing technical plans (para 3).
                  \item Maintenance done to prevent breakdowns (para 3).
              \end{enumerate}
        \item Grammar in context: Explain the difference in meaning between "\cle{few} of them understood" (para 1) and "I learned \cle{a great deal}" (para 2). Why not "A few of them understood"?
    \end{enumerate}
    \tcblower
    \textbf{Solution détaillée :}
    \begin{enumerate}
        \item \textbf{Right/Wrong/Not Mentioned :}
              \begin{enumerate}
                  \item \textbf{R (Right).} \textit{Justification :} "When I finished my Probatoire in Electrical Engineering at the \cle{Lycée Technique de Douala}..."
                  \item \textbf{W (Wrong).} \textit{Justification :} Text says: "I earned \cle{a little} money". \emph{A little} = une petite quantité (pas "a lot"). Note: "Little money" (sans 'a') serait négatif (peu d'argent). Ici "a little" est positif mais quantité petite.
                  \item \textbf{W (Wrong).} \textit{Justification :} "Today, I work for a \cle{manufacturing plant in Bassa}." Pas un bureau.
              \end{enumerate}
        \item \textbf{Questions ouvertes :}
              \begin{enumerate}
                  \item \textbf{He chose it for the satisfaction of seeing machines work thanks to his skills.} (Citation: "\emph{the satisfaction of seeing a complex machine start up because of your wiring}").
                  \item \textbf{"Grueling" means extremely tiring and difficult.} \textit{Clue :} "\emph{night shifts, high voltage risks, and strict safety protocols}".
              \end{enumerate}
        \item \textbf{Vocabulaire :}
              \begin{enumerate}
                  \item \cle{Apprenticeship} (l'apprentissage / stage professionnel).
                  \item \cle{Blueprints} (les plans / schémas techniques).
                  \item \cle{Preventive maintenance} (la maintenance préventive).
              \end{enumerate}
        \item \textbf{Grammaire en contexte :}
              \begin{itemize}
                  \item "\cle{Few of them understood}" (sans 'a') $\rightarrow$ \textbf{Sens négatif} : "Presque aucun d'entre eux n'a compris" / "Très peu ont compris". L'auteur insiste sur le manque de compréhension de ses camarades.
                  \item "\cle{A few of them understood}" (avec 'a') $\rightarrow$ \textbf{Sens positif} : "Quelques-uns ont compris (et c'est déjà bien)". Cela contredirait l'opposition "Many chose office jobs, BUT few understood...".
                  \item "\cle{A great deal}" $\rightarrow$ Quantificateur pour nom non comptable (\emph{theory, knowledge}), synonyme de \emph{a lot of} / \emph{much}, registre plus formel/écrit. Signifie "beaucoup / énormément".
              \end{itemize}
    \end{enumerate}
    \textbf{Conclusion :} La lecture active (surligner, annoter les quantificateurs, connecteurs) permet de valider Vrai/Faux et de deviner le vocabulaire technique (\emph{apprenticeship, blueprints, PLC}) sans dictionnaire.
\end{exoresolu}

\courssub{Writing about Trades and Professions (Writing \& Integration)}

\begin{methode}[Rédiger une composition sur un métier/une profession (Writing)]
    \textbf{Objectif :} Produire un texte cohérent (150-200 mots) développant l'intérêt pour un métier technique, structuré en paragraphes, utilisant le vocabulaire spécifique, les quantificateurs et les connecteurs logiques.
    \begin{enumerate}
        \item \textbf{Analyser le sujet :} Identifier le métier visé (ou choisir le sien : \emph{welder, nurse, IT technician, accountant, mechanic}). Identifier l'angle : \emph{Why I chose it / A typical day / Challenges and Rewards / Advice to youth}.
        \item \textbf{Plan type (4 paragraphes) :}
              \begin{enumerate}
                  \item \textbf{Introduction (Hook + Thesis) :} Accroche (citation, statistique, anecdote) $\rightarrow$ Présentation du métier $\rightarrow$ Phrase thèse : \emph{"I am interested in [Job] because it offers [Reason 1], [Reason 2], and [Reason 3]."}
                  \item \textbf{Corps 1 : Formation et Compétences (Training \& Skills).} \emph{To become a..., one needs...} (Vocab: \cle{apprenticeship}, \cle{vocational training}, \cle{technical drawing}, \cle{problem-solving}, \cle{teamwork}). Utiliser \emph{must/have to} (obligation).
                  \item \textbf{Corps 2 : Quotidien, Défis et Récompenses (Daily Life, Challenges, Rewards).} \emph{A typical day involves...} (Vocab: \cle{shift work}, \cle{deadlines}, \cle{troubleshooting}, \cle{hands-on}). Utiliser quantificateurs : \emph{Many tasks, Most of the time, A few challenges, Little free time}. Contraste : \emph{Although it is hard, it is rewarding.}
                  \item \textbf{Conclusion :} Résumé + Ouverture/Conseil. \emph{In conclusion, [Job] is vital for the economy. I encourage young people to...} (Utiliser \emph{Let us / Let me urge you to...}).
              \end{enumerate}
        \item \textbf{Check-list linguistique (Self-Correction) :}
              \begin{itemize}
                  \item \textbf{Quantificateurs :} Ai-je utilisé \emph{some/any/much/many/few/a few/little/a little/both/each/every} correctement ?
                  \item \textbf{Connecteurs :} \emph{Firstly, Furthermore, However, Therefore, For instance, In conclusion.}
                  \item \textbf{Temps :} Présent simple (faits généraux), Passé simple (expérience personnelle), Futur (projets).
                  \item \textbf{Registre :} Semi-formel (pas de SMS, pas de contractions excessives en écrit formel \emph{do not} vs \emph{don't}).
              \end{itemize}
    \end{enumerate}
\end{methode}

\begin{exoresolu}[Rédaction guidée : "Why I want to be a..."]
    \textbf{Énoncé :} Écrivez une composition de 150-180 mots sur le sujet suivant : \textbf{"Developing an interest in a technical trade: The case of Industrial Maintenance."} Utilisez au moins 5 quantificateurs différents étudiés (surlignez-les dans votre copie) et 3 connecteurs logiques.
    \tcblower
    \textbf{Proposition de rédaction corrigée (Modèle Bac) :}
    \begin{quote}
        \textbf{Building the Nation: My Interest in Industrial Maintenance}
        \par
        \cle{Every} year, \cle{many} young Cameroonians hesitate between university and technical schools. Personally, I have developed a strong interest in \textbf{Industrial Maintenance} because it combines practical skills, job security, and service to the industry.
        \par
        \textbf{Firstly}, to enter this field, \cle{each} student must master mechanics, hydraulics, and PLC programming. \cle{Most of} the training happens in workshops, not just classrooms. \cle{A few} months of internship in a plant like ALUCAM or SOCAPALM are compulsory. It is \cle{not} an easy path; \cle{little} free time remains during apprenticeship, but \cle{a little} sacrifice builds a solid career.
        \par
        \textbf{However}, the rewards are significant. \cle{Both} preventive and corrective maintenance are vital. \cle{No} production line runs without us. \cle{Few} jobs offer such immediate visibility of results: a stopped machine restarts, the factory saves money. \cle{Some} days are stressful with emergency breakdowns, \cle{but} the team spirit solves everything.
        \par
        \textbf{In conclusion}, Industrial Maintenance is the backbone of our economy. \cle{Let me} advise my peers: choose a trade with \cle{both} hands and brain. \cle{Our} future needs skilled technicians.
    \end{quote}
    \textbf{Analyse des items requis (Vérification) :}
    \begin{itemize}
        \item \textbf{Quantificateurs (9 utilisés) :} \cle{Every} (l.1), \cle{many} (l.1), \cle{each} (l.5), \cle{Most of} (l.6), \cle{A few} (l.6), \cle{little} (l.7), \cle{a little} (l.7), \cle{No} (l.10), \cle{Few} (l.10), \cle{Some} (l.11), \cle{Both} (l.9, l.13). ✓
        \item \textbf{Connecteurs (4 utilisés) :} \cle{Firstly}, \cle{However}, \cle{But} (coordination), \cle{In conclusion}. ✓
        \item \textbf{Structure :} 4 paragraphes clairs (Intro, Formation, Réalités/Récompenses, Conclusion). ✓
        \item \textbf{Vocabulaire technique :} \emph{PLC programming, hydraulics, internship, preventive/corrective maintenance, production line, breakdowns}. ✓
        \item \textbf{Grammaire :} \emph{Let me advise} (politesse/requête), \emph{must master} (obligation), \emph{combines} (présent simple).
    \end{itemize}
    \textbf{Conclusion :} Ce modèle montre l'intégration naturelle des points de grammaire (quantificateurs, requêtes) dans une écriture argumentée sur un métier technique.
\end{exoresolu}

\begin{attention}[Les 5 erreurs qui coûtent des points au Bac (Unité 2)]
    \begin{enumerate}
        \item \textbf{Confusion \cle{Few} vs \cle{A Few} / \cle{Little} vs \cle{A Little} :} Écrire "I have \emph{few} money" (NC $\rightarrow$ \emph{little}) ou "There are \emph{a few} water" (NC). Surtout : ne pas saisir la nuance \textbf{négative} (\emph{few/little} = "presque pas, insuffisant") vs \textbf{positive} (\emph{a few/a little} = "quelques/une petite quantité, existe"). \textit{Astuce :} Si la phrase exprime une plainte ou un manque $\rightarrow$ \emph{few/little}. Si elle exprime une possibilité ou un espoir $\rightarrow$ \emph{a few/a little}.
        \item \textbf{\cle{Some} en phrase négative / \cle{Any} en phrase affirmative (sauf offre/demande) :} \emph{"I don't have some money"} (Faux $\rightarrow$ \emph{any}). \emph{"Do you want any coffee?"} (Faux si offre polie $\rightarrow$ \emph{some}). Règle : \emph{Some} = Affirmatif + Offre/Demande polie. \emph{Any} = Négatif + Interrogatif standard + Conditionnel.
        \item \textbf{Oubli de \cle{of} après quantificateur + Déterminant (the, my, this, these) :} \emph{"Most students"} (OK) MAIS \emph{"Most \textbf{of} the students"} (Obligatoire). \emph{"All the work"} (OK courant) MAIS \emph{"All \textbf{of} the work"} (Préféré formel). \emph{"Both my friends"} (OK) MAIS \emph{"Both \textbf{of} my friends"} (Standard). \emph{None of, Some of, Many of} : \textbf{OF} toujours obligatoire devant déterminant.
        \item \textbf{Accord verbe avec \cle{None of} / \cle{All of} / \cle{Most of} + Pluriel :} \emph{"None of the students \textbf{has} arrived"} (Acceptable formel singulier) MAIS \emph{"None of the students \textbf{have} arrived"} (Standard moderne, pluriel attraction). Au Bac : \textbf{Pluriel} recommandé après \emph{None of / Most of / All of} + Nom Pluriel. \emph{Each / Every / No} + Singulier $\rightarrow$ Verbe \textbf{Singulier} impératif.
        \item \textbf{Mauvaise formation des requêtes polies (Speaking/Writing) :} \emph{"Let me \textbf{to} help you"} (Faux $\rightarrow$ \textbf{Bare infinitive} sans \emph{to}). \emph{"Could you \textbf{to} pass the salt?"} (Faux $\rightarrow$ Bare infinitive). \emph{"Can you help me?"} (Trop familier pour un examen formel $\rightarrow$ Préférer \emph{Could you / Would you mind + V-ing}). Oubli de \emph{please} en fin de phrase.
    \end{enumerate}
\end{attention}

\newpage

\coursec{Exercices}

\courssub{Je vérifie mes connaissances}

\begin{exercice}[QCM : At the market]
Choose the correct answer (A, B or C).
\begin{enumerate}
\item I would like \textbf{........} of bananas, please.
\begin{itemize}
\item[A.] a loaf \item[B.] a bunch \item[C.] a tin
\end{itemize}
\item \textbf{........} you show me that shirt, please?
\begin{itemize}
\item[A.] Could \item[B.] So \item[C.] Much
\end{itemize}
\item \textbf{........} money do you spend every week?
\begin{itemize}
\item[A.] How many \item[B.] How much \item[C.] How few
\end{itemize}
\item ``I like saving money.'' --- ``\textbf{........}''
\begin{itemize}
\item[A.] Neither do I \item[B.] So do I \item[C.] So I do
\end{itemize}
\item There is \textbf{........} rice left in the bag; we must buy more.
\begin{itemize}
\item[A.] a few \item[B.] little \item[C.] many
\end{itemize}
\end{enumerate}
\end{exercice}

\begin{exercice}[True or False]
Say \textbf{True} or \textbf{False} and justify your answer with a correct sentence.
\begin{enumerate}
\item We say ``a loaf of milk''.
\item ``Could you help me, please?'' is a polite request.
\item We use ``how many'' with uncountable nouns like money.
\item ``Neither did he'' expresses agreement with a negative sentence.
\item ``A few'' and ``few'' have exactly the same meaning.
\end{enumerate}
\end{exercice}

\begin{exercice}[Definitions]
Give a short definition in English (one sentence each) of the following words:
\begin{enumerate}
\item a customer
\item a bargain
\item a tontine
\item an apprentice
\item a budget
\end{enumerate}
\end{exercice}

\begin{exercice}[Quick calculations]
Mama Rose sells tomatoes at Mokolo market. Answer with a complete sentence.
\begin{enumerate}
\item One kilogram of tomatoes costs 500 FCFA. How much do 3 kilograms cost?
\item A customer buys a bunch of plantains at 1\,500 FCFA and a tin of tomatoes at 800 FCFA. How much does he pay in total?
\item The customer gives a 5\,000 FCFA note. How much change does Mama Rose give back?
\item Jean saves 2\,000 FCFA every week in his tontine. How much does he save in 6 weeks?
\end{enumerate}
\end{exercice}

\courssub{Je m'entraîne}

\begin{exercice}[Matching : foodstuffs and partitives]
Match each foodstuff (1--10) with the correct partitive expression (a--j).

\begin{center}
\begin{tabularx}{\linewidth}{|l|X|l|X|}
\hline
\rowcolor{popblueL} \textbf{N°} & \textbf{Foodstuff} & \textbf{Letter} & \textbf{Partitive} \\
\hline
1 & bread & a & a bunch of \\
\hline
2 & bananas & b & a tin of \\
\hline
3 & rice (quantity) & c & a loaf of \\
\hline
4 & sardines & d & a cup of \\
\hline
5 & water & e & a bag of \\
\hline
6 & eggs & f & a bottle of \\
\hline
7 & groundnuts & g & a packet of \\
\hline
8 & sugar (for tea) & h & a dozen \\
\hline
9 & palm oil & i & a glass of \\
\hline
10 & plantains & h & a heap of \\
\hline
\end{tabularx}
\end{center}

Write your answers like this : 1 -- c.
\end{exercice}

\begin{exercice}[Dialogue completion : At the shop]
Complete this dialogue between a customer and a shopkeeper in Yaoundé. Use : \emph{could you, how much, a tin of, a loaf of, please, let me, a bunch of}.

\textbf{Shopkeeper :} Good morning, madam. What can I do for you?

\textbf{Customer :} Good morning. (1) \textbf{........} have \textbf{........} (2) bread and \textbf{........} (3) sardines, please.

\textbf{Shopkeeper :} Here you are. Anything else?

\textbf{Customer :} Yes. (4) \textbf{........} are those bananas?

\textbf{Shopkeeper :} 500 FCFA the bunch.

\textbf{Customer :} (5) \textbf{........} take two, (6) \textbf{........}. That is all, thank you.
\end{exercice}

\begin{exercice}[Grammar : determiners in context]
Fill in the blanks with : \emph{some, any, much, few, a little, each, none of, both, every, many}.

Mr Ndi is a trader at the central market in Douala. (1) \textbf{........} morning, he opens his shop at 6 a.m. He sells (2) \textbf{........} foodstuffs : rice, oil and beans. He does not have (3) \textbf{........} free time during the day because he has (4) \textbf{........} customers. (5) \textbf{........} customer receives a receipt. (6) \textbf{........} his sons help him after school. Unfortunately, (7) \textbf{........} of his neighbours wants to learn the trade. He saves (8) \textbf{........} money every month, but there are (9) \textbf{........} banks near his house, so he keeps his savings in a tontine. He does not spend (10) \textbf{........} money on useless things.
\end{exercice}

\begin{exercice}[Sentence transformation : So do I / Neither did he]
Express agreement with each sentence using \emph{So + auxiliary + subject} or \emph{Neither + auxiliary + subject}.

Example : ``Amina saves money every month.'' (Paul) $\rightarrow$ \emph{So does Paul.}

\begin{enumerate}
\item ``I buy my food at the market.'' (my brother)
\item ``Clarisse did not spend her salary on clothes.'' (Eric)
\item ``We are learning a trade at the technical college.'' (they)
\item ``My father does not borrow money from the bank.'' (my uncle)
\item ``The apprentices worked hard last week.'' (the students)
\item ``I have never sold goods on credit.'' (she)
\end{enumerate}
\end{exercice}

\courssub{Je me prépare au Bac}

\begin{exercice}[Reading comprehension : The apprentice dressmaker]
Read the text carefully and answer the questions.

\textbf{Text :} \emph{Mariam is a nineteen-year-old girl from Bafoussam. After her BEPC, she decided to become a dressmaker because she has always loved beautiful clothes. Every morning, she walks to Madame Fotso's workshop in the town centre, where she is an apprentice with four other girls. She pays 5\,000 FCFA every month for her training. At the workshop, she learns to take measurements, cut fabric and use a sewing machine. In the beginning, she made many mistakes, but now she can sew a complete dress in two days. After three years of training, Mariam wants to open her own shop near the market. She has already saved 150\,000 FCFA in a tontine with her friends. ``A trade is better than a diploma without a job,'' she often says. Her parents are very proud of her determination.}

\begin{enumerate}
\item Say \textbf{True} or \textbf{False} and justify with a quotation from the text.
\begin{enumerate}
\item[a.] Mariam is an only apprentice in the workshop.
\item[b.] Her training is free of charge.
\item[c.] She can now sew a dress in two days.
\end{enumerate}
\item Answer these questions in complete sentences.
\begin{enumerate}
\item[a.] Why did Mariam choose dressmaking?
\item[b.] What does she learn at the workshop? (Give two activities.)
\item[c.] How much has she saved for her future shop?
\end{enumerate}
\item Find in the text the synonyms of : (a) \emph{errors} ; (b) \emph{pleased}.
\item What do these words refer to in the text : (a) ``she'' (line 3) ; (b) ``her'' (last line)?
\end{enumerate}
\end{exercice}

\begin{exercice}[Listening comprehension : The tontine meeting]
Your teacher will read a dialogue twice about a tontine in a quarter of Bafoussam. Listen carefully and complete the grid, then answer the questions.

\begin{center}
\begin{tabularx}{\linewidth}{|l|X|X|}
\hline
\rowcolor{popblueL} \textbf{Member} & \textbf{Amount contributed} & \textbf{Frequency} \\
\hline
Mr Kamga & ................ & ................ \\
\hline
Mrs Abena & ................ & ................ \\
\hline
Paul & ................ & ................ \\
\hline
\end{tabularx}
\end{center}

\begin{enumerate}
\item How many members are there in the tontine?
\item Who receives the money this month?
\item What will Mrs Abena do with the money?
\item How much money is collected in total at each meeting?
\item Write two polite requests you heard in the dialogue.
\end{enumerate}
\end{exercice}

\begin{exercice}[Essay : My future profession]
The Ministry of Employment and Vocational Training organises an essay competition in your school on the theme : \emph{``Choosing a trade for a better future.''}

Write a composition of about \textbf{120 words} in which you :
\begin{enumerate}
\item name the trade or profession you would like to embrace ;
\item explain two reasons for your interest in this trade ;
\item describe the training required and where you can receive it in Cameroon ;
\item say how this trade will help you and your community.
\end{enumerate}

Use at least \textbf{three} quantifiers studied in the lesson (\emph{some, many, a few, much, each, both...}) and \textbf{one} polite structure (\emph{let me...} / \emph{could you...}).
\end{exercice}

\newpage

\coursec{Corrigés des exercices}

\begin{corrige}[Exercice 1 — QCM : At the market]
\begin{enumerate}
\item \textbf{B — a bunch.} We say \emph{a bunch of bananas}. A \emph{loaf} is used with bread, and a \emph{tin} with sardines or tomatoes.
\item \textbf{A — Could.} \emph{Could you show me that shirt, please?} is a polite request. \emph{So} and \emph{Much} cannot introduce a request.
\item \textbf{B — How much.} \emph{Money} is an uncountable noun, so we ask \emph{How much money...?} \emph{How many} is used with countable nouns.
\item \textbf{B — So do I.} To agree with a positive sentence in the present simple, we say \emph{So + do/does + subject}. \emph{Neither do I} would agree with a negative sentence.
\item \textbf{B — little.} \emph{Rice} is uncountable. \emph{Little rice left} means ``almost no rice'', so they must buy more. \emph{A few} and \emph{many} are used with countable nouns.
\end{enumerate}
\textbf{Points de vigilance :} distinguish \emph{how much / how many} (uncountable vs countable) and \emph{So do I / Neither do I} (positive vs negative agreement).
\end{corrige}

\begin{corrige}[Exercice 2 — True or False]
\begin{enumerate}
\item \textbf{False.} We say \emph{a loaf of bread} and \emph{a bottle} (or \emph{a packet}) \emph{of milk}. A loaf is a partitive for bread only.
\item \textbf{True.} \emph{Could you help me, please?} is a polite request : \emph{could you} + bare infinitive + \emph{please} is the polite form studied in the lesson.
\item \textbf{False.} We use \emph{how much} with uncountable nouns like money : \emph{How much money do you have?} \emph{How many} is for countable nouns.
\item \textbf{True.} \emph{Neither did he} expresses agreement with a negative sentence in the past : ``I did not go.'' --- ``Neither did he.''
\item \textbf{False.} \emph{A few} means ``some, a small number'' (positive idea), while \emph{few} means ``almost none'' (negative idea) : \emph{I have a few friends} (I am happy) / \emph{I have few friends} (I am lonely).
\end{enumerate}
\textbf{Points de vigilance :} the justification must be a complete, grammatically correct sentence, not just ``True'' or ``False''.
\end{corrige}

\begin{corrige}[Exercice 3 — Definitions]
Suggested answers (any correct one-sentence definition is accepted) :
\begin{enumerate}
\item \textbf{A customer} is a person who buys goods or services from a shop or a trader.
\item \textbf{A bargain} is something bought at a price lower than its real value (a good deal).
\item \textbf{A tontine} is a group of people who contribute money regularly and give the total amount to one member in turn.
\item \textbf{An apprentice} is a person who learns a trade by working with a skilled worker or craftsman.
\item \textbf{A budget} is a plan that shows how much money you earn and how you will spend it.
\end{enumerate}
\textbf{Points de vigilance :} a definition must contain a general word (\emph{a person who...}, \emph{a plan that...}, \emph{a group of people who...}) and one complete sentence only.
\end{corrige}

\begin{corrige}[Exercice 4 — Quick calculations]
\begin{enumerate}
\item $3 \times 500 = 1\,500$. \textbf{Three kilograms of tomatoes cost 1\,500 FCFA.}
\item $1\,500 + 800 = 2\,300$. \textbf{He pays 2\,300 FCFA in total.}
\item $5\,000 - 2\,300 = 2\,700$. \textbf{Mama Rose gives back 2\,700 FCFA.}
\item $2\,000 \times 6 = 12\,000$. \textbf{Jean saves 12\,000 FCFA in six weeks.}
\end{enumerate}
\textbf{Points de vigilance :} each answer must be a complete sentence with the unit (FCFA) ; show the operation before the result.
\end{corrige}

\begin{corrige}[Exercice 5 — Matching : foodstuffs and partitives]
Correct answers :
\begin{center}
\begin{tabularx}{\linewidth}{|l|X|l|}
\hline
\rowcolor{popblueL} \textbf{N°} & \textbf{Foodstuff} & \textbf{Answer} \\
\hline
1 & bread & 1 -- c : a loaf of bread \\
\hline
2 & bananas & 2 -- a : a bunch of bananas \\
\hline
3 & rice (quantity) & 3 -- e : a bag of rice \\
\hline
4 & sardines & 4 -- b : a tin of sardines \\
\hline
5 & water & 5 -- i : a glass of water \\
\hline
6 & eggs & 6 -- h : a dozen eggs \\
\hline
7 & groundnuts & 7 -- j : a heap of groundnuts \\
\hline
8 & sugar (for tea) & 8 -- d : a cup of sugar \\
\hline
9 & palm oil & 9 -- f : a bottle of palm oil \\
\hline
10 & plantains & 10 -- g : a packet of plantains \\
\hline
\end{tabularx}
\end{center}
\textbf{Points de vigilance :} \emph{a dozen} is used for eggs (twelve) ; \emph{a heap of} suits groundnuts sold in piles at the market ; \emph{a glass of water} but \emph{a bottle of oil}.
\end{corrige}

\begin{corrige}[Exercice 6 — Dialogue completion : At the shop]
Completed dialogue :

\textbf{Shopkeeper :} Good morning, madam. What can I do for you?

\textbf{Customer :} Good morning. (1) \textbf{Let me} have \textbf{a loaf of} (2) bread and \textbf{a tin of} (3) sardines, please.

\textbf{Shopkeeper :} Here you are. Anything else?

\textbf{Customer :} Yes. (4) \textbf{How much} are those bananas?

\textbf{Shopkeeper :} 500 FCFA the bunch.

\textbf{Customer :} (5) \textbf{Could you} take two, (6) \textbf{please}. That is all, thank you.

\textbf{Points de vigilance :} \emph{let me} is followed by the bare infinitive (\emph{let me have}) ; \emph{how much} asks for the price ; \emph{please} at the end keeps the request polite. Note : \emph{a bunch of} was a distractor — the bananas are already sold ``the bunch'' by the shopkeeper.
\end{corrige}

\begin{corrige}[Exercice 7 — Grammar : determiners in context]
Completed text :

Mr Ndi is a trader at the central market in Douala. (1) \textbf{Every} morning, he opens his shop at 6 a.m. He sells (2) \textbf{some} foodstuffs : rice, oil and beans. He does not have (3) \textbf{any} free time during the day because he has (4) \textbf{many} customers. (5) \textbf{Each} customer receives a receipt. (6) \textbf{Both} his sons help him after school. Unfortunately, (7) \textbf{none} of his neighbours \textbf{want} to learn the trade. He saves (8) \textbf{a little} money every month, but there are (9) \textbf{few} banks near his house, so he keeps his savings in a tontine. He does not spend (10) \textbf{much} money on useless things.

Justifications :
\begin{itemize}
\item (1) \emph{Every} + singular noun for a repeated habit ; (2) \emph{some} in a positive sentence ; (3) \emph{any} after a negative verb ; (4) \emph{many} + countable plural (\emph{customers}) ; (5) \emph{each} + singular noun (\emph{customer}) ; (6) \emph{both} = the two sons ; (7) \emph{none of} + plural noun = zero (verb agrees with plural \emph{neighbours} $\rightarrow$ \emph{want}) ; (8) \emph{a little} + uncountable (\emph{money}) = a small amount, positive idea ; (9) \emph{few} + countable = almost no banks (negative idea, that is why he uses a tontine) ; (10) \emph{much} + uncountable in a negative sentence.
\end{itemize}
\textbf{Points de vigilance :} check countable/uncountable before choosing (\emph{much/many}, \emph{little/few}) and the positive/negative meaning of \emph{a little} vs \emph{little}, \emph{a few} vs \emph{few} ; watch verb agreement after \emph{none of} + plural noun.
\end{corrige}

\begin{corrige}[Exercice 8 — Sentence transformation : So do I / Neither did he]
\begin{enumerate}
\item ``I buy my food at the market.'' (my brother) $\rightarrow$ \textbf{So does my brother.} (present simple, third person : auxiliary \emph{does})
\item ``Clarisse did not spend her salary on clothes.'' (Eric) $\rightarrow$ \textbf{Neither did Eric.} (negative sentence, past : \emph{neither} + \emph{did})
\item ``We are learning a trade at the technical college.'' (they) $\rightarrow$ \textbf{So are they.} (auxiliary \emph{be} in the present continuous)
\item ``My father does not borrow money from the bank.'' (my uncle) $\rightarrow$ \textbf{Neither does my uncle.} (negative, present simple, third person)
\item ``The apprentices worked hard last week.'' (the students) $\rightarrow$ \textbf{So did the students.} (positive sentence, past simple : auxiliary \emph{did})
\item ``I have never sold goods on credit.'' (she) $\rightarrow$ \textbf{Neither has she.} (\emph{never} makes the sentence negative ; auxiliary \emph{has})
\end{enumerate}
\textbf{Points de vigilance :} the auxiliary in the short answer must copy the tense and auxiliary of the first sentence (\emph{do/does/did/be/have}) ; use \emph{Neither} whenever the first sentence is negative or contains \emph{never}, \emph{not}, \emph{no}.
\end{corrige}

\begin{corrige}[Exercice 9 — Reading comprehension : The apprentice dressmaker]
\begin{enumerate}
\item True or False :
\begin{enumerate}
\item[a.] \textbf{False.} ``She is an apprentice with four other girls'' : there are five apprentices, so she is not alone.
\item[b.] \textbf{False.} ``She pays 5\,000 FCFA every month for her training'' : the training is not free.
\item[c.] \textbf{True.} ``Now she can sew a complete dress in two days.''
\end{enumerate}
\item Questions :
\begin{enumerate}
\item[a.] She chose dressmaking because she has always loved beautiful clothes.
\item[b.] At the workshop, she learns to take measurements and to cut fabric (she also learns to use a sewing machine).
\item[c.] She has already saved 150\,000 FCFA in a tontine for her future shop.
\end{enumerate}
\item Synonyms : (a) \emph{errors} = \textbf{mistakes} ; (b) \emph{pleased} = \textbf{proud}.
\item Reference words : (a) ``she'' (line 3) refers to \textbf{Mariam} ; (b) ``her'' (last line) refers to \textbf{Mariam} (her determination).
\end{enumerate}
\textbf{Barème indicatif (sur 10) :} Q1 : 3 pts (1 pt par réponse justifiée par une citation) ; Q2 : 3 pts ; Q3 : 2 pts ; Q4 : 2 pts.

\textbf{Points de vigilance :} at the Probatoire/Baccalauréat technique, a True/False answer \emph{without a quotation} earns no point ; answers must be complete sentences ; for reference words, name the precise person, not just ``a girl''.
\end{corrige}

\begin{corrige}[Exercice 10 — Listening comprehension : The tontine meeting]
\textbf{Script lu par l'enseignant (à lire deux fois) :}

\emph{Mr Kamga : Good evening, everybody. Welcome to our tontine meeting. Let me open the session, please. As usual, I contribute 10\,000 FCFA every month.}

\emph{Mrs Abena : Good evening. Could you write my name for 5\,000 FCFA, please? I contribute every month too. This month, it is my turn to receive the money. I will buy a sewing machine for my tailoring workshop.}

\emph{Paul : And here is my contribution : 2\,000 FCFA every week, so 8\,000 FCFA for this month.}

\emph{Mr Kamga : Thank you. Could you all sign the register, please? See you next month.}

\textbf{Corrected grid :}
\begin{center}
\begin{tabularx}{\linewidth}{|l|X|X|}
\hline
\rowcolor{popblueL} \textbf{Member} & \textbf{Amount contributed} & \textbf{Frequency} \\
\hline
Mr Kamga & 10\,000 FCFA & every month \\
\hline
Mrs Abena & 5\,000 FCFA & every month \\
\hline
Paul & 2\,000 FCFA (8\,000 FCFA this month) & every week \\
\hline
\end{tabularx}
\end{center}

\textbf{Answers :}
\begin{enumerate}
\item There are \textbf{three} members in the tontine (Mr Kamga, Mrs Abena and Paul).
\item \textbf{Mrs Abena} receives the money this month (``it is my turn to receive the money'').
\item She will \textbf{buy a sewing machine} for her tailoring workshop.
\item Total collected this month : $10\,000 + 5\,000 + 8\,000 = 23\,000$. \textbf{23\,000 FCFA} are collected at this meeting.
\item Two polite requests heard : \emph{``Could you write my name for 5\,000 FCFA, please?''} and \emph{``Could you all sign the register, please?''} (also accepted : \emph{``Let me open the session, please.''})
\end{enumerate}
\textbf{Barème indicatif (sur 10) :} grille : 3 pts (1 pt par ligne complète) ; Q1--Q3 : 3 pts ; Q4 : 2 pts (calcul exact exigé) ; Q5 : 2 pts.

\textbf{Points de vigilance :} for question 4, Paul contributes weekly, so his monthly amount is $2\,000 \times 4 = 8\,000$ FCFA — do not add 2\,000 only ; polite requests must be quoted exactly with \emph{could you... please} or \emph{let me...}.
\end{corrige}

\begin{corrige}[Exercice 11 — Essay : My future profession]
\textbf{Model essay (about 120 words) :}

\emph{Choosing a trade for a better future}

Let me tell you about the trade I would like to embrace : \textbf{carpentry}. I have \textbf{many} reasons for this choice, but I will give two. First, I love working with wood and making beautiful furniture with my hands. Second, \textbf{each} family in my quarter needs tables, chairs and doors, so a carpenter always has customers.

To learn this trade, I can train for three years at a technical college or in a workshop in Douala, where \textbf{a few} experienced carpenters accept apprentices. The training is not expensive, and I will not need \textbf{much} money to start.

With this trade, I will earn a living, create jobs for \textbf{some} young people and build good furniture for my community. Could you imagine a better future?

\textbf{Barème indicatif (sur 20) :}
\begin{itemize}
\item Respect of the task (the 4 points treated) : 6 pts ;
\item Coherence and organisation (introduction, development, conclusion) : 4 pts ;
\item Vocabulary of trades and buying/selling : 4 pts ;
\item Grammar accuracy (quantifiers, polite structure, tenses) : 4 pts ;
\item Length (about 120 words) and presentation : 2 pts.
\end{itemize}
\textbf{Points de vigilance :} the three quantifiers and the polite structure must be used \emph{correctly} (a wrong quantifier earns no grammar point) ; do not copy the model word for word ; stay within the theme and the word limit ; write in paragraphs, not in a list.
\end{corrige}

\newpage

\coursec{Fiche bilan}

\begin{popfigure}
\begin{center}
\begin{center}\tcbox[colback=popink,colframe=popink,coltext=white,arc=9pt,boxsep=4pt,left=6pt,right=6pt]{\bfseries\small Economic Life \& Occupations}\end{center}
\vspace{-2pt}
\begin{tcbraster}[raster columns=3,raster equal height=rows,raster column skip=6pt,raster row skip=6pt,enhanced,blankest]
\begin{tcolorbox}[enhanced,colback=white,colframe=popblue,boxrule=0.9pt,arc=3pt,title={Buying \& Selling},fonttitle=\bfseries\scriptsize,coltitle=white,colbacktitle=popblue,left=4pt,right=4pt,top=2pt,bottom=2pt,boxsep=1pt,toptitle=1.5pt,bottomtitle=1.5pt,halign=left]
\begin{itemize}[nosep,leftmargin=*,itemsep=1pt,font=\scriptsize]
\item market words
\item bargaining
\end{itemize}
\end{tcolorbox}
\begin{tcolorbox}[enhanced,colback=white,colframe=poporange,boxrule=0.9pt,arc=3pt,title={Polite Requests},fonttitle=\bfseries\scriptsize,coltitle=white,colbacktitle=poporange,left=4pt,right=4pt,top=2pt,bottom=2pt,boxsep=1pt,toptitle=1.5pt,bottomtitle=1.5pt,halign=left]
\begin{itemize}[nosep,leftmargin=*,itemsep=1pt,font=\scriptsize]
\item Could you...?
\item Let me + V
\end{itemize}
\end{tcolorbox}
\begin{tcolorbox}[enhanced,colback=white,colframe=popgreen,boxrule=0.9pt,arc=3pt,title={Partitives},fonttitle=\bfseries\scriptsize,coltitle=white,colbacktitle=popgreen,left=4pt,right=4pt,top=2pt,bottom=2pt,boxsep=1pt,toptitle=1.5pt,bottomtitle=1.5pt,halign=left]
\begin{itemize}[nosep,leftmargin=*,itemsep=1pt,font=\scriptsize]
\item a bunch of
\item a cup of
\end{itemize}
\end{tcolorbox}
\begin{tcolorbox}[enhanced,colback=white,colframe=poppurple,boxrule=0.9pt,arc=3pt,title={Personal Finances},fonttitle=\bfseries\scriptsize,coltitle=white,colbacktitle=poppurple,left=4pt,right=4pt,top=2pt,bottom=2pt,boxsep=1pt,toptitle=1.5pt,bottomtitle=1.5pt,halign=left]
\begin{itemize}[nosep,leftmargin=*,itemsep=1pt,font=\scriptsize]
\item budget, save
\item So do I / Neither did he
\end{itemize}
\end{tcolorbox}
\begin{tcolorbox}[enhanced,colback=white,colframe=popgold,boxrule=0.9pt,arc=3pt,title={Quantifiers},fonttitle=\bfseries\scriptsize,coltitle=white,colbacktitle=popgold,left=4pt,right=4pt,top=2pt,bottom=2pt,boxsep=1pt,toptitle=1.5pt,bottomtitle=1.5pt,halign=left]
\begin{itemize}[nosep,leftmargin=*,itemsep=1pt,font=\scriptsize]
\item some, any, much
\item few / little
\item each, every, both, none
\end{itemize}
\end{tcolorbox}
\begin{tcolorbox}[enhanced,colback=white,colframe=poppink,boxrule=0.9pt,arc=3pt,title={Trades \& Writing},fonttitle=\bfseries\scriptsize,coltitle=white,colbacktitle=poppink,left=4pt,right=4pt,top=2pt,bottom=2pt,boxsep=1pt,toptitle=1.5pt,bottomtitle=1.5pt,halign=left]
\begin{itemize}[nosep,leftmargin=*,itemsep=1pt,font=\scriptsize]
\item job vocabulary
\item composition
\end{itemize}
\end{tcolorbox}
\end{tcbraster}

\end{center}
\legende{Carte mentale de la leçon : les six blocs de savoirs à maîtriser (mind map of the lesson).}
\end{popfigure}

\begin{bilan}
\textbf{1. Vocabulary — Buying and selling (au marché / à la boutique)}
\begin{itemize}
\item \cle{stall} (étal), \cle{customer} (client), \cle{seller/vendor} (vendeur), \cle{price} (prix), \cle{bargain} (marchander), \cle{discount} (remise), \cle{change} (la monnaie), \cle{receipt} (reçu).
\item Useful phrases : \emph{How much is it?} — \emph{It costs 500 FCFA.} — \emph{Can you reduce the price?} — \emph{Here is your change.}
\end{itemize}

\textbf{2. Grammar — Polite requests and suggestions}
\begin{center}
\begin{tabularx}{\linewidth}{|l|X|X|}\hline
\rowcolor{popblueL} \textbf{Structure} & \textbf{Example} & \textbf{Use} \\ \hline
Could you / Can you + V...? & Could you give me a bag, please? & polite request \\ \hline
Let me + bare infinitive & Let me help you. & suggestion / offer \\ \hline
Please + imperative & Please weigh the tomatoes. & polite order \\ \hline
\end{tabularx}
\end{center}

\textbf{3. Partitive articles (pour les quantités)}
a \cle{cup} of tea, a \cle{bunch} of plantains, a \cle{loaf} of bread, a \cle{bottle} of oil, a \cle{kilogram} of rice, a \cle{piece} of soap, a \cle{tin} of milk.

\textbf{4. Managing personal finances (gérer son argent)}
\begin{itemize}
\item \cle{budget} (budget), \cle{save money} (économiser), \cle{spend} (dépenser), \cle{earn} (gagner), \cle{debt} (dette), \cle{expenses} (dépenses), \cle{income} (revenu).
\item Agreement : \emph{I save money.} $\rightarrow$ \cle{So do I}. / \emph{He didn't spend much.} $\rightarrow$ \cle{Neither did he}.
\item Adjectives as nouns : \emph{the rich} (les riches), \emph{the poor} (les pauvres), \emph{the unemployed} (les chômeurs).
\end{itemize}

\textbf{5. Determiners and quantifiers (à connaître par c\oe ur)}
\begin{center}
\begin{tabularx}{\linewidth}{|l|X|X|}\hline
\rowcolor{popblueL} \textbf{Quantifier} & \textbf{Used with} & \textbf{Example} \\ \hline
some / any & plural or uncountable & some rice / any money? \\ \hline
much / many & uncountable / plural & much water / many customers \\ \hline
few / a few & plural countable & few traders / a few traders \\ \hline
little / a little & uncountable & little time / a little money \\ \hline
each / every & singular & each student / every day \\ \hline
both (of) / none (of) & two / zero & both shops / none of them \\ \hline
all (of) / most (of) & plural or uncountable & all the food / most sellers \\ \hline
\end{tabularx}
\end{center}

\textbf{6. Trades and professions — Reading \& Writing}
\begin{itemize}
\item \cle{trade} (métier), \cle{carpenter} (menuisier), \cle{tailor} (tailleur), \cle{mechanic} (mécanicien), \cle{hairdresser} (coiffeur), \cle{apprentice} (apprenti), \cle{workshop} (atelier), \cle{skill} (compétence).
\item Writing method : introduction (the trade chosen) $\rightarrow$ body (reasons, skills, training) $\rightarrow$ conclusion (future plans).
\end{itemize}
\end{bilan}

\begin{aretenir}[Checklist avant le Bac]
\begin{itemize}
\item[$\square$] Je connais le vocabulaire du marché : price, bargain, discount, change, stall.
\item[$\square$] Je sais formuler une demande polie avec \emph{Could you...?} et \emph{Can you...?}.
\item[$\square$] J'utilise \emph{Let me + infinitif} pour proposer de l'aide.
\item[$\square$] Je maîtrise les partitifs : a bunch of, a cup of, a loaf of, a bottle of.
\item[$\square$] Je sais répondre avec \emph{So do I} et \emph{Neither did he}.
\item[$\square$] Je connais le vocabulaire des finances : budget, save, spend, earn, income, expenses.
\item[$\square$] Je distingue \emph{much} (indénombrable) et \emph{many} (dénombrable).
\item[$\square$] Je distingue \emph{few/a few} et \emph{little/a little}.
\item[$\square$] J'emploie correctement \emph{each, every, both, none of}.
\item[$\square$] Je connais le vocabulaire des métiers : trade, apprentice, workshop, skill.
\item[$\square$] Je sais rédiger une composition en trois parties sur un métier qui m'intéresse.
\item[$\square$] Je relis ma copie : orthographe, accords, ponctuation.
\end{itemize}
\end{aretenir}

\findecours

\end{document}
